%~Mouliné par MaN_auto v.0.39.2 (60cbe62a) 2026-02-26 13:30:51
\documentclass[AHL,Unicode,longabstracts,MathTools]{cedram}


\usepackage{tikz}
\usepackage{tikz-cd}
\usepackage{mathrsfs}
\usepackage[scr=boondox]{mathalfa}	
\usetikzlibrary{intersections}
\usetikzlibrary{calc}


\DeclareMathOperator{\bus}{b}
\DeclareMathOperator{\spa}{span}
\DeclareMathOperator{\class}{C}
\DeclareMathOperator{\diam}{diam}
\DeclareMathOperator{\id}{id}
\DeclareMathOperator{\inte}{int}
\DeclareMathOperator{\Hff}{dim_{Hf{}f}}
\DeclareMathOperator{\ann}{Ann}
\DeclareMathOperator{\codim}{codim}
\DeclareMathOperator{\Diff}{Diff}

\newcommand{\ii}{\mathrm{i}}
\newcommand{\rmB}{\mathrm{B}}

\newcommand{\grass}{\mathrm{Gr}}
\newcommand{\Fundeb}{\mathrm{Leb}}
\newcommand{\Ann}{\mathrm{Ann}}
\newcommand{\cartan}{a}


\newcommand{\Ndiff}{{\mathsf{NDiff}}}
\DeclareMathOperator{\SL}{{\mathsf{SL}}}
\DeclareMathOperator{\PSL}{{\mathsf{PSL}}}
\DeclareMathOperator{\GL}{{\mathsf{GL}}}
\DeclareMathOperator{\SO}{{\mathsf{SO}}}
\DeclareMathOperator{\PGL}{{\mathsf{PGL}}}
\DeclareMathOperator{\PO}{{\mathsf{PO}}}
\DeclareMathOperator{\Sp}{\mathsf{Sp}}
\DeclareMathOperator{\SU}{\mathsf{SU}}


\newcommand{\scr}{\mathscr}
\newcommand{\hol}{\scr{b}}
%\newcommand{\hJ}[1]{{\scr{h}_{#1}}}
%\newcommand{\hC}[1]{{\scr{h}^{#1}}}
\newcommand{\hJ}[1]{{\scr h}_{#1}}
\newcommand{\hC}[1]{{\scr h}^{#1}}
\renewcommand{\varXi}{\mathscr{G}}
\newcommand{\cc}{\scr{v}}


\renewcommand{\sf}[1]{{\mathsf{#1}}}

\newcommand{\sroot}{{\sf{a}}}
\newcommand{\bb}{{\sf{b}}}
\newcommand{\sfR}{\mathsf{R}}
\newcommand{\sfG}{\sf{G}}

\renewcommand{\flat}{\sf{b}}
\newcommand{\poids}{\sf{\Pi}}
\newcommand{\G}{\sf{\Gamma}}
\newcommand{\roots}{\mathsf{\Phi}}


\newcommand{\cal}{\mathcal}
\DeclareMathOperator{\Ext}{\cal{H}}
\newcommand{\Gr}{\cal{G}}
\newcommand{\cone}{\cal{C}}
\newcommand{\posgen}{\cal{F}^{(2)}}
\newcommand{\Weyl}{\cal{W}}

\newcommand{\calF}{\mathcal{F}}
\newcommand{\calL}{\mathcal{L}}
\newcommand{\calQ}{\mathcal{Q}}

\newcommand{\Z}{\mathbb{Z}}
\newcommand{\R}{\mathbb{R}}
\newcommand{\RR}{\R}
\newcommand{\C}{\mathbb{C}}
\newcommand{\N}{\mathbb{N}}
\renewcommand{\P}{\mathbb{P}}
\newcommand{\T}{\mathbb{T}}
\newcommand{\I}{\mathbb{I}}
\renewcommand{\H}{\mathbb{H}}
\newcommand{\K}{\mathbb{K}}
\newcommand{\HH}{\mathbb{H}}
\newcommand{\bbS}{\mathbb{S}}

\renewcommand{\circle}{\mathbb{S}^1}

\newcommand{\II}{\mathbf{I}}

\newcommand{\eps}{\varepsilon}

\newcommand{\g}{\gamma}
\renewcommand{\t}{\theta}
\newcommand{\bord}{\partial}
\newcommand{\grupo}{\varLambda}
\newcommand{\wk}{\check}
\newcommand{\piS}{\pi_1S}
\newcommand{\ps}{\mu}
\newcommand{\simple}{{\sf \Delta}}
\newcommand{\Fund}{\Phi}
\newcommand{\peso}{\varpi}
\newcommand{\rfr}{\beta}
\newcommand{\BM}{\Omega}
\newcommand{\df}{\omega}
\newcommand{\vi}{{\varphi^\infty_\hol}}
\newcommand{\slroot}{{\tau}}

%%%%
\newcommand{\ov}{\overline}
%\newcommand{\ov}{\bar}\cc

\newcommand{\ovcc}{\bar{\cc}}
\newcommand{\ovrho}{\overline{\!\rho}}
\newcommand{\ovsloot}
{\bar{\slroot}}

\newcommand{\ovxi}{\overline{\!\xi}}
\newcommand{\ovd}{\bar{d}}

\renewcommand{\epsilon}{\varepsilon}
\renewcommand{\/}{\backslash}
\renewcommand{\angle}{\measuredangle}

\newcommand{\<}{\langle}
\renewcommand{\>}{\rangle}


\newcommand{\p}{\mathfrak{p}}

\renewcommand{\ge}{\mathfrak{g}}
\renewcommand{\a}{\mathfrak{a}}
\renewcommand{\k}{\mathfrak{k}}
\newcommand{\fkX}{\mathfrak{X}}

\newcommand{\fksl}{\mathfrak{sl}}
\renewcommand{\sl}{\fksl}
\newcommand{\fkgl}{\mathfrak{gl}}




\newcommand{\cdc}[1]{\bigl(\calL_{#1}\bigr)^*}



%\numberwithin{equation}{section}



%\lb

%%%%%%%%%%%%%%%%%%%%%%%%%%%%%%%%%%%%%%%%%
%Better widebar

\makeatletter
\let\save@mathaccent\mathaccent
\newcommand*\if@single[3]{%
  \setbox0\hbox{${\mathaccent"0362{#1}}^H$}%
  \setbox2\hbox{${\mathaccent"0362{\kern0pt#1}}^H$}%
  \ifdim\ht0=\ht2 #3\else #2\fi
  }
%The bar will be moved to the right by a half of \macc@kerna, which is computed by amsmath:
\newcommand*\rel@kern[1]{\kern#1\dimexpr\macc@kerna}
%If there's a superscript following the bar, then no negative kern may follow the bar;
%an additional {} makes sure that the superscript is high enough in this case:
\newcommand*\widebar[1]{\@ifnextchar^{{\wide@bar{#1}{0}}}{\wide@bar{#1}{1}}}
%Use a separate algorithm for single symbols:
\newcommand*\wide@bar[2]{\if@single{#1}{\wide@bar@{#1}{#2}{1}}{\wide@bar@{#1}{#2}{2}}}
\newcommand*\wide@bar@[3]{%
  \begingroup
  \def\mathaccent##1##2{%
%Enable nesting of accents:
    \let\mathaccent\save@mathaccent
%If there's more than a single symbol, use the first character instead (see below):
    \if#32 \let\macc@nucleus\first@char \fi
%Determine the italic correction:
    \setbox\z@\hbox{$\macc@style{\macc@nucleus}_{}$}%
    \setbox\tw@\hbox{$\macc@style{\macc@nucleus}{}_{}$}%
    \dimen@\wd\tw@
    \advance\dimen@-\wd\z@
%Now \dimen@ is the italic correction of the symbol.
    \divide\dimen@ 3
    \@tempdima\wd\tw@
    \advance\@tempdima-\scriptspace
%Now \@tempdima is the width of the symbol.
    \divide\@tempdima 10
    \advance\dimen@-\@tempdima
%Now \dimen@ = (italic correction / 3) - (Breite / 10)
    \ifdim\dimen@>\z@ \dimen@0pt\fi
%The bar will be shortened in the case \dimen@<0 !
    \rel@kern{0.6}\kern-\dimen@
    \if#31
      \overline{\rel@kern{-0.6}\kern\dimen@\macc@nucleus\rel@kern{0.4}\kern\dimen@}%
      \advance\dimen@0.4\dimexpr\macc@kerna
%Place the combined final kern (-\dimen@) if it is >0 or if a superscript follows:
      \let\final@kern#2%
      \ifdim\dimen@<\z@ \let\final@kern1\fi
      \if\final@kern1 \kern-\dimen@\fi
    \else
      \overline{\rel@kern{-0.6}\kern\dimen@#1}%
    \fi
  }%
  \macc@depth\@ne
  \let\math@bgroup\@empty \let\math@egroup\macc@set@skewchar
  \mathsurround\z@ \frozen@everymath{\mathgroup\macc@group\relax}%
  \macc@set@skewchar\relax
  \let\mathaccentV\macc@nested@a
%The following initialises \macc@kerna and calls \mathaccent:
  \if#31
    \macc@nested@a\relax111{#1}%
  \else
%If the argument consists of more than one symbol, and if the first token is
%a letter, use that letter for the computations:
    \def\gobble@till@marker##1\endmarker{}%
    \futurelet\first@char\gobble@till@marker#1\endmarker
    \ifcat\noexpand\first@char A\else
      \def\first@char{}%
    \fi
    \macc@nested@a\relax111{\first@char}%
  \fi
  \endgroup
}
\makeatother

\let\oldbar\bar
\renewcommand*{\bar}[1]{{\mathchoice{\widebar{#1}}{\widebar{#1}}{\widebar{#1}}{\oldbar{#1}}}}

%%%%%%%%%%%%%%%%%%%%%
\renewcommand*{\to}{\mathchoice{\longrightarrow}{\rightarrow}{\rightarrow}{\rightarrow}}

\let\oldleftrightarrow\leftrightarrow
\renewcommand*{\leftrightarrow}{\mathchoice{\longleftrightarrow}{\oldleftrightarrow}{\oldleftrightarrow}{\oldleftrightarrow}}

\let\oldmapsto\mapsto
\renewcommand*{\mapsto}{\mathchoice{\longmapsto}{\oldmapsto}{\oldmapsto}{\oldmapsto}}

\renewcommand*{\implies}{\mathchoice{\Longrightarrow}{\Rightarrow}{\Rightarrow}{\Rightarrow}}


%~ This fixes spacing issues around \left( and \right) and other similar symbols
\let\originalleft\left 
\let\originalright\right
\renewcommand{\left}{\mathopen{}\mathclose\bgroup\originalleft}
\renewcommand{\right}{\aftergroup\egroup\originalright}

%%%%%%

%%%%%%%%%%%%%%%%%%%%%%%%%%%%%%%%%%%%%%%%%%%%%%%%%%%%%%%%%%%%%%%%%%%%%%%%%%%
\graphicspath{{./figures/}}

\newcommand*{\mk}{\mkern -1mu}
\newcommand*{\Mk}{\mkern -2mu}
\newcommand*{\mK}{\mkern 1mu}
\newcommand*{\MK}{\mkern 2mu}

\hypersetup{urlcolor=purple, linkcolor=blue, citecolor=red}

\newcommand*{\relabel}{\renewcommand{\labelenumi}{(\theenumi)}}
\newcommand*{\romanenumi}{\renewcommand*{\theenumi}{\roman{enumi}}\relabel}
\newcommand*{\Romanenumi}{\renewcommand*{\theenumi}{\Roman{enumi}}\relabel}
\newcommand*{\alphenumi}{\renewcommand*{\theenumi}{\alph{enumi}}\relabel}
\newcommand*{\Alphenumi}{\renewcommand*{\theenumi}{\Alph{enumi}}\relabel}
\let\oldtilde\tilde
\renewcommand*{\tilde}[1]{\mathchoice{\widetilde{#1}}{\widetilde{#1}}{\oldtilde{#1}}{\oldtilde{#1}}}
\let\oldexists\exists
\renewcommand*{\exists}{\mathrel{\oldexists}}
\let\oldforall\forall
\renewcommand*{\forall}{\mathrel{\oldforall}}

\newcommand{\llbracket}{\mathopen{[\mkern -2.7mu[}}
\newcommand{\rrbracket}{\mathclose{]\mkern -2.7mu]}}
\newcommand{\lb}{\llbracket}
\newcommand{\rb}{\rrbracket}
%%%%%%%%%%%%%%%%%%%%%%%%%%%%%%%%%%%%%%%%%%%%%%%%%%%%%%%%%%%%%%%%%%%%%%%%%%%%%%%%%%%%%%%%%%%%%%%%%%%%%%%%%%%%%%%%%%%%%%%%%%%%%%%%%%%%%%%%%

\title[Metric properties of boundary maps]{Metric properties of boundary maps, Hilbert entropy and non-differentiability}
\alttitle{Propriétés métriques des applications au bord à l'infini, entropie de Hilbert et non-différentiabilité}

\subjclass{22E40, 51F30, 37DXX}
\keywords{Anosov representations, Hausdorff dimension, higher rank Teichm\"uller Theory, Patterson--Sullivan measures}


\author[\initial{M. B.} \lastname{Pozzetti}]{\firstname{Maria Beatrice} \lastname{Pozzetti}}
\address{Università di Bologna\\ 
Dipartimento di Matematica\\
Piazza di Porta San Donato 5\\
40126 Bologna (Italy)}
\email{beatrice.pozzetti@unibo.it}


\author[\initial{A.} \lastname{Sambarino}]{\firstname{Andrés} \lastname{Sambarino}}
\address{Sorbonne Université\\
Institut de Mathématiques de Jussieu\\
Paris Rive Gauche (IMJ-PRG)\\ 
CNRS UMR 7586\\ 
4 place Jussieu\\
75005 Paris (France)}
\email{andres.sambarino@imj-prg.fr}

\thanks{B.~P. is supported by the Deutsche Forschungsgemeinschaft under Germany’s Excellence Strategy EXC-2181/1 - 390900948 (the Heidelberg STRUCTURES Cluster of Excellence), and acknowledges further support by DFG grant 338644254 (within the framework of SPP2026) and 427903332 (in the Emmy Noether program). A.~S. was partially financed by ANR DynGeo ANR-16-CE40-0025. Part of this work was carried out in Oberwolfach, we thank the institute for its great hospitality. B.~P. additionally thanks Prof. Farkas and Humboldt University for their hospitality, the Familienzimmer des mathematiches Institut made it possible to work on this paper while taking care of a~small child.}
\CDRGrant[EXC-2181/1]{390900948}
\CDRGrant[DFG]{338644254}
\CDRGrant[DFG]{SPP2026}
\CDRGrant[DFG]{427903332}
\CDRGrant[DynGeo]{ANR-16-CE40-0025}



\begin{abstract}
We interpret the Hilbert entropy of a convex projective structure on a closed higher-genus surface as the Hausdorff dimension of the non-differentiability points of the limit set in the full flag space $\mathcal{F}(\mathbb{R}^3)$. Generalizations for regularity properties of boundary maps between locally conformal representations are also discussed. An ingredient for the proofs is the concept of \emph{hyperplane conicality} that we introduce for a $\theta$-Anosov representation into a~reductive real-algebraic Lie group $\mathsf{G}$. In contrast with directional conicality, hyperplane-conical points always have full mass for the corresponding Patterson--Sullivan measure.
\end{abstract}

\begin{altabstract}
On interprète l'entropie de Hilbert d'une structure projective convexe sur une surface de genre minoré par 2 comme la dimension de Hausdorff des points de non-differentiabilité de l'ensemble limite dans l'espace des drapeaux complets $\mathcal{F}(\mathbb{R}^3)$. Des généralisations concernant des propriétés de régularité des applications limites au bord à l'infini pour des représentations localement conformes sont également traitées. Un ingrédient de la démonstration de ces résultats est le concept de \emph{point limite conique suivant un hyperplan} que l'on introduit pour une représentation $\theta$-anosovienne à valeurs dans un groupe réductif réel-algébrique $\mathsf{G}$. Contrairement aux points coniques suivant une direction, l'ensemble des points coniques suivant un hyperplan est toujours de masse totale pour la mesure de Patterson--Sullivan correspondante.
\end{altabstract}

\datereceived{2024-01-26}
\daterevised{2025-10-23}
\dateaccepted{2025-11-24}

\editors{C.~Sorger and V.~Colin}

%\begin{DefTralics}
%\renewcommand{\t}{\theta}
%\newcommand{\sfG}{\sf{G}}
%\newcommand{\R}{\mathbb{R}}
%\newcommand{\calF}{\mathcal{F}}
%\renewcommand{\sf}[1]{{\mathsf{#1}}}
%\end{DefTralics}
%%%%%%%%%%%%%%%%%%%%%%%%%%%%%%%%%%%%%%%%%%%%%%%%%%%%%%%%%%%%%%%%%%%%%%%%%%%%%%%%%%%%%%%%%%%%%%%%%%%%%%%%%%%%%%%%%%%%%%%%%%%%%%%%%%%%%%%%%
\setcounter{tocdepth}{1}
\setcounter{secnumdepth}{3}
\hypersetup{bookmarksdepth = 3}

\dateposted{2026-04-07}
\begin{document}
\maketitle
\tableofcontents

\section{Introduction}\label{intro}

Consider a closed connected orientable surface $S$ of genus at least two, and let $\rho:\piS\to\PSL(3,\R)$ be a faithful representation preserving an open convex set $\Omega=\Omega_\rho\subset\P(\R^3)$, properly contained in an affine chart. The group $\rho(\piS)$ is necessarily discrete and acts co-compactly on $\Omega$: one says that $\rho$ \emph{divides} $\Omega$.



The geometry of such convex set $\Omega$ is well studied, by Benoist~\cite{convexes1} it is strictly convex with $\class^{1+\nu}$ boundary $\bord\Omega$ (that is not $\class^2$ unless it is an ellipse), and the Hilbert metric of $\Omega$ is Gromov-hyperbolic. The geodesic flow of $\Omega/\rho(\piS)$ is an Anosov flow and its topological entropy, the \emph{Hilbert entropy} $\hJ{\sf H}={(\hJ{\sf H})}_\rho$, satisfies
\[
\hJ{\sf H}\leq1,
\]
an inequality proved by Crampon~\cite{crampon} that is strict if $\Omega$ is not an ellipse.



A consequence of Theorem~\ref{tutti} below is a new geometric interpretation of the Hilbert entropy which we now explain. For each $x\in\bord\Omega$ let $\Xi(x)\in\grass_{2}(\R^3)$ be the unique plane whose projectivization is tangent to $\bord\Omega$ at $x$. By~\cite{convexes1}, the image curve $\Xi(\bord\Omega)\subset\grass_2(\R^3)\simeq\P((\R^3)^*)$ is also the boundary of a strictly convex divisible set $\Omega^*$ and is thus again a $\class^{1+\nu}$-circle.
The \emph{full-flag-curve}
\[
\{(x,\Xi(x)):x\in\bord\Omega\}\subset\calF(\R^3),
\]
is the graph of a monotone map between $\class^1$ circles and thus is a Lipschitz submanifold that is therefore differentiable almost everywhere. We establish the following:



\begin{coro}\label{hilbert}
Let $\rho:\piS\to\PSL(3,\R)$ divide a strictly convex set that is not an ellipse. Then, the set of non-differentiability points of the full flag curve
has Hausdorff dimension ${(\hJ{\sf H})}_\rho$.
\end{coro}

Throughout the paper the Hausdorff dimension is computed with respect to a(ny) Riemannian metric on the flag space. When $\Omega$ is an ellipse the result does not apply as the associated curve is differentiable everywhere while $\hJ{\sf H}=1$.


A classical result by Choi--Goldman~\cite{choigoldman} states that the space of representations dividing a convex set forms a connected component of the character variety $\fkX(\piS,\PSL(3,\R))$ of homomorphisms up to conjugation. This component is known today as \emph{the Hitchin component} of $\PSL(3,\R)$ and is diffeomorphic to a ball of dimension $-8\chi(S)$. Nie~\cite{Nie0} and Zhang~\cite{ZhangSL3} have found paths $(\rho_t)$ in this Hitchin component such that $({\hJ{\sf H}})_{\rho_t}\to 0$ as $t\to\infty$. Together with Corollary~\ref{hilbert} this suggest that the closer $\Omega$ is to being an ellipse (the \emph{Fuchsian locus}), the less differentiable the flag curve is whilst the furthest away from Fuchsian locus, the more regular the flag curve becomes.



The proof of Corollary~\ref{hilbert} is outlined in Section~\ref{corAProof} and serves as a guide path for the strategy on the general case (Theorems~\ref{LCdiff} and~\ref{tutti}).


\subsection{Locally conformal representations and concavity properties}\label{2hyper}

Let $\K$ be $\R$, $\C$ or the non-commutative field of Hamilton's quaternions $\H$. Denote by
\[
\a=\left\{(a_1,\dots,a_d)\in\R^d:\sum_i a_i=0\right\}
\]
the Cartan subspace of the real-algebraic group $\SL(d,\K)$, by
\begin{equation}\label{e.slroot}
\slroot_i(a_1,\dots,a_d)=a_i-a_{i+1}
\end{equation}
the $i^{\, \rm th}$ simple root and by $\a^+\subset \a$ the Weyl chamber whose associated set of simple roots is $\simple=\{\slroot_i:i\in\lb1,d-1\rb\}$. Let $\cartan:\SL(d,\K)\to\a^+$ be the \emph{Cartan projection} with respect to the choice of an inner (or Hermitian) product on $\K^d$. The $e^{\cartan_i(g)}$'s are the \emph{singular values} of the matrix $g$, namely the square roots of the modulus of the eigenvalues of the matrix $gg^*$. We also let $d_\P$ denote the distance on $\P(\K^d)$ induced by the chosen Hermitian product.


Let $\G$ be a finitely generated word-hyperbolic group, consider a finite symmetric generating set and let $|\,|$ be the associated word-length. For $k\in\lb1,d-1\rb$, a representation $\rho:\G\to\SL(d,\K)$ is $\{\slroot_k\}$-\emph{Anosov} if there exist positive constants $\mu$ and $c$ such that for all $\g\in\G$ one has
\[
\slroot_k\bigl(\cartan(\rho(\g))\bigr)\geq\mu|\g|-c.
\]

A $\{\slroot_k\}$-{Anosov} representation is also $\{\slroot_{d-k}\}$-Anosov. Under such assumption there exists an equivariant H\"older-continuous map
\[
\xi^{k}_\rho:\bord\G\to\grass_k(\K^d),
\]
called the \emph{limit map} in the Grassmannian $\grass_k(\K^d)$ of $k$-dimensional subspaces of~$\K^d$, which is a homeomorphism onto its image. If $k\leq l\in\lb1,d-1\rb$ and $\rho$ is also\linebreak $\{\slroot_l\}$-Anosov then the limit maps are {compatible}, i.e. $\xi^k_\rho(x)\subset\xi^l_\rho(x)$ $\forall x$, see Section~\ref{Anosov} for references and details.

\begin{defi}\label{d.hyp}
Fix $p\in\lb2,d-1\rb$. A $\{\slroot_1,\slroot_{d-p}\}$-Anosov representation\linebreak $\rho:\G\to\SL(d,\K)$ is \emph{$(1,1,p)$-hyperconvex} if for every pairwise distinct triple $x,y,\linebreak z\in\bord\G$ one has
\begin{equation}\label{hypdefi}
\bigl(\xi^{1}_\rho(x)+\xi^{1}_\rho(y)\bigr)\cap\xi^{{d-p}}_\rho(z)=\{0\}.
\end{equation}
If in addition one has $\cartan_2(\rho(\g))=\cartan_{p}(\rho(\g))$ $\forall\g$, we say that $\rho$ is \emph{locally conformal}.
\end{defi}

Hyperconvex representations form an open subset of the character variety
\[
\fkX\bigl(\G,\SL(d,\K)\bigr)=\hom\bigl(\G,\SL(d,\K)\bigr)/\SL(d,\K)
\]
and appear naturally. For example, when $\K=\R$, strictly convex divisible sets give rise to $(1,1,d-1)$-hyperconvex representations, while higher rank Teichm\"uller theory provides many examples of $(1,1,2)$-hyperconvex representations of surface groups, see Example~\ref{ex.hyperconvex}.


When $p=2$ the second part of the definition is trivially true, so $(1,1,2)$-hypercon\-vex representations over $\K$ \emph{are} locally conformal, when $p>2$ the assumption constrains the Zariski closure of $\rho(\G)$. However, Zariski-dense locally conformal representations exist (and form open sets) for the groups locally isomorphic to $\SL(n,\R)$, $\SL(n,\C)$, $\SL(n,\H)$, $\SU(1,n)$, $\Sp(1,n)$, $\SO(p,q)$, see~\cite[Section~8]{PSW1} %P.-S.-Wienhard
for details, and, of course, $\SO(1,n)$ where every convex co-compact representation is locally conformal.

A concrete example in $\SU(1,n)$ consists on considering a convex co-compact group in $\H^n_\C$ whose limit set intersects the projectivization of any complex line in at most 2~points. These subgroups are locally conformal (\cite[Proposition~8.3]{PSW1}) and their limit set (though fractal) is tangent to the contact distribution of $\bord\H^n_\C$.



Consider also $\bar{\K}\in\{\R,\C,\H\}$ and positive integers $d$ and $\ovd$. Throughout the paper we mainly deal with a pair of locally conformal representations
\[
\rho:\G\to\SL(d,\K)\quad\text{and}\quad\ovrho:\G\to\SL(\ovd,\bar{\K}),
\]
with equivariant maps $\xi=\xi^1_\rho$ and $\ovxi=\xi^1_{\ovrho}$, and we study regularity properties of the equivariant H\"older-continuous homeomorphism
\[
\Xi=\ovxi\circ\xi^{-1}:\xi(\bord\G)\to\ov{\xi}(\bord\G).
\]

To avoid confusion we denote the simple roots of $\SL(\ovd,\bar{\K})$ by $\{\ov{\slroot}_i:i\in\lb1,\ovd-1\rb\}$, and to simplify notation we identify $\g$ with $\rho(\g)$ and we let $\ov\g=\ovrho(\g)$. We consider also the graph of $\Xi$, or equivalently the \emph{graph map},
\[
\varXi:\bigl(\xi,\ov{\xi}\bigr):\bord\G\to\P(\K^d)\times\P(\bar{\K}{}^{\ovd}).
\]

\pagebreak{}
\begin{defi}
Fix $\hol\in(0,1]$. We will say that $\Xi$ is $\hol$\emph{-concave} at $x\in\bord\G$, or that $x$ is a $\hol$\emph{-concavity point for} $\Xi$, if there exists a sequence $(y_k)$ converging to $x$ as $k\to\infty$ such that the incremental quotient
\begin{equation}\label{incri}
\frac{d_\P\bigl(\ovxi(x),\ovxi(y_k)\bigr)}{d_\P\bigl(\xi(x),\xi(y_k)\bigr)^\hol}
\end{equation}
is bounded away from $\{0,\infty\}$. The set of $\hol$-concavity points is denoted by $\Ext_{\rho,\ovrho}^\hol$.
\end{defi}

Observe that $\Xi$ can be $\hol$-concave at $x$ for several $\hol$'s and that it is a $1$-concave point if one has $y_k\to x$ such that $d(\xi(x),\xi(y_k))$ and $d_\P(\ovxi(x),\ovxi(y_k))$ are comparable.

\begin{figure}[h]
\centering
\begin{tikzpicture}
\begin{scope}[scale=1.5]
\begin{scope}[rotate=-90]
\draw[name path=Q] (0,0) parabola bend (0,0) (-1.5,1.5);
\fill[color=gray!20] (0,0) parabola bend (0,0) (-1.5,1.5)--(-1,1.5)--cycle;
\fill[color=white] (0,0) parabola bend (0,0) (-1,1.5)--cycle;
\draw (0,0) parabola bend (0,0) (-1,1.5);
\draw[name path=Q] (0,0) parabola bend (0,0) (-1.5,1.5);
\node at (-1.5,1.5) [above right] {$Cs^\hol$};
\node at (-1,1.5) [below right] {$cs^\hol$};
\end{scope}

\draw[gray!60] (0.8,0)--(0.8,0.8)--(0,0.8);
\draw[gray!60] (0.4,0)--(0.4,0.67)--(0,0.67);
\draw[gray!60] (0.1,0)--(0.1,0.3)--(0,0.3);
\draw[gray!60] (0.05,0)--(0.05,0.215)--(0,0.215);

\draw (0,0)--(2,0);
\node[below] at (2,0) {$d_\P(\xi(x),\xi(z))$};
\draw(0,0)--(0,2);
\node[left] at (0,2) {$d_\P(\ovxi(x),\ovxi(z))$};

\draw [fill] (0.8,0) circle [radius = 0.02];
\draw [fill] (0.8,0.8) circle [radius = 0.02];
\draw [fill] (0,0.80) circle [radius = 0.02];


\draw [fill] (0.4,0) circle [radius = 0.02];
\draw [fill] (0.4,0.67) circle [radius = 0.02];
\draw [fill] (0,0.67) circle [radius = 0.02];


\draw [fill] (0.1,0.3) circle [radius = 0.017];
\draw [fill] (0.1,0) circle [radius = 0.017];
\draw [fill] (0,0.3) circle [radius = 0.017];

\draw [fill] (0.05,0) circle [radius = 0.008];
\draw [fill] (0.05,0.215) circle [radius = 0.008];
\draw [fill] (0,0.215) circle [radius = 0.008];

\begin{scope}[rotate=-90]
\draw (0,0) parabola bend (0,0) (-1.5,1.5);

\draw (0,0) parabola bend (0,0) (-1,1.5);
\end{scope}
\end{scope}
\end{tikzpicture}
\caption{A $\hol$-concave point $x$. The marked points on the axis' represent $d_\P(\xi(x),\xi(y_k))$ and $d_\P(\ovxi(x),\ovxi(y_k))$ respectively.}
\end{figure}



In what follows we will compute the Hausdorff dimension of $\varXi(\Ext_{\rho,\ovrho}^\hol)$ with respect to the product metric on $\P(\K^d)\times\P(\bar{\K}{}^{\ovd})$ for $\hol$ lying on an interval that we now define. The \emph{dynamical intersection} between $\rho$ and $\ovrho$ with respect to $\ovsloot_1$ and $\slroot_1$ is defined by
\[
\II_{\ovsloot_1}(\slroot_1)=\lim_{t\to\infty}\frac1{\#\sf R_t(\ovsloot_1)}\sum_{\g\in \sf R_t(\ovsloot_1)}\frac{\slroot_1(\lambda(\g))}{\ovsloot_1(\lambda(\ov\g))},
\]
where $\sfR_t(\ovsloot_1)=\{[\g]\in[\G]:\ovsloot_1(\lambda(\ov\g))\leq t\}$ and $\lambda:\SL(d,\K)\to\a^+$ is the Jordan projection. This concept (from Bridgeman et al.~\cite{pressure},
%--Canary--Labourie--Sambarino
Burger~\cite{burger}, Knieper~\cite{kni95}, among others) generalizes Bonahon's intersection number between two elements in Teichm\"uller space.

Let us say that $\rho$ and $\ovrho$ are \emph{gap-isospectral} if for all $\g\in\G$ one has
\[
\slroot_1\bigl(\lambda(\g)\bigr)=\ovsloot_1\bigl(\lambda(\ov\g)\bigr).
\]
Corollary~\ref{corraiz} (a consequence of~\cite{pressure} together with Proposition~\ref{nonA}) implies that if $\rho$ and $\ovrho$ are not gap-isospectral, then $\II_{\slroot_1}(\ovsloot_1)>(\II_{\ovsloot_1}(\slroot_1))^{-1}$. We will study $\hol$-concavity for any $\hol\in(0,1]$ with
\[
\II_{\slroot_1}(\ovsloot_1)>\hol>\bigl(\II_{\ovsloot_1}(\slroot_1)\bigr)^{-1}.
\]

Finally, consider the critical exponents
\begin{alignat*}{2}\hC{\slroot_1} & =\lim_{t\to\infty}\frac1t\log\#\bigl\{\g\in\G:\slroot_1(\cartan(\g))\leq t\bigr\},\\
\hC{\infty,\hol} & =\lim_{t\to\infty}\frac1t\log\#\bigl\{\g\in\G:\max\bigl\{\hol\slroot_1(\cartan(\g)),\ov{\slroot}_1(\cartan(\ov\g))\bigr\}\leq t\bigr\}.
\end{alignat*}

\begin{theo}[Theorem~\ref{tutti.LCdiff}]\label{LCdiff}
Let $\{\K,\bar{\K}\}\subset\{\R,\C\}$ and let $\rho:\G\to\SL(d,\K)$ and $\ovrho:\G\to\SL(\ovd,\bar{\K})$ be locally conformal, $\R$-irreducible and not gap-isospectral. Then for any $\hol\in(0,1]$ with $\II_{\slroot_1}(\ovsloot_1)>\hol>(\II_{\ovsloot_1}(\slroot_1))^{-1}$, one has
\begin{align*}
\hol\hC{\infty,\hol}\leq\Hff\left(\varXi\left({\Ext}_{\rho,\ovrho}^\hol\right)\right)
&\leq\min\left\{\hC{\infty,\hol},\hol\hC{\infty,\hol}+1-\hol \right\}\\
&<\min\left\{\hC{\ovsloot_1},\hC{\slroot_1}/\hol\right\}\\
&\leq\Hff(\varXi(\bord\G))\\
&=\max\left\{\hC{\slroot_1},\hC{\ovsloot_1}\right\}.
\end{align*} If $\K=\H$ (resp. $\bar{\K}=\H$) we further assume that the Zariski closure if $\rho$ (resp. $\ovrho$) does not have compact factors, then the same conclusion holds.
\end{theo}

The proof of the above Theorem is completed in Section~\ref{ProofA}. For representations in $\PSL(2,\C)$ we can furthermore give a geometric interpretation of the 1-weakly-bi-H\"older points, see Section~\ref{klein}.

\subsection{Surface-group representations}

Observe that the first line of inequalities in Theorem~\ref{LCdiff} becomes an equality when $\hol=1$. We pursue now this situation while further restricting the source and ambient groups.

Let then $\K=\R$ and assume $\bord\G$ is homeomorphic to a circle. Real representations of $\G$ that are $(1,1,2)$-hyperconvex are necessarily locally conformal and form the prototype example of Anosov representations with $\class^1$ limit sets: indeed we have the following result from Pozzetti et al.~\cite{PSW1} and Zhang--Zimmer~\cite{ZZ}.

\begin{theo}\label{t.C1}
Assume $\bord\G$ is homeomorphic to a circle and let $\rho:\G\to\PGL(d,\R)$ be $\{\slroot_1\}$-Anosov.
\begin{itemize}
\item \cite{PSW1,ZZ}:  If $\rho$ is $(1,1,2)$-hyperconvex, then $\xi^1(\bord\G)\subset\P(\R^{d})$ is\linebreak a~$\class^1$-submanifold tangent at $\xi^1(x)$ to $\xi^2(x)$.
\item {\cite{ZZ}}: If $\rho$ is irreducible and $\xi(\bord\G)$ is a $\class^{1+\alpha}$ circle then $\rho$ is $(1,1,2)$-hyperconvex.
\end{itemize}
\end{theo}



The graph map $\varXi=(\xi,\ov{\xi}):\bord\G\to\P(\R^{d})\times\P(\R^{\ovd})$ has image contained in the $\class^{1+\nu}$ torus $\xi(\bord\G)\times\ov{\xi}(\bord\G)$ and $\varXi(\bord\G)$ is the graph of $\Xi$, a H\"older-continuous homeomorphism between $\class^{1+\nu}$-circles. By monotonicity of $\Xi$, $\varXi(\bord\G)$ is a Lipschitz curve and is thus differentiable almost everywhere. We let
\[
\Ndiff_{\rho,\ovrho}\subset\varXi(\bord\G)
\]
be the subset of points where the curve $\varXi(\bord\G)$ is not differentiable. The combination of Lemma~\ref{generalcase} and Corollary~\ref{nondiffbconicalR} establishes that in the current situation (with mild additional assumptions)
\[
\varXi\bigl(\Ext^1_{\rho,\ovrho}\bigr)=\Ndiff_{\rho,\ovrho},
\]
whence with Theorem~\ref{LCdiff} one obtains the following:

\begin{theo}\label{tutti}
Assume $\bord\G$ is homeomorphic to a circle and let $\rho:\G\to\SL(d,\R)$ and $\ovrho:\G\to\SL(\ovd,\R)$ be $(1,1,2)$-hyperconvex and not gap-isospectral. Then,
\[
\Hff\big(\Ndiff_{\rho,\ovrho}\big)=\hC{\infty,1}<1.
\]
\end{theo}

We emphasize that no irreducibility assumption is made on the representations $\rho$ and $\ovrho$. On the other hand, if the representations are irreducible and gap-isospectral, we show that there exists an isomorphism between the Zariski closures of $\rho(\G)$ and of~$\ovrho(\G)$ intertwining the two representations. It follows then that $\varXi(\bord\G)$ is the diagonal of the $\class^{1+\nu}$ torus, and thus differentiable everywhere. To prove this we give the following preliminary classification of Zariski-closures, established in Section~\ref{s.Zd}.

Recall that if $\sfG$ is a semi-simple real-algebraic group of non-compact type, then irreducible proximal representations $\Fund:\sfG\to\PGL(V)$ are determined by their highest restricted weight $\chi_\Fund$. A special subset of dominant weights are the so-called \emph{fundamental weights} $\{\peso_\sroot:\sroot\in\simple\}$, and are indexed by the set of simple roots $\simple$ of~$\sfG$ (see Section~\ref{representaciones} for definitions and details).

\begin{theo}\label{t.Zcl}
Assume $\bord\G$ is homeomorphic to a circle and let $\rho:\G\to\PGL(d,\R)$ be irreducible and $(1,1,2)$-hyperconvex. Then the Zariski closure $\sfG$ of $\rho(\G)$ is simple and the highest weight of the induced representation $\Fund:\sfG\to\PGL(d,\R)$ is a multiple of a fundamental weight associated to a root whose root-space is one-dimensional.
\end{theo}


In light of the following examples it is unclear if further restrictions can occur.



\begin{exam}\label{ex.hyperconvex}
Any pair of representations $\rho:\pi_1S\to\sfG$ and $\Fund:\sfG\to\PGL(V)$ in each of the following classes (and small deformations), gives rise to a $(1,1,2)$-hyperconvex representation via post-composition $\Fund\circ\rho$. In particular the limit set of $\rho$ in the specified flag manifold of $\sfG$ is a $\class^{1+\nu}$ curve:
\begin{itemize}
\item $\sfG$ is split, $\rho:\piS\to\sfG$ is Hitchin, and $\Fund$ satisfies $\chi_\Fund=n\peso_\sroot$ for any $\sroot\in\simple$ and $n\in\N_{>0}$. This is non-trivial and requires results from Fock--Goncharov~\cite{FG} and Labourie~\cite{labourie} together with Lusztig's canonical basis~\cite[Proposition~3.2]{Lusztig-TP} (see Sambarino~\cite[Section~5.8]{clausurasPos} for details). As a result the limit set of $\rho$ in any maximal flag manifold $\calF_{\{\sroot\}}$ of $\sfG$ is a $\class^{1+\nu}$ curve.

\item $\rho:\piS\to\PO(p,q)$ is $\Theta$-positive and $\Fund$ has highest weight $\peso_\sroot$ for any root $\sroot$ in the interior\footnote{I.e. $\sroot$ is only connected to roots in $\Theta$ in the Dynkin diagram of $\simple$.} of $\Theta$ (\cite[Theorem~10.3]{PSW2}, see also Beyrer--Pozzetti~\cite[Remark~4.6]{BP}). In particular the limit set in any flag manifold of the form ${\sf{Is}}_k(\R^{p,q})$ for $k\leq p-2$ is a $\class^{1+\nu}$-curve. When $\rho$ is moreover Zariski-dense, we can consider any $\Fund$ with $\chi_\Fund^+=n\peso_\sroot$ for any $\sroot \in\inte\Theta$ and $n\in\N_{>0}$.
\item For all $k\geq 1$, $k$-positive representations $\rho:\piS\to\PSL(d,\R)$ introduced in~\cite{BeyP2} are $(1,1,2)$-hyperconvex.
\end{itemize}
\end{exam}



For these examples also the following applies:


\begin{coro}\label{cor.B}
Assume $\bord\G$ is homeomorphic to a circle, let $\sfG$ be a simple Lie group and let $\rho:\G\to\sfG$ have Zariski-dense image. Assume there exist $\{\sroot,\bb\}\subset\simple$ distinct such that both $\Fund_\sroot\circ\rho$ and $\Fund_\bb\circ\rho$ are $(1,1,2)$-hyperconvex. Then:
\begin{enumerate}
\item\label{coro1.9.1} The image of $\xi^{\{\sroot,\bb\}}:\bord\G\to\calF_{\{\sroot,\bb\}}$ is Lipschitz and the Hausdorff dimension of the points where it is non-differentiable is $\hC{\max\{\sroot,\bb\}}$.
\item\label{coro1.9.2} If the opposition involution $\ii$ on $\ge$ is non-trivial and $\bb=\ii\sroot$ then
\[
\hC{\max\{\sroot,\bb\}}=\hC{(\sroot+\bb)/2}.
\]
\end{enumerate}
\end{coro}



\begin{rema}
A different approach to Theorem~\ref{tutti}, relying on Theorems~\ref{t.Zcl} and~\ref{t.C1}, would be to code the action of $\piS$ on $\bord\piS$ via Bowen-series and apply Jordan et al.~\cite[Theorem~1.1]{JKPS}. This method, followed by Pollicott--Sharp~\cite{PS} for two representations in the Teichm\"uller space of $S$, is not applicable for groups other than $\piS$, in particular this approach cannot be used in the generality of Theorem~\ref{LCdiff}.
\end{rema}


\subsection{Hyperplane vs directional conicality}

To prove Theorems~\ref{LCdiff} and~\ref{tutti} we introduce the concept of \emph{hyperplane conicality}, a~generalization of directional conicality from Burger et al.~\cite{BLLO}.

Let $\sfG$ be a real-algebraic semi-simple Lie group of non-compact type, $\a\subset\ge$ a~Cartan subspace, $\roots\subset\a^*$ the associated root system and $\simple\subset\roots$ a choice of simple roots with associated Weyl chamber $\a^+$.


Consider a non-empty $\t\subset\simple$ and let $\a_\t$ be the associated Levi space. Fix a $\t$-Anosov representation $\rho:\G\to\sfG$ and denote by $\cal L_{\t,\rho}\subset\a_\t$ its $\t$-limit cone. We will recall in Section~\ref{rfr} that, when $\rho(\G)$ is Zariski-dense, there are natural bijections
\begin{align*}
\inte\P(\cal L_{\t,\rho})& \leftrightarrow\calQ_{\t,\rho}=
\left\{\varphi \in(\a_\t)^*:\hJ\varphi=1\right\} \\
&\leftrightarrow \left\{\text{Patterson--Sullivan measures supported on $\xi^\t(\bord\G)$}\right\}.
\end{align*}
For $\varphi\in\calQ_{\t,\rho}$ we let $\sf u_\varphi\in \inte\P(\cal L_{\t,\rho})$ be the associated direction and $\mu^\varphi$ the associated Patterson--Sullivan measure.


Consider now a hyperplane $\sf W\subset\a_\t$ and assume, for the notion to be interesting, that $\sf W$ intersects the relative interior of $\calL_{\t,\rho}$. Then $x\in\bord\G$ is \emph{$\sf W$-conical} if there exists a conical sequence $(\g_n)_0^\infty\subset\G$ converging to $x$, a constant $K$ and a sequence $(w_n)_0^\infty\in\sf W$ such that for all $n$ one has
\[
\bigl\|\cartan_\t(\rho(\g_n))- w_n\bigr\|\leq K,
\]
where $\cartan_\t:\sfG\to\a^+_\t$ is the $\t$-Cartan projection. The set of $W$-conical points will be denoted by $\bord_{\sf W,\rho}\G=\bord_{\sf W}\G$. Inspired by~\cite{BLLO}, in Theorem~\ref{ps-sandwich} we show the following.



\begin{theo}\label{Ahyper}
Let $\rho:\G\to\sfG$ be a Zariski-dense $\t$-Anosov representation and $\sf W$ be a hyperplane of $\a_\t$ intersecting non-trivially the interior of $\calL_{\t,\rho}$. Then for every $\varphi\in\calQ_{\t,\rho}$ with $\sf u_\varphi\in\P(\sf W)$ one has $\mu^\varphi(\bord_\sf W\G)=1$.
\end{theo}



\subsection{Strategy of the proof of Corollary~\ref{hilbert}}\label{specialcases}
\label{corAProof}
Corollary~\ref{hilbert} is a consequence of Theorem~\ref{tutti} where $\ovrho$ is the dual representation of~$\rho$. We sketch a direct proof of Corollary~\ref{hilbert} serving as a guide-path for the general result.


Let $\rho:\piS\to\SL(3,\R)$ be the holonomy of a strictly convex projective structure dividing the convex set $\Omega$. We consider the $L^\infty$ distance on the product $(\P(\R^3), d_\P)\times\linebreak(\P((\R^3)^*),d_{\P})$, which is equivalent to the Riemannian distance, and thus induces the same Hausdorff dimension.


As a replacement of Sullivan's shadows we use \emph{coarse cone type at infinity}, inspired by Cannon's work on \emph{cone types}~\cite{CannonCones} (see also Section~\ref{cont}). Fix a finite symmetric generating set on $\piS$ and let $|\,|$ be the associated word length. For $\g\in\piS$ and $c>0$, the \emph{coarse cone type at infinity} $\cone^{c}_\infty(\g)$ of $\gamma$ is the set of endpoints at infinity of $(c,c)$-quasi geodesic rays based at $\g^{-1}$ passing through the identity. See Figure~\ref{figure:conetype}.

\begin{figure}[hh]
\centering
\begin{tikzpicture}

\begin{scope}[scale=0.4, rotate=75]
\draw circle [radius=.5];
\draw [thick] circle [radius = 3];
\draw [fill] circle [radius = 0.03];
\draw[thick] (-0.1,-0.3) -- (0.5,0.1);
\node [below left] at (0,0) {$\g$};
\node [below left] at (0,1.2) {$\g\eta$};
\node [above] at (1.2,3.2) {$\g\cone^{c}_{\infty}(\g)$};
\draw [thick] (0.004,-1.3) -- (0.1,-0.7) -- (-0.1,-0.3) -- (0,0);
\draw [thick] (0,0) -- (-0.3,1) -- (0.2,2) -- (0,3);
\draw [thick] (0.2,2) -- (0.4, 2.7) -- (0.7,2.91);
\draw [thick] (0.2,2) -- (0.4, 2.7);
\draw [thick] (1,0.2) -- (0.89,1.37);
\draw [thick] (1.8,0.8) -- (2.7,1.3);
\draw [thick] (0,0) -- (1,0.2) -- (1.8,0.8) -- (2.846,0.948);
\node [below] at (0.004,-1.3) {$e$};
\draw [fill] (0.7,2.91) circle [radius = 0.03];
\draw [fill] (2.7,1.3) circle [radius = 0.03];
\draw [fill] (0.004,-1.3) circle [radius = 0.03];
\draw [fill] (2.12,2.12) circle [radius = 0.03];
\draw [fill] (0,3) circle [radius = 0.03];
\draw [fill] (2.846,0.948) circle [radius = 0.03];
\draw [thick] (0,0) -- (0.89,1.37) -- (1.2,2) -- (2.12,2.12);
\draw [thick] (1.2,2) -- (1.4,2.65);
\draw [fill] (1.4,2.65) circle [radius = 0.03];
\node [right] at (3,-2) {$\piS$};
\filldraw [red](0,3) circle [radius=2pt];
\filldraw [red](0.7,2.91) circle [radius=2pt];
\filldraw [red](2.846,0.948) circle [radius=2pt];
\filldraw [red](2.7,1.3) circle [radius=2pt];
\filldraw [red](2.12,2.12) circle [radius=2pt];
\filldraw [red] (1.4,2.65) circle [radius=2pt];
\end{scope}

\begin{scope}[scale=0.4,shift = {(-12,0)}]
\draw [thick] circle [radius = 3];
\draw circle [radius=.75];
\draw [fill] circle [radius = 0.03];
\node [left] at (0,0) {$e$};
\draw[thick] (-0.1,-0.3) -- (0.5,0.1);
\draw [thick] (0.004,-1.3) -- (0.1,-0.7) -- (-0.1,-0.3) -- (0,0);
\draw [thick] (0,0) -- (-0.3,1) -- (0.2,2) -- (0,3);
\draw [thick] (0.2,2) -- (0.4, 2.7) -- (0.7,2.91);
\draw [thick] (1,0.2) -- (0.89,1.37);
\draw [thick] (0.2,2) -- (0.4, 2.7);
\draw [thick] (1.8,0.8) -- (2.7,1.3);
\draw [thick] (0,0) -- (1,0.2) -- (1.8,0.8) -- (2.846,0.948);
\node [below] at (0.004,-1.3) {$\g^{-1}$};
\draw [fill] (0.7,2.91) circle [radius = 0.03];
\draw [fill] (2.7,1.3) circle [radius = 0.03];
\draw [fill] (0.004,-1.3) circle [radius = 0.03];
\draw [fill] (2.12,2.12) circle [radius = 0.03];
\draw [fill] (0,3) circle [radius = 0.03];
\draw [fill] (2.846,0.948) circle [radius = 0.03];
\draw [thick] (0,0) -- (0.89,1.37) -- (1.2,2) -- (2.12,2.12);
\draw [thick] (1.2,2) -- (1.4,2.65);
\draw [fill] (1.4,2.65) circle [radius = 0.03];
\node [right] at (3,-2) {$\piS$};
\draw [->] (3.5,3) to [out = 180,in = 90] (1.5,1.2);
\node [right] at (3.5,3) {$\cone^{c}(\g)$};
\end{scope}
\end{tikzpicture}
\caption{The coarse cone type of $\g\in\G$ (left). The set $\g\cdot\cone^{c}_\infty(\g)$ (right). Pictures from~\cite{PSW1}. %Pozzetti--Sambarino--Wienhard
}\label{figure:conetype}
\end{figure}



We let $\xi:\bord\piS\to\bord\Omega$ be the natural identification via the action of $\rho(\piS)$ on~$\Omega$, and analogously $\ovxi:\piS\to\bord\Omega^*$. We denote by $\varXi:=(\xi,\ovxi): \piS\to\bord\Omega\times\bord\Omega^*$ the flag curve. Consider $x\in\bord\piS$ and let $\alpha_i\to x$ be a geodesic ray on $\piS$. The following fact is a consequence of Proposition~\ref{conetypesBalls}.

\begin{enonce}
{Fact} For big enough $i$, the subset $\xi(\alpha_i\cone^{c}_\infty(\alpha_i))\subset\bord\Omega$ is coarsely the intersection of a ball of radius $e^{-\slroot_1(\alpha_i)}$ about $\xi(x)$ with $\bord\Omega$. By duality, one has $\ovxi(\alpha_i\cone^{c}_\infty(\alpha_i))\subset\bord\Omega^*$ is coarsely the intersection of a ball of radius $e^{-\slroot_2(\alpha_i)}$ about $\ovxi(x)$ with $\bord\Omega^*$.
\end{enonce}

The coarse constants and the minimal length $i$ required in the above statement depend only on the representation and not on the point $x$.



For any point $x\in\bord\piS$ we distinguish two complementary situations that don't depend on the choice of the geodesic ray $(\alpha_i)_{i\in\N}$ converging to $x$:
\begin{enumerate}\romanenumi
\item For all $R>0$ there exists $N\in\N$ with $|\slroot_1(\cartan(\alpha_i))-\slroot_2(\cartan(\alpha_i))|\geq R$ for all $i\geq N;$
\item There exists $R>0$ and an infinite set of indices $\I\subset\N$ such that for all $k\in\I$ one has $|\slroot_1(\cartan(\alpha_k))-\slroot_2(\cartan(\alpha_k))|\leq R$. We say in this case that $x$ is $\flat$\emph{-conical} ($\flat$ stands for 'barycenter of the chamber').
\end{enumerate}


\begin{figure}[!h]
\centering
\begin{tikzpicture}


\begin{scope}[rotate=0,scale=0.7,shift={(-3,-2)}]


\draw (0,0) rectangle (6,4);
\fill (3,0) circle (2pt);
\fill (0,2) circle (2pt);
\fill (3,2) circle (2pt);



\draw[name path=c1,-] (0,0).. controls (1,2) and (2,1.7).. (3,2);
\draw[name path=c2,-] (3,2).. controls (4.5,2.5) and (4.5,3).. (6,4);


\draw[name path=l1,white] (2,0.2)--(2,3.8);
\draw[name path=l2,white] (4,0.2)--(4,3.8);

\draw[name intersections={of=c1 and l1, by={a}}, name intersections={of=c2 and l2, by={b}}] (a) rectangle (b);


\node [above left] at (b) {\footnotesize{$e^{-\slroot_1(\rho\alpha_i)}$}};
\node [below right] at (b) {\footnotesize{$e^{-\slroot_2(\rho\alpha_i)}$}};



\begin{scope}[shift={(3,2)}]
\end{scope}

\node [above right] at (3,0) {$\footnotesize{\xi(x)}$};
\node [above right] at (0,2) {$\footnotesize{\ovxi(x)}$};
\end{scope}
\end{tikzpicture}
\caption{The image of the cone type $\alpha_i\cone^{c}_\infty(\alpha_i)$ by the graph curve $\varXi$ in the $\class^{1+\nu}$-torus $\bord\Omega\times\bord\Omega^*$.}\label{figg}
\end{figure}

In the first case one is easily convinced by looking at Figure~\ref{figg} that the rectangle becomes flatter along one of its sides (see Section~\ref{s.Hffproof} for details in the general case). Furthermore, since $\slroot_1(\cartan(\alpha_i))-\slroot_1(\cartan(\alpha_{i+1}))$ is uniformly bounded, its sign is eventually constant, and thus the longer side only depends on the point. As a result $x$ is necessarily a differentiability point of the graph curve $\varXi$, with either horizontal or vertical derivative.

We are thus bound to understand the set of $\flat$-conical points. We show (see Corollary~\ref{nondiffbconicalR}):

\begin{enonce}{Fact}
The non-differentiability points of the curve $\varXi(\bord\piS)$ and the $\flat$-conical points coincide.
\end{enonce}

The main idea for this is to show that if a $\flat$-conical point $x$ were a differentiability point, then the derivative could not be horizontal nor vertical, and thus (by Proposition~\ref{c.indepPer}) $\Xi$ would be bi-Lipschitz. In turn, this would force the periods of the two roots to agree, which in turn would imply that the representation is Fuchsian, contradicting the assumption that $\Omega$ is not an ellipse.

It remains to understand the Hausdorff dimension of the set of $\flat$-conical points. The upper bound (Proposition~\ref{upper})
\begin{equation}\label{u}
\Hff(\{\flat-\text{conical}\})\leq\hC{\max\{\slroot_1,\slroot_2\}}
\end{equation}
follows readily: since for a $\flat$-conical point the lengths $e^{-\slroot_1(\alpha_k)}$ and $e^{-\slroot_2(\alpha_k)}$ are comparable independently on $k\in\I$, one can replace the rectangle in Figure~\ref{figg} by the (smaller) square of length
\[
e^{-\max\{\slroot_1(\cartan(\alpha_k)),\slroot_2(\cartan(\alpha_k))\}}
\]
and still get a covering\footnote{Choosing the longer side $e^{-\min\{\slroot_1(\cartan(\alpha_k)),\slroot_2(\cartan(\alpha_k))\}}$ gives the bound {$\Hff\varXi(\bord\piS)\leq 1$.}} (this time by balls on the $L^\infty$ metric) of the set $\{\flat-\text{conical}\}$. Standard arguments on Hausdorff dimension give Eq.~\eqref{u}.



Finding a lower bound for the Hausdorff dimension is more subtle; we use here an appropriate Patterson--Sullivan measure to study how the mass of a ball of radius $r$ scales with $r$.

Since $\varXi(\bord\piS)$ is a subset the full flag space $\calF(\R^3)$ and
\[
\|v\|_\infty:=\max\{|\slroot_1(v)|,|\slroot_2(v)|\}
\]
is a norm on $\a_{\PSL(3,\R)}$, we can apply results by Quint~\cite{Quint-Div} to determine a linear form $\varphi^\infty_\flat\in\a^*$ whose associated growth direction is the barycenter $\scr{b}=\ker(\slroot_1-\slroot_2)$. By~\cite[Proposition~3.3.3]{Quint-Div}
\[
\hC{\max\{\slroot_1,\slroot_2\}}=\|\varphi^\infty_\flat\|^1,
\]
where $\|\,\|^1$ is the operator norm on $\a^*$ defined by $\|\,\|_\infty$, which turns out to be the $L^1$ norm $\|a\slroot_1+b\slroot_2\|^1=|a|+|b|$. The form $\varphi^\infty_\flat$ additionally admits an associated \emph{Patterson--Sullivan} probability measure, namely a measure $\ps^\infty$ such that for all $\g\in\piS$ one has (see Corollary~\ref{existe})
\begin{equation}\label{radi}
\ps^\infty\bigl(\varXi(\g\cone^{c}_\infty(\g))\bigr)\leq Ce^{-\varphi^\infty_\flat(\cartan(\g))}.
\end{equation}

A key extra information available in the case of $\PSL(3,\R)$ is that the form $\varphi^\infty_\flat$ is explicit and, up to scaling, doesn't depend on $\rho$. For this we need a small parenthesis on the \emph{critical hypersurface} $\calQ_\rho$ of $\rho$, depicted in Figure~\ref{calQ},
and characterized by
\[
\calQ_\rho=\bigl\{\varphi\in\a^*: \hC\varphi =1\bigr\},
\]
where the \emph{critical exponent} of a functional $\varphi\in\a^*$ is
\[
\hC\varphi:=\lim_{t\to\infty}\frac1t\log\#\bigl\{\g\in\piS:\varphi(\cartan(\g))\leq t\bigr\}\in(0,\infty].
\]
The critical hypersurface $\calQ_\rho\subset\a^*$ is a closed analytic curve that bounds a strictly convex set (\cite{exponentecritico,exponential}),
%(Sambarino~\cite{exponential} and Potrie--Sambarino~\cite{exponentecritico}),
and thus by~\cite{Quint-Div}, the linear form $\varphi^\infty_\flat$ is uniquely determined by
\begin{equation}\label{vi=norm}
\|\varphi^\infty_\flat\|^1=\inf\big\{\|\varphi\|^1:\varphi\in\calQ_\rho\big\}.
\end{equation}


\begin{figure}[h]\centering
\begin{tikzpicture}

\begin{scope}[scale=0.7]

\begin{scope}[rotate=-30,scale=0.5]

\begin{scope}[shift={(0,3)},scale=0.85]
\draw[name path global=Q, thick] (-13,3).. controls (0,-2.5).. (13,3);
\node at (13,3) [right] {$\calQ_\rho$};
\end{scope}

\draw[name path global=ii, dashed] (0,0) -- (0,6);
\fill[name intersections={of=Q and ii, by={J}}] (J) circle (1pt);

\begin{scope}[shift={(J)}]
\draw [blue] (-6,0)--(6,0);
\end{scope}

\draw[<->] (-1,6) to [bend left] (1,6);
\node at (0,7) {$\ii$};
\end{scope}

\fill (J) circle (2pt);

\draw [name path=slroot2] (0,0) -- (120:5cm);
\draw (0,0) -- (180:2cm);

\draw (0,0) -- (300:2cm);
\draw [name path=slroot1] (0,0) -- (360:5cm);
\node at (5,.5) [above] {$\a^*$};

\fill [name intersections={of=Q and slroot2, by={a}}] (a) circle (2pt);
\node [above right] at (a) {$\slroot_2$};

\fill [name intersections={of=Q and slroot1, by={b}}] (b) circle (2pt);
\node [below] at (b) {$\slroot_1$};

\draw [name path=borde, blue] (a) -- (b);

\fill [name intersections={of=borde and ii, by={c}}] (c) circle (2pt);
\node [right] at (c) {\footnotesize{$\sf H=\frac{\slroot_1+\slroot_2}2$}};

\node [left] (h') at (J) {};

\node [above left] (def) at (-3.5,0.5) {$\varphi^\infty_\flat=\hJ{\sf H}\sf H$};

\draw[->, orange] (def).. controls (2,2.5) and (-3.5,-.5).. (h');
\end{scope}
\end{tikzpicture}
\caption{The critical hypersurface of a strictly convex projective structure on $S$. Since $\sf H$ is a convex combination of $\{\slroot_1,\slroot_2\}$ one has $\|\sf H\|^1=1$ and thus $\|\varphi^\infty_\flat\|^1=\hC{\sf H}$.}\label{calQ}
\end{figure}


\pagebreak{}
Again by~\cite{exponentecritico} one has $\{\slroot_1,\slroot_2\}\subset\calQ_\rho$. Since both $\calQ_\rho$ and the norm $\|\,\|^1$ are invariant by the opposition involution $\ii$ (see again Figure~\ref{calQ}) we deduce that, if we let $\sf H=(\slroot_1+\slroot_2)/2$, then
\begin{equation}\label{>min}
\varphi^\infty_\flat=\hC{\sf H}\cdot\sf H\geq \hC{\sf H}\min\{\slroot_1,\slroot_2\}.
\end{equation}
In particular, using Eq.~\eqref{vi=norm}, we obtain that $\hC{\max\{\slroot_1,\slroot_2\}}=\hC{\sf H}$. Moreover, since the geodesic flow is Anosov (by Benoist~\cite{convexes1}) we can apply Bowen's characterization of entropy~\cite{axiomA} (and Remark~\ref{r.entropy=crit}) to obtain that the Hilbert entropy $\hJ{\sf H}=\hC{\sf H}$.


After this small parenthesis on the critical hypersurface, we come back to the lower bound on the Hausdorff dimension. Since $\varXi$ is a graph, $\varXi(\bord\piS)$ has the same intersection with the rectangle in Figure~\ref{figg} than with the larger square of size
\[
e^{-\min\{\slroot_1(\cartan(\alpha_i)),\slroot_2(\cartan(\alpha_i))\}};
\]
this square is now a ball (for the $L^\infty$ metric) of radius $e^{-\min\{\slroot_1(\cartan(\alpha_i)),\slroot_2(\cartan(\alpha_i))\}}$. Thus for all $i$, $\varXi(\alpha_i\cone^{c}_\infty(\alpha_i))$ is coarsely a ball of the latter radius and one has
\begin{align*}
\ps^\infty\bigl(B(\varXi(x),e^{-\min\{\slroot_1(\cartan(\alpha_i)),\slroot_2(\cartan(\alpha_i))\}}\bigr) & \leq \ps^\infty\bigl(\varXi(\alpha_i\cone^{c}_\infty(\alpha_i))\bigr)\leq Ce^{-\varphi^\infty_\flat(\cartan(\alpha_i))}\\
& \leq C\bigl(e^{-\min\{\slroot_1(\cartan(\alpha_i)),\slroot_2(\cartan(\alpha_i))\}}\bigr)^{\hC{\sf H}},
\end{align*}

\noindent where the last inequalities follow from Eqs.~\eqref{radi} and~\eqref{>min}. This gives a possibly bigger constant $C'$ such that, for all $r$,
\[
\ps^\infty\big(B(\varXi(x),r)\big)\leq C'r^{\hC{\sf H}}.
\]
Again, classical Hausdorff dimension arguments (c.f. Corollary~\ref{edgar} below) give that, for any measurable subset $E\subset\varXi(\bord\piS)$ with full $\ps^\infty$ mass, one has $\Hff(E)\geq \hC{\sf H}$.


Since $\PSL(3,\R)$ has rank smaller than 3 and $\rho$ is $\simple$-Anosov we can apply~\cite[Theorem~1.6]{BLLO}
%Burger--Landesberg--Lee--Oh 
to obtain that $\ps^\infty(\{\flat\text{-conical}\})=1$ and thus we have the desired lower bound
\[
\Hff\bigl(\{\flat\text{-conical}\}\bigr)\geq \hC{\sf H},
\]
which combined with the upper bound~\eqref{u} and the equality $\hC{\max\{\slroot_1,\slroot_2\}}=\hC{\sf H}$, gives the proof of Corollary~\ref{hilbert}.


In the general case~\cite[Theorem~1.6]{BLLO} is not applicable and we replace it with Theorem~\ref{Ahyper}.
\qed
\medskip


\subsection*{Structure of the paper}

The preliminaries of the paper are standard facts about linear algebraic groups, recalled in Section~\ref{Liegroups}, the work of Sambarino in~\cite{dichotomy} about linear cocycles over the boundary of a hyperbolic group (in Section~\ref{cocycles}), as well as basic facts about Anosov representations and their Patterson--Sullivan theory recalled from~\cite{BPS, GW, PSW1, dichotomy} in the first part of Section~\ref{Anosov}. In the rest of Section~\ref{Anosov} we prove Theorem~\ref{hyper} a more precise statement than Theorem~\ref{Ahyper}, discussing the Patterson--Sullivan measure of $(\sf W,\varphi)$-conical points. The heart of the proof is to construct and study a rank 2 flow whose recurrence set is related to $(\sf W,\varphi)$-conical points.

In Section~\ref{s.Hffbcon} we consider two locally conformal representations. We prove Theorem~\ref{Hffconical}, stating that for such a pair the Hausdorff dimension of the set of $\hol$-conical points belongs to
\[
\left[\hol\hC{\infty,\hol},\min\left\{\hC{\infty,\hol},\hol\hC{\infty,\hol}+1-\hol\right\}\right].
\]
The lower bound is obtained by analyzing properties of the linear form $\vi$ whose associated growth direction is $(\hol,1)$; its Patterson--Sullivan measure $\mu^{\vi}$ gives full mass to the set of $\hol$-conical points thanks to Theorem~\ref{hyper}. Using cone-types we can show that for a fine set of balls $\mu^{\vi}(B(x,r))\leq Cr^{-\hol\hC{\infty,\hol}}$. The upper bound uses results of~\cite{PSW1} to construct a fine covering of the set of $\hol$-conical points with balls of radius $e^{-\max\{\hol\slroot,\ovsloot\}}$. In Section~\ref{teoLCdiff} we prove Theorem~\ref{LCdiff}.

In Section~\ref{s.6} we prove that if the graph map between $\R$-hyperconvex representations has an oblique derivative, then the map is bi-Lipschitz (Proposition~\ref{c.indepPer}). This only relies on basic properties of hyperconvex representations, and is crucial for the proof of Theorem~\ref{tutti}, achieved in Section~\ref{s.Hffproof}, as it allows the identification of $\flat$-conical points and points of non-differentiability.



\subsection*{Acknowledgments}

We thank Katie~Mann, Gabriele~Viaggi, Anna~Wienhard and Maxime~Wolff for insightful conversations and Andrés~Navas for pointing us to useful literature.



\section{Linear algebraic groups}\label{Liegroups}

Throughout the text $\sfG$ will denote a real-algebraic semi-simple Lie group of non-compact type and $\ge$ its Lie algebra.


\subsection{Linear algebraic groups}

Fix a Cartan involution $o:\ge\to\ge$ with associated Cartan decomposition $\ge=\k\oplus\p$. Let $\a\subset\p$ be a maximal abelian subspace and let $\roots\subset\a^*$ be the set of restricted roots of $\a$ in $\ge$. For $\sroot\in\roots$, we denote by
\[
\ge_\sroot=\{u\in\ge:[a,u]=\sroot(a)u\;\forall a\in\a\}
\]
its associated root space. The (restricted) root space decomposition is $\ge=\ge_0\oplus\bigoplus_{\sroot\in\roots}\ge_\sroot$, where $\ge_0$ is the centralizer of $\a$. Fix a Weyl chamber $\a^+$ of $\a$ and let $\roots^+$ and $\simple$ be, respectively, the associated sets of positive and simple roots. Let $\Weyl$ be the Weyl group of $\roots$ and $\ii:\a\to\a$ be the opposition involution: if $u:\a\to\a$ is the unique element in $\Weyl$ with $u(\a^+)=-\a^+$ then $\ii=-u$.

We denote by $(\cdot,\cdot)$ both the Killing form of $\ge$, its restriction to $\a$, and its associated dual form on $\a^*$, the dual of $\a$. For $\chi,\psi\in\a^*$ let
\[
\<\chi,\psi\>=2\frac{(\chi,\psi)}{(\psi,\psi)}.
\]
The \emph{restricted weight lattice} is defined by
\[
\poids=\left\{\varphi\in\a^*:\<\varphi,\sroot\>\in\Z\;\forall\sroot\in\roots\right\}.
\]
It is spanned by the \emph{fundamental weights} $\{\peso_\sroot:\sroot\in\simple\}$, defined by
\begin{equation}\label{pesoFund}
\<\peso_\sroot,\bb\>=d_\sroot\delta_{\sroot\bb}
\end{equation}
for every $\sroot,\bb\in\simple$, where $d_\sroot=1$ if $2\sroot\notin\roots^+$ and $d_\sroot=2$ otherwise.



A subset $\t\subset\simple$ determines a pair of opposite parabolic subgroups $\sf P_\t$ and $\check{\sf P}_\t$ whose Lie algebras are
\begin{align*}
\p_\t &=\bigoplus_{\sroot\in\roots^+\cup\{0\}}\ge_{\sroot} \oplus\bigoplus_{\sroot\in\<\simple-\t\>}\ge_{-\sroot},\\
\check{\p_\t}
&=\bigoplus_{\sroot\in\roots^+\cup\{0\}}\ge_{-\sroot} \oplus\bigoplus_{\sroot\in\<\simple-\t\>}\ge_{\sroot}.
\end{align*}
The group $\check{\sf P}_\t$ is conjugated to the parabolic group $\sf P_{\ii\t}$. We denote the \emph{flag space} associated to $\t$ by $\calF_\t=\sfG/\sf P_\t$. The $\sfG$ orbit of the pair $([\sf P_{\t}],[\check{\sf P}_{\t}])$ is the unique open orbit for the action of $\sfG$ in the product $\calF_\t\times\calF_{\ii\t}$ and is denoted by $\posgen_\t$.



\subsection{Cartan and Jordan projection}


Denote by $\sf K=\exp\k$ and $\sf A=\exp\a$. The \emph{Cartan decomposition} asserts the existence of a continuous map $\cartan:\sfG \to \a^+$, called the \emph{Cartan projection}, such that every $g\in\sfG$ can be written as $g=ke^{a(g)}l$ for some $k,l\in\sf K$.


We will need the following uniform continuity of the Cartan projection:


\begin{prop}[{\cite[Proposition~5.1]{Benoist-HomRed}}]\label{p.Ben}
For any compact $ L\subset \sfG$ there exists a compact set $H\subset \a$ such that, for every $g\in \sfG$, one has
\[
\cartan(LgL)\subset \cartan(g)+H.
\]
\end{prop}



By the Jordan's decomposition, every element $g\in\sfG$ can be uniquely written as a commuting product $g=g_eg_{ss}g_u$ where $g_e$ is conjugate to an element in $\sf K$, $g_{ss}$ is conjugate to an element in $\exp(\a^+)$ and $g_u$ is unipotent. The \emph{Jordan projection} $\lambda=\lambda_{\sfG}:\sfG\to\a^+$ is the unique map such that $g_{ss}$ is conjugated to $\exp\big(\lambda(g)\big)$.


\begin{defi}
Let $\Gamma\subset\sfG$ be a discrete subgroup, then its \emph{limit cone} $\cal L_\Gamma$ is the smallest closed cone of the closed Weyl chamber $\a^+$ that contains $\{\lambda(g):g\in\Gamma\}$.
\end{defi}


We will need the following results by Benoist.

\begin{theo}[{Benoist~\cite{limite,benoist2}}]\label{densidad}
Let $\Gamma\subset\sfG$ be a Zariski-dense sub-semigroup, then its limit cone $\cal L_\Gamma$ has non-empty interior. Moreover, the group generated by the Jordan projections $\{\lambda(g):g\in\Gamma\}$ is dense in $\a$.
\end{theo}



\subsection{Representations of \texorpdfstring{$\sfG$}{sf G}}\label{representaciones}
The standard references for the following are Fulton--Harris~\cite{FultonHarris}, Humphreys \cite{james} and Tits~\cite{Tits}.



Let $\Fund:\sfG\to\PGL(V)$ be a finite dimensional rational\footnote{Namely a rational map between algebraic varieties.} irreducible representation and denote by $\phi_\Fund:\ge\to\fksl(V)$ the Lie algebra homomorphism associated to $\Fund$. The \emph{weight space} associated to $\chi\in\a^*$ is the vector space
\[
V_\chi=\{v\in V:\phi_\Fund(a) v=\chi(a) v\ \forall a\in\a \}.
\]

\noindent We say that $\chi\in\a^*$ is a \emph{restricted weight} of $\Fund$ if $V_\chi\neq0$. Theorem~7.2 from~\cite{Tits} states that the set of weights has a unique maximal element with respect to the partial order $\chi\succ\psi$ if $\chi-\psi$ is a $\N$-linear combination of positive roots. This is called \emph{the highest weight} of $\Fund$ and denoted by $\chi_\Fund$. By definition, for every $g\in\sfG$ one has
\begin{equation}\label{eq:spectralrep}
\lambda_1(\Fund(g))=\chi_\Fund(\lambda(g)),
\end{equation}
where $\lambda_1$ is the logarithm of the spectral radius of $\Fund(g)$.

We denote by $\poids(\phi)$ the set of restricted weights of the representation $\phi_\Fund$
\[
\poids(\phi)=\bigl\{\chi\in\a^*:V_\chi\neq\{0\}\bigr\},
\]
these are all bounded above by $\chi_\Fund$ (see for example~\cite[Section~13.4, Lemma~B]{james}),
namely every weight $\chi\in\poids(\phi)$ has the form
\[
\chi_\Fund-\sum_{\sroot\in\simple} n_\sroot\sroot\quad \text{for}\quad n_\sroot\in\N.
\]
The \emph{level} of a weight $\chi$ is the integer $\sum_\sroot n_\sroot$, the highest weight is thus the only weight of level zero.
Additionally, if $\chi\in\poids(\phi_\Fund)$ and $\sroot\in\roots^+$ then the elements of $\poids(\phi_\Fund)$ of the form $\chi+j\sroot,\, j\in\Z$ form an unbroken string
\[
\chi+j\sroot,\: j\in\lb-r,q\rb
\]
and $r-q=\<\chi,\sroot\>$. One can then recover algorithmically the set $\poids(\phi_\Fund)$ level by level starting from $\chi_\Fund$, as follows:
\begin{itemize}
\item Assume the set of weights of level at most $ k$ is known and consider a weight $\chi$ of level $k$.
\item For each $\sroot\in\simple$ compute $\<\chi,\sroot\>$, this gives the length $r-q$ of the $\sroot$-string through $\chi$. The weights of the form $\chi +j\sroot$, for positive $j$, have level smaller than $k$ and are thus known, thus we can decide whether $\chi-\sroot$ is a weight or not, determining the set of weights of level $k+1$.
\end{itemize}

The following lemma follows at once from the algorithmic description above. Let $\ge=\bigoplus_i\ge_i$ be the decomposition in simple factors of a semi-simple real Lie algebra of non-compact type. Recall that if $\a_i \subset\ge_i$ is a Cartan subspace, then $\a=\bigoplus_i\a_i$ is a Cartan subspace of $\ge$. Any $\varphi\in(\a_i)^*$ extends to a functional on $\a$, still denoted~$\varphi$, by vanishing on the remaining factors. The restricted root system of $\ge$ is then $\simple_\ge=\bigcup \simple_{\ge_i}$. The \emph{associated simple factor} to $\sroot\in\simple_\ge$ is $\ge_i$ such that $\sroot\in\simple_i$.



\begin{lemm}\label{simple}
Let $\ge$ be a semi-simple real Lie algebra of non-compact type and $\phi$ be an irreducible representation of $\ge$ whose highest restricted weight is a multiple of a fundamental weight, $\chi_\phi=k\peso_\sroot$ for some $\sroot\in\simple$. Then $\phi$ factors as a representation of the simple factor associated to $\sroot$.
\end{lemm}

\begin{proof}
Proceeding by induction on the levels of $\phi$, one readily sees that for every $\slroot\in\simple_j$ for $j\neq i$ and all $\chi\in\poids(\phi)$ one has $\<\chi,\slroot\>=0$. Thus the associated root space $(\ge_j)_{-\slroot}$ acts trivially on every weight space of $\phi$ and so the whole factor $\ge_j$ acts trivially.
\end{proof}

The following set of simple roots plays a special role in representation theory.


\begin{defi}\label{trep}
Let $\Fund:\sfG\to\PGL(V)$ be a representation. We denote by $\t_\Fund$ the set of simple roots $\sroot\in\simple$ such that $\chi_\Fund-\sroot$ is still a weight of $\Fund$. Equivalently
\begin{equation}\label{peso<}
\t_\Fund=\bigl\{\sroot\in\simple:\<\chi_\Fund,\sroot\>\neq0
\bigr\}.
\end{equation}
\end{defi}


The following lemma will be needed in the proof of Theorem~\ref{t.Zcl}.

\begin{lemm}\label{not}
Let $\ge$ be semi-simple of non-compact type and $\phi:\ge\to\fkgl(V)$ an irreducible representation. For $\sroot\in\t_\phi$ and $v\in V_{\chi_\phi}-\{0\}$, the map $n\mapsto \phi(n)v$ is injective when defined on $\ge_{-\sroot}$.
\end{lemm}


\begin{proof}
By definition of $\chi_\phi$ every $n\in\ge_{\sroot}$ acts trivially on $V_{\chi_\phi}$. For $y\in\ge_{-\sroot}-\{0\}$, there exists $x\in\ge_\sroot$ such that $\{x,y,h_\sroot\}$ spans a Lie algebra isomorphic to $\sl_2(\R)$, where $h_\sroot$ is defined by $\varphi(h_\sroot)=\<\varphi,\sroot\>$ for all $\varphi\in\a^*$. If $\phi(y)v=0$ then, since $\phi(x)V_{\chi_\phi}=0$ one concludes $\phi(h_\sroot)v=0$ and since $V_{\chi_\phi}$ is a weight-space one has $\phi(h_\sroot)V_{\chi_\phi}=0$. This in turn implies that
\[
\<\chi_\phi,\sroot\>=\chi_\phi(h_\sroot)=0,
\]
contradicting that $\sroot\in\t_\phi$.
\end{proof}


We denote by $\|\,\|_\Fund$ an Euclidean norm on $V$ invariant under $\Fund \sf K$ and such that $\Fund\sf A $ is self-adjoint, see for example Benoist--Quint's book~\cite[Lemma~6.33]{BQLibro}. By definition of $\chi_\Fund$ and $\|\,\|_\Fund$, and Eq.~\eqref{eq:spectralrep} one has, for every $g\in \sfG, $ that
\begin{equation}\label{eq:normayrep}
\log\|\Fund g\|_\Fund=\chi_\Fund(\cartan(g)).
\end{equation}
Here, with a slight abuse of notation, we denote by $\|\,\|_\Fund$ also the induced operator norm, which doesn't depend on the scale of $\|\,\|_\Fund$.


Denote by $W_{\chi_\Fund}$ the $\Fund\sf A $-invariant complement of $V_{\chi_\Fund}$. The stabilizer in $\sfG $ of $W_{\chi_\Fund}$ is $\wk{\sf P}_{\t_\Fund}$, and thus one has a map of flag spaces
\begin{equation}\label{maps}
(\zeta_\Fund,\zeta^*_\Fund):\calF_{\t_\Fund}^{(2)}(\sfG)\to \grass_{\dim V_{\chi_\Fund}}^{(2)}(V).
\end{equation}
This is a proper embedding which is an homeomorphism onto its image. Here, as above, $\grass_{\dim V_{\chi_\Fund}}^{(2)}(V)$ is the open $\PGL (V)$-orbit in the product of the Grassmannian of $(\dim V_{\chi_\Fund})$-dimensional subspaces and the Grassmannian of $(\dim V-\dim V_{\chi_\Fund})$-dimensional subspaces. One has the following proposition (see also~\cite[Chapter~XI]{LAG}).

\begin{prop}[{\cite{Tits}}]\label{prop:titss}
For each $\sroot\in\simple$ there exists a finite dimensional rational irreducible representation $\Fund_\sroot:\sfG\to\PSL(V_\sroot)$, such that $\chi_{\Fund_\sroot}$ is an integer multiple $l_\sroot\peso_\sroot$ of the fundamental weight and $\dim V_{\chi_{\Fund_\sroot}}=1$.
\end{prop}


We will fix from now on such a set of representations and call them, for each $\sroot\in\simple$, the \emph{Tits representation associated to $\sroot$}.

\subsection{The center of the Levi group \texorpdfstring{$\sf P_{\t}\cap\check{\sf P}_{\t}$}{sf P\_ t cap check sf P\_ t}}\label{s.Levi}


We now consider the vector subspace
\[
\a_\t=\bigcap_{\sroot\in\simple-\t}\ker\sroot.
\]
Denoting by $\Weyl_\t=\{w\in \Weyl:w(v)=v\quad \forall v\in\a_\t\}$ the subgroup of the Weyl group generated by reflections associated to roots in $\simple-\t$, there is a unique projection $\pi_\t:\a\to\a_\t$ invariant under $\Weyl_\t$.

The dual $(\a_\t)^*$ is canonically identified with the subspace of $\a^*$ of $\pi_\t$-invariant linear forms. Such space is spanned by the fundamental weights of roots in $\t$,
\[
(\a_\t)^*=\left\{\varphi\in\a^*:\varphi\circ\pi_\t=\varphi
\right\}=\left\<\peso_\sroot|\a_\t:\sroot\in\t\right\>.
\]
We will denote, respectively, by
\begin{alignat*}{2}\cartan_\t& =\pi_\t\circ \cartan:\sfG\to \a_\t\ \\ \lambda_\t & =\pi_\t\circ\lambda:\sfG\to \a_\t,
\end{alignat*} the compositions of the Cartan and Jordan projections with $\pi_\t$.

\subsection{The Busemann--Iwasawa cocycle}

The \emph{Iwasawa decomposition} of $\sfG$ states that every $g\in\sfG$ can be written uniquely as a product $lzu$ with $l\in\sf K$, $z\in\sf A$ and $u\in\sf U_\simple$, where $\sf U_\simple$ is the unipotent radical of $\sf P_{\simple}$.

The \emph{Busemann--Iwasawa cocycle} of $\sfG$ is the map $\bus:\sfG\times\calF\to\a$ such that, for all $g\in\sfG$ and $k[\sf P_\simple]\in\calF$,
\[
\bus(g,k[\sf P_\simple])=\log(z)
\]
where $\log:\sf A\to \a$ denotes the inverse of the exponential map, and $gk=lzu$ is the Iwasawa decomposition of $gk$. Lemmes~6.1 and 6.2 from~\cite{Quint-PS} proved that the function $\bus_\t=\pi_\t\circ\bus$ factors as a cocycle $\bus_\t:\sfG\times \calF_\t\to\a_\t$.


The Busemann--Iwasawa cocycle can also be read from the representations of $\sfG$. Indeed, \cite[Lemme~6.4]{Quint-PS} shows that for every $g\in\sfG$ and $x\in\calF_\t$ one has
\begin{equation}\label{busnorma}
l_\sroot\peso_\sroot(\bus(g,x))=\log\frac{\|\Fund_\sroot(g)v\|_\Fund}{\|v\|_\Fund},
\end{equation}
where $v\in\zeta_{\Fund_\sroot}(x)\in\P(\sf V_\sroot)$ is non-zero, and $l_\sroot$ is as in Proposition~\ref{prop:titss}.



\subsection{Gromov product and Cartan attractors}\label{GryBus}

Let $\K$ be either $\C$ or $\R$. For a decomposition $\K^d=\ell\oplus V$ into a line $\ell$ and a~hyperplane $V$ together with an inner (Hermitian) product $o$ on $\K^d$, one defines the \emph{Gromov product} by
\[
\Gr(\ell,V)=\Gr^o(\ell,V):=\log\frac{|\varphi(v)|}{\|\varphi\|\|v\|}=\log\sin\angle_o(\ell,V),
\]
for any non-vanishing $v\in\ell$ and $\varphi\in (\K^d)^*$ with $\ker\varphi=V$.



This induces, for any semisimple Lie group $\sfG$ and subset $\t<\Delta$, a \emph{Gromov product} $\Gr_\t:\posgen_\t\to\a_\t$ defined, for every $(x,y)\in\posgen_\t$ and $\sroot\in\t$, by
\[
l_\sroot\peso_\sroot(\Gr_\t(x,y))= \Gr^{\Fund_\sroot}(\zeta_{\Fund_\sroot} x,\zeta_{\Fund_\sroot}^*y)=\log\sin\angle_o\left(\zeta_{\Fund_\sroot} y,\zeta_{\Fund_\sroot}^*x\right),
\]
where $\zeta^*_{\Fund_\sroot}$ and $\zeta_{\Fund_\sroot}$ are the equivariant maps from Eq.~\eqref{maps}, and the Hermitian product~$o$ is induced by an Euclidean norm $\|\,\|_{\Fund_\sroot}$ invariant under $\Fund_\sroot\sf K$.

From~\cite[Lemma~4.12]{orbitalcounting} one has, for all $g\in\sfG$ and $(x,y)\in\posgen_\t$,
\begin{equation}\label{forGr}
\Gr_\t(gx,gy)-\Gr_\t(x,y)=-\big(\bus_\t(g,x)+\ii\bus_{\ii\t}(g,y)\big).
\end{equation}



If $g=k\exp(\cartan(g))l$ is a Cartan decomposition of $g\in\sfG$ we define its $\t$-\emph{Cartan attractor} (resp. \emph{repeller}) by
\[
U_\t(g)=k[\sf P_\t]\in\calF_\t\quad\text{and}\quad U_{\ii\t}(g^{-1})=l^{-1}[\check{\sf P}_\t]\in\calF_{\ii\t}.
\]
The \emph{Cartan basin of $g$} is defined, for $\alpha>0$, by
\begin{equation}\label{e.CartBasin}
B_{\t,\alpha}(g)= \left\{x\in\calF_\t:\peso_\sroot \Gr_\t\left(x,U_{\ii\t}(g^{-1})\right)>-\alpha,\; \forall \sroot\in\theta\right\}.
\end{equation}

\begin{rema}\label{fty}
Observe that a statement of the form $\peso_\sroot\Gr_\t(x,y) \geq-\kappa$ for all $\sroot\in\t$ is a quantitative version (depending on the choice of $\sf K$) of the transversality between $x$ and $y;$ in particular it implies that $x$ and $y$ are transverse.
\end{rema}

Neither the Cartan attractor nor its basin are uniquely defined unless for all $\sroot\in\t$ one has $\sroot\big(\cartan(g)\big)>0$, regardless one has the following:

\begin{rema}
Given $\alpha>0$ there exists a constant $K_\alpha$ such that if $y\in\calF_{\t}$ belongs to $B_{\t,\alpha}(g)$ then one has
\begin{equation}\label{comparision-shadow}
\big\|\cartan_\t(g)-\bus_\t(g,y)\big\|\leq K_\alpha.
\end{equation}
Indeed, using Tits's representations of $\sfG$ and Eqs.~\eqref{eq:normayrep} and~\eqref{busnorma} this boils down to the elementary fact that if $A\in\GL_d(\R)$ verifies\footnote{Recall from Eq.~\eqref{e.slroot} that we denote by $\slroot_i$ the simple roots of $\GL_d(\RR)$.} $\slroot_1(\cartan(A))>0$ then for every $v\in\R^d$ one has
\[
\log\frac{\|Av\|}{\|v\|}\geq \log\|A\|+\log\sin\angle\big(\R\cdot v,U_{d-1}(A^{-1})\big)
\]
(see for example~\cite[Lemma~A.3]{BPS}).
\end{rema}



\section{H\"older cocycles on \texorpdfstring{$\bord\G$}{bord G}}\label{cocycles}

Let $\G$ be a finitely generated group, and fix a finite generating set $S$. A group $\G$ is \emph{Gromov hyperbolic} if its Cayley graph $\sf{Cay}(\G, S)$ is a Gromov hyperbolic geodesic metric space. In this case we denote by $\bord\G$ its Gromov boundary, namely the equivalence classes of (quasi)-geodesic rays. It is well known that, up to H\"older homeomorphism, $\bord\G$ doesn't depend on the choice of the generating set $S$. We will denote by $\bord^2\G$ the set of distinct pairs in $\bord\G$:
\[
\bord^2\G:=\{(x,y)\in\bord\G\times\bord\G|\, x\neq y\}.
\]


For a finitely generated, non-elementary, word-hyperbolic group $\G$ we denote by $\sf g=(\sf g_t:\sf U\G\to\sf U\G)_{t\in\R}$ the \emph{Gromov--Mineyev geodesic} flow of $\G$ (see~\cite{gromov, mineyev}).
%Gromov~\cite{gromov} and Mineyev~\cite{mineyev}
Throughout this section we will have the same assumptions as in~\cite[Section~3]{dichotomy}, namely that $\sf g$ is metric-Anosov and that the lamination induced on the quotient by $\widetilde{\cal W}^{cu}=\{(x,\cdot,\cdot)\in\widetilde{\sf U\G}\}$ is the central-unstable lamination of $\sf g$.

Since we will mostly recall needed results
% S.
from~\cite[Section~3]{dichotomy} we do not overcharge the paper with the definitions of metric-Anosov and central-unstable lamination: by~\cite{pressure}, %Bridgeman--Canary--Labourie--Sambarino 
word-hyperbolic groups admitting an Anosov representation verify the required assumptions.


\begin{defi}
Let $V$ be a finite dimensional real vector space. A \emph{H\"older cocycle} is a function $c:\G\times\bord\G\to V$ such that:
\begin{itemize}
\item for all $\g,h\in \G$ one has $c(\g h,x)=c(h,x)+c(\g,h(x))$,
\item there exists $\alpha\in(0,1]$ such that for every $\g\in\G$ the map $c(\g,\cdot)$ is $\alpha$-H\"older continuous.
\end{itemize}
\end{defi}



Recall that every \emph{hyperbolic element}\footnote{Namely an infinite order element.} $\g\in\G$ has two fixed points on $\bord\G$, the attracting $\g_+$ and the repelling $\g_-$. If $x\in\bord\G-\{\g_-\}$ then $\g^nx\to\g_+$ as $n\to\infty$. The \emph{period} of a H\"older cocycle for a hyperbolic $\g\in\G$ is $\ell_c(\g):=c(\g,\g^+)$. A cocycle $c^*:\G\times\bord\G\to\R$ is \emph{dual to $c$} if for every hyperbolic $\g\in\G$ one has
\[
\ell_{c^*}(\g)=\ell_c(\g^{-1}).
\]



\subsection{Real-valued cocycles}\label{Sreal}



Assume now $V=\R$ and consider a cocycle $\kappa$ with non-negative (and not all vanishing) periods. We denote by $[\G]$ the set of conjugacy classes in $\G$ and, for an element $\g\in\G$ we denote by $[\g]$ its conjugacy class. For $t>0$ we let
\[
\sf R_t(\kappa)=\bigl\{[\g]\in[\G]\text{ hyperbolic}:\ell_\kappa(\g)\leq t\bigr\}
\]
and define the \emph{entropy} of $\kappa$ by
\[
\hJ\kappa=\limsup_{t\to\infty}\frac1t\log\#\sf R_t(\kappa)\in(0,\infty].
\]
For such a cocycle consider the action of $\G$ on $\bord^2\G\times\R$ via $\kappa$:
\begin{equation}\label{gaction}
\g\cdot(x,y,t)=\left(\g x,\g y,t-\kappa\left(\g,y\right)\right).
\end{equation}


The following is a straightforward consequence of~\cite[Theorem~3.2.2]{dichotomy}.


\begin{prop}\label{-infty}
Let $\kappa$ be a H\"older cocycle with non-negative periods and finite entropy. Then, the above action of $\G$ on $\bord^2\G\times\R$ is properly-discontinuous and co-compact. If moreover $c$ is another H\"older cocycle with non-negative periods and finite entropy then there exists a $\G$-equivariant bi-H\"older-continuous homeomorphism $E:\bord^2\G\times\R\to\bord^2\G\times\R$ which is an orbit equivalence between the $\R$-translation actions.
\end{prop}


We recall the notion of \emph{dynamical intersection}, a~concept from~\cite{pressure} %Bridgeman--Canary--Labourie--Sambarino 
for H\"older functions over a metric-Anosov flow, that can be pulled back to this setting via the existence of the \emph{Ledrappier potential} of $\kappa$ from~\cite[Section~3.1]{dichotomy}.


The \emph{dynamical intersection} of two real valued cocycles $\kappa, c$ is
\begin{equation}\label{defiII}
\II(\kappa,c)=\lim_{t\to\infty}\frac1{\sf R_t(\kappa)}\sum_{\g\in\sf R_t(\kappa)}\frac{\ell_c(\g)}{\ell_\kappa(\g)}.
\end{equation}
We record in the following proposition various needed facts about $\II$:

\begin{prop}[{\cite[Section~3]{pressure}}]\label{ineq}
The dynamical intersection defined above is well defined, linear in the second variable and for all positive $s$ satisfies $\II(s\kappa,c)=\II(\kappa,c)/s$. If also $c$ has non-negative periods and finite entropy then $\II(\kappa,c)\geq \hJ\kappa/\hJ c$. Moreover, if $\II(\kappa,c)= \hJ\kappa/\hJ c$ then for every $\g\in\G$ one has $\hJ\kappa\ell_\kappa(\g)=\hJ c\ell_c(\g)$.
\end{prop}



We will also need the following definitions.

\begin{defi}\leavevmode
\begin{itemize}
\item A \emph{Patterson--Sullivan measure for $\kappa$ of exponent} $\delta\in\R_+$ is a probability measure $\ps$ on $\bord\G$ such that for every $\g\in\G$ one has
\begin{equation}\label{defP}
\frac{d \g_*\ps}{d\ps}(\cdot)=e^{-\delta\cdot \kappa\left(\g^{-1},\,\cdot\,\right)}.
\end{equation}
\item Let $\kappa^*$ be a cocycle dual to $\kappa$, then a \emph{Gromov product} for the ordered pair $(\kappa^*, \kappa)$ is a function $[\cdot,\cdot]:\bord^2\G\to\R$ such that for all $\g\in\G$ and $(x,y)\in\bord^2\G$ one has
\[
[\g x,\g y]-[x,y]=-\bigl(\kappa^*(\g,x)+\kappa(\g,y)\bigr).
\]
\end{itemize}
\end{defi}



\subsection{The critical hypersurface and intersection}\label{QQ}

Let now $c:\G\times\bord\G\to V$ be a H\"older cocycle. Its \emph{limit cone} is denoted by
\[
\calL_c=\overline{\bigcup_{\g\in\G}\R_+\cdot\ell_c(\g)}
\]
and its \emph{dual cone} by $(\calL_c)^*=\{\psi\in V^*:\psi|_{\calL_c}\geq0\}$. Observe that for every $\varphi\in\inte(\calL_c)^*$, $\varphi\circ c$ is a real-valued cocycle, so the concepts from Section~\ref{Sreal} apply. We denote by
\begin{alignat}{2}\label{qh1}
\calQ_c & =\left\{\varphi\in\inte\cdc{c}:\hJ{\varphi\circ c}=1\right\},\\
\cal D_c & = \left\{\varphi\in\inte\cdc{c}:\hJ{\varphi\circ c}\in(0,1)\right\},
\end{alignat}
respectively the \emph{critical hypersurface} and the \emph{convergence domain} of $c$.

For $\varphi\in\inte(\calL_c)^*$ we consider the linear map $\II_{\varphi}=\II_{\varphi}^c:V^*\to\R$ defined by
\[
\II_{\varphi}^c(\psi):= \II(\varphi\circ c,\psi\circ c),
\]
as in Eq.~\eqref{defiII}. The natural identification between the set of hyperplanes in $V^*$ and $\P(V)$ is used in the next proposition.

\begin{coro}[{\cite[Corollary~3.4.3]{dichotomy}}]
\label{strictly}
Assume $\calL_c$ has non-empty interior and that there exists $\psi\in(\calL_c)^*$ such that $\hJ{\psi}<\infty$. Then $\cal D_c$ is a strictly convex set with boundary $\calQ_c$. The latter is an analytic co-dimension one sub-manifold of~$V^*$. The map $\sf u^c:\calQ_c\to\P(V)$ defined by
\[
\varphi\mapsto\sf u^c_\varphi:=\sf T_\varphi\calQ_c=\ker\II_\varphi
\]
is an analytic diffeomorphism between $\calQ_c$ and $\inte(\P(\calL_c))$.
\end{coro}



\subsection{Ergodicity of directional flows}\label{directional}


It follows from Proposition~\ref{-infty} that if there exists $\psi\in(\calL_c)^*$ with $\hJ\psi<\infty$ then the $\G$-action $\bord^2\G\times V$
\[
\g(x,y,v)=\big(\g x,\g y,v-c(\g,y)\big)
\]
is properly discontinuous.

\begin{defi}
A H\"older cocycle $c$ is \emph{non-arithmetic} if the periods of $c$ generate a dense subgroup in $V$.
\end{defi}

We fix $\varphi\in\calQ_c$ and denote by $u_\varphi\in\sf u_\varphi$ the unique vector in $\cal L_c\cap\sf u_\varphi$ with $\varphi(u_\varphi)=1$. We define then the \emph{directional flow}
\begin{align*}
\df^\varphi &=\left(\df^\varphi_t:\G\/\left(\bord^2\G\times V\right)\to\G\/\left(\bord^2\G\times V\right)\right)_{t\in\R}\\
\intertext{by}
t\cdot(x,y,v) &=(x,y,v-tu_\varphi).
\end{align*}



\begin{enonce}{Assumption}\label{Pat-SulEx}
We assume there exists:
\begin{itemize}
\item a dual cocycle $(\varphi\circ c)^*$,
\item a Gromov product $[\,,\,]_\varphi$ for such a pair,
\item Patterson--Sullivan measures, $\ps^\varphi$ and ${\ov\ps}^\varphi$, respectively for each of the cocycles; (the exponent of both measures is then necessarily $\hJ\varphi=1$~\cite[Proposition~3.3.2]{dichotomy}).
\end{itemize}
\end{enonce}



Consider then the $\varphi$-\emph{Bowen--Margulis} measure $\BM^\varphi$ on $\G\/(\bord^2\G\times V)$ defined as the measure induced on the quotient by the measure
\begin{equation}\label{forBM}
e^{-[\cdot,\cdot]_\varphi}{\ov\ps}^\varphi\otimes\ps^\varphi\otimes \Fundeb_{V},
\end{equation}
for a $V$-invariant Lebesgue measure on $V$. We denote by $\cal K(\df^\varphi)$ the \emph{recurrence set} of the directional flow $\df^\varphi$:
\[
\cal K(\df^\varphi):=\left\{p\in\G\/\left(\bord^2\G\times V\right)\,\middle|\, \exists B \text{ open bounded},\; t_n\to \infty \;\text{ with }\; \df^\varphi_{t_n}(p)\in B \right\}.
\]

\begin{coro}[{\cite[Corollary~3.6.1]{dichotomy}}]
\label{d=2}
Assume that $c$ is non-arithmetic, and that there exists $\varphi\in\calQ_c$ satisfying Assumptions~\ref{Pat-SulEx}. If $\dim V\leq2$ then the directional flow $\df^\varphi$ is $\BM^\varphi$-ergodic, and $\cal K(\df^\varphi)$ has total mass. If $\dim V\geq4$ then $\cal K(\df^\varphi)$ has measure zero.
\end{coro}



\section{Subspace conicality for Anosov representations: Theorem~\ref{Ahyper}}\label{Anosov}

\subsection{Gromov hyperbolic groups and cone types}\label{cont}
Let $\G=\langle S\rangle$ be a finitely generated non-elementary Gromov hyperbolic group, and recall from Section~\ref{cocycles} that we denote by $\bord^2\G$ the set of distinct pairs in its Gromov boundary~$\bord\G$.



\begin{defi}\label{conicalsequence}
A divergent sequence $\{\g_n\}_{n\in\N}\subset\G$ converges to a point $x\in\bord\G$ \emph{conically} if for every $y\in\bord\G-\{x\}$ the sequence $(\g_n^{-1}y,\g_n^{-1}x)$ remains on a~compact set of $\bord^2\G$.
\end{defi}

\begin{rema}\label{r.conicalgeodesic}
It is easy to verify that a sequence $\{\g_n\}_{n\in\N}$ converges conically to $x\in\bord\G$ if and only if it lies in an uniform neighborhood of any geodesic ray $(\alpha_n)_0^\infty$ converging to $x$, namely there exists $K>0$ and a subsequence $\{\alpha_{n_k}\}$ such that for all $k$ one has $d_\G(\alpha_{n_k},\g_{k})<K$.
\end{rema}



Given $\g\in\G$ we denote by $\cone(\g)$ the \emph{cone type} of $\g\in\G$, namely
\[
\cone(\g):=\left\{h\in\G\,\middle|\, d(e,\g h)=d(e,\g)+d(e,h)\right\}.
\]
Cannon~\cite{CannonCones} showed the set of cone types of a Gromov hyperbolic group is finite, see for example Bridson--Haefliger's book~\cite[p.~455]{BH}. We denote by $\cone_\infty(\g)\subset{\bord\G}$ the set of points $x$ that can be represented by a geodesic ray contained in $\cone(\g)$.


We will also need a coarse version of these. Recall that a sequence $(\alpha_j)_0^\infty$ is a $(c,C)$-quasigeodesic if for every pair $j,l$ it holds
\[
\frac1{c}|j-l|-C\leq d_\G(\alpha_j,\alpha_l)\leq c|j-l|+C.
\]
The \emph{coarse cone type at infinity} of an element $\g$ is the set of endpoints at infinity of quasi-geodesic rays based at $\g^{-1}$ passing through the identity:
\[
\cone^{c}_\infty(\g) = \left\{\left[(\alpha_j)_0^\infty\right]\in\bord\G\,\middle|\, (\alpha_j)_0^\infty \quad\text{is a $(c,c)$-quasi-geodesic,}\; \alpha_0=\g^{-1},\; e\in\{\alpha_j\}\right\}.
\]
\begin{figure}[hbt]
\centering
\begin{tikzpicture}[scale = 0.6]
\draw [thick] circle [radius = 3];\draw [fill] circle [radius = 0.03];
\draw circle [radius = 0.78];
\node [left] at (-0.65,0) {$B_{c}(e)$};
\draw [thick] (0.7,-2.2) -- (0.5,-1.7) -- (0.004,-1.3) -- (0.1,-0.7) -- (-0.2,-0.3) -- (0,0);
\node [left] at (0.7,-2.2) {$\g^{-1}$}; \draw [fill] (0.7,-2.2) circle [radius = 0.03];
\draw [thick] (1,0.2) -- (0.1,-0.7);
\draw [thick] (-0.2,-0.3) -- (-0.3,1);
\draw [thick] (0,0) -- (-0.3,1) -- (0.2,2) -- (0,3);
\draw [thick] (0.2,2) -- (0.4, 2.7) -- (0.7,2.91);
\draw [thick] (1,0.2) -- (0.89,1.37);
\draw [thick] (1.8,0.8) -- (2.7,1.3);
\draw [thick] (0,0) -- (1,0.2) -- (1.8,0.8) -- (2.846,0.948);

\draw [fill] (0.7,2.91) circle [radius = 0.03];
\draw [fill] (2.7,1.3) circle [radius = 0.03];
\draw [fill] (2.12,2.12) circle [radius = 0.03];
\draw [fill] (0,3) circle [radius = 0.03];
\draw [fill] (2.846,0.948) circle [radius = 0.03];
\draw [thick] (0,0) -- (0.89,1.37) -- (1.2,2) -- (2.12,2.12);
\draw [thick] (1.2,2) -- (1.4,2.65);
\draw [fill] (1.4,2.65) circle [radius = 0.03];
\node [right] at (3,-2) {$\G$};
\node [right] at (2,3) {$\cone^{c}_\infty(\g)$};
\end{tikzpicture}
\caption{The coarse cone type at infinity, picture from~\cite{PSW2}. %P.-S.-Wienhard
}\label{fig:1}
\end{figure}



\subsection{Anosov representations}\label{s.Anosov}
Fix a subset $\t\subset\simple$. Let $\G$ be a finitely generated group and denote by $|\,|$ the word-length associated to a finite generating set $S$.

\begin{defi}\label{AnosovDefi}
Following\footnote{
%See also~\cite{BPS,GGKW}.%Bochi--Potrie--Sambarino
See also~\cite{BPS} and Guéritaud--Guichard--Kassel--Wienhard
\cite{GGKW}.
} Kapovich--Leeb--Porti in~\cite{KLP}, a representation $\rho:\G\to\sfG$ is \emph{$\t$-Anosov} if there exist positive constants $C$ and $\mu$ such that for all $\g\in\G$ and $\sroot\in\t$ one has
\[
\sroot(\cartan(\rho\g))\geq \mu|\g|-C.
\]
The constants $\mu$ and $C$ are usually referred to as \emph{the domination constants} of $\rho$.
If $\sfG=\PGL(d,\R)$ and $\t=\{\slroot_1\}$ we say that $\rho$ is \emph{projective Anosov}. In order to easy the notation we will identify in what follows $\g$ with $\rho(\g)$.
\end{defi}



Anosov representations were introduced by Labourie~\cite{labourie} and further developed by Guichard--Wienhard~\cite{GW}. They have played a central role in understanding the Hitchin component of split groups (see below) and are considered nowadays as the higher-rank generalization of convex co-compact groups. We refer the reader to the surveys by Kassel~\cite{FannyICM18} and Wienhard~\cite{Anna-ICM} for further information.


\begin{rema}\label{wallsAnosov}
A Zariski-dense representation $\rho:\G\to\sfG$ is $\t$-Anosov if and only if $\rho$ is a quasi-isometric embedding and its limit cone $\calL_\rho$ does not meet any wall $\ker\sroot$ for $\sroot\in\t:$ this follows from the definition since by Benoist~\cite{limite}, if $\rho(\G)$ is Zariski-dense then the limit cone $\calL_\rho$ equals the asymptotic cone.
\end{rema}



A useful property of $\t$-Anosov representations is that their limit set $\Lambda_\G\subset \calF_{\theta}$, namely the minimal $\G$-invariant subset in $\calF_{\theta}$, is parametrized by the Gromov boundary of the group $\G$, see~\cite{GGKW,KLP}.
%Kapovich--Leeb--Porti Guéritaud--Guichard--Kassel--Wienhard~\cite{}
We will need the following precise statement.



\begin{prop}[{%Bochi--Potrie--Sambarino~
\cite[Proposition~4.9]{BPS}}]\label{conical}
If $\rho:\G\to\sfG$ is $\t$-Anosov, then for any geodesic ray $(\alpha_n)_0^\infty$ with endpoint $x$, the limits
\begin{align*}
&\xi^\t_\rho(x):=\lim_{n\to\infty}U_\t(\alpha_n),  &&\xi^{\ii\t}_\rho(x):=\lim_{n\to\infty}U_{\ii\t}(\alpha_n)
\end{align*}
exist and do not depend on the ray; they define continuous $\rho$-equivariant transverse maps $\xi^{\t}:\bord\G\to\calF_\t$, $\xi^{\ii\t}:\bord\G\to\calF_{\ii\t}$. If $\g\in\G$ is hyperbolic, then $\g$ is $\t$-proximal with attracting point $\xi^\t(\g^+)=(\g)_\t^+$.
\end{prop}


We conclude the section with a number of quantitative results that will be needed in the paper. For an Anosov representation $\rho$ there exists a constant $\delta_\rho$ quantifying transversality of Cartan-attractors along (quasi)-geodesic rays:

\begin{prop}[{%Bochi-Potrie-S.~
\cite[Lemma~2.5]{BPS}}]
\label{fundamentalConstant}
If $\rho:\G\to\sfG$ is $\t$-Anosov and $c>0$ is given, then there exist $\ov L\in\N$ and $\delta_{\rho,c}>0$, depending only $c$ and the domination constants of $\rho$, such that for every $(c,c)$-quasi-geodesic segment through the identity $\{\alpha_i\}_{-m}^k$ with $k,m\geq L$ one has, for all $\sroot\in\t, $ that
\[
\peso_\sroot\Gr_\t\bigl(U_\t(\rho\alpha_k), \;U_{\ii\t}(\alpha_{-m})\bigr)\geq\log\delta_{\rho,c}.
\]
\end{prop}

Combining Proposition~\ref{conical} and Proposition~\ref{fundamentalConstant} we obtain:
\begin{coro}\label{c.fundamentalConstant}
Up to decreasing $\delta_{\rho,c}$, for every $\g\in\G$ and every $x\in\cone_\infty^c(\g)$ one has
\[
\peso_\sroot\Gr_\t\bigl(\xi^\t_\rho(x),\;U_{\ii\t}(\g^{-1})\bigr)\geq\log\delta_{\rho,c}.
\]
In particular, if we let $\alpha=-\log\delta_{\rho,c}$ then (recall Eq.~\eqref{e.CartBasin})
\begin{equation}\label{e.conesinCartan}
\xi^\theta_\rho(\cone_{\infty}^{c}(\g))\subset B_{\theta,\alpha}(\g).
\end{equation}
\end{coro}

\begin{defi}\label{FConstant}
Let $\rho:\G\to\sfG$ be $\t$-Anosov and $c>0$, then the constant $\delta_{\rho,c}$ verifying both Proposition~\ref{fundamentalConstant} and Corollary~\ref{c.fundamentalConstant} will be called \emph{the fundamental constant} of $\rho$ and $c$. If we consider geodesics instead of quasi-geodesics (i.e. $(c,C)=(1,0)$) we let $\delta_\rho$ be the \emph{fundamental constant} associated to $\rho$.
\end{defi}

The following two results will be needed in Section~\ref{s.6.1}.

\begin{prop}[{cfr.~\cite[Section~5.1]{PSW1}}]%P.-S.-Wienhard~
\label{p.coneinBall}
\label{Lipschitz-compatibleA}
Let $\rho:\G\to\SL(d,\K)$ be projective Anosov and consider $c>0$. Then there exists a constant $K$, depending on $c$ and on $\rho$ such that for every large enough $\g\in\G$ one has
\[
\xi^1_\rho\bigl(\g \cone_\infty^{c}(\g)\bigr)\subset B\left(U_1(\g),\; Ke^{-\slroot_1(\cartan(\g))}\right).
\]
\end{prop}


\begin{proof}
Using Corollary~\ref{c.fundamentalConstant} for $\theta=\{\slroot_1\}$, the result follows as in~ \cite[Section~5.1]{PSW1}.% P.--S.--Wienhard
\end{proof}



\begin{prop}\label{Lipschitz-compatibleB}
Let $\rho:\G\to\SL(d,\K)$ be projective Anosov. For every $\alpha>0$ there exist $C$ and $\mu>0$ such that for every $\ell_1,\ell_2\in\P(\K^d)$ with
\[
\Gr\left(\ell_i,U_{d-1}(\g^{-1})\right) \geq-\alpha, \quad i=1,2
\]
it holds $d_\P(\rho(\g)\ell_1,\rho(\g)\ell_2)\leq Ce^{-\mu|\g|}d(\ell_1,\ell_2)$.
\end{prop}



\begin{proof}
For an Hermitian product on $\C^d$, and every $\alpha>0$ there exists $C>0$ such that if $h\in\GL(d,\C)$ is such that $\slroot_1(\cartan(h))>0$, then for all $\ell_1,\ell_2\in\P(\C^d)$ with {$\angle(\ell_i,U_{d-1}(h^{-1}))>\alpha$} one has
\[
d_\P(h\ell_1,h\ell_2)\leq Ce^{-\slroot_1(\cartan(h))}d_\P(\ell_1,\ell_2),
\]
(a proof follows, for instance, by applying~\cite[Lemma~2.8]{PSW1} to $g=h^{-1}$, $P=U_1(h)$ and $Q=hU_{d-1}(h)$). The result then follows by applying Definition~\ref{AnosovDefi}.
\end{proof}

The following technical result will be useful in the proof of Proposition~\ref{conK}. Given an Anosov representation, we can use the Gromov product to determine the endpoint of a conical sequence (recall Definition~\ref{conicalsequence}):

\begin{lemm}\label{UdU1}
Let $\rho:\G\to\sfG$ be $\t$-Anosov. If $\{\g_n\}\subset\G$ is a conical sequence, $x\in\bord\G$, and there exists $\sroot\in\t$ such that $\peso_\sroot\Gr_\t(\xi^\t(x),U_{\ii\t}(\g_n))\to-\infty$, then $\g_n\to x$.
\end{lemm}

\begin{proof}
We denote by $y$ the endpoint of the conical sequence $\g_n$. Proposition~\ref{conical} and Remark~\ref{r.conicalgeodesic} imply that $U_{\ii\t}(\g_n)\to\xi^{\ii\t}(y)$. Since, however, $\peso_\sroot\Gr_\t(\xi^\t(x),U_{\ii\t}(\g_n))\to-\infty$, we deduce that $\xi^{\ii\t}(y)$ is not transverse to $\xi^\t(x)$ (recall Remark~\ref{fty}). Since $\xi^\t$ is transverse, we deduce that $x=y$.
\end{proof}

It will be useful in the proof of Proposition~\ref{conK} to know that the endpoints of conical sequences belong to pushed Cartan basins:

\begin{lemm}\label{l.conicbasin}
Let $\rho:\G\to G$ be $\theta$-Anosov, $x\in\bord\G$. If $\g_n\to x$ conically, then there exists $\alpha$ only depending on the sequence and the representation $\rho$ such that for every $n$, $\xi^\theta(x)\in \g_nB_{\t,\alpha}(\g_n)$.
\end{lemm}

\begin{proof}
We know from Remark~\ref{r.conicalgeodesic} that $\g_n$ is contained in a neighbourhood of a geodesic ray to $x$, or equivalently there exist a constant $c$ such that $\gamma^{-1}x\in\cone^{c}_\infty(\gamma_n)$. The result is then a consequence of Eq.~\eqref{e.conesinCartan}.
\end{proof}

\subsection{Patterson--Sullivan theory of Anosov representations}\label{rfr}
If $\rho$ is a $\t$-Anosov representation, then we can pullback the Busemann--Iwasawa cocycle of $\sfG$ using the equivariant maps: the \emph{refraction cocycle} associated to a $\t$-Anosov representation $\rho:\G\to\sfG$ is $\rfr:\G\times\bord\G\to\a_\t$ given by
\[
\rfr(\g,x)=\rfr_{\t,\rho}(\g,x)=\bus_\t\bigl(\rho(\g),\;\xi^\t_\rho(x)\bigr).
\]
%Bridgeman--Canary--Labourie-S.
Theorem~1.10 in~\cite{pressure} show that the Mineyev geodesic flow of a group $\G$ admitting an Anosov representations is metric-Anosov, and thus Section~\ref{cocycles} applies to~$\rfr$. Moreover, the following fact places $\rfr$ in the assumptions required in Sections~\ref{Sreal} and~\ref{QQ}, see~\cite{dichotomy} for details.

\begin{enonce}{Fact}
The periods of the refraction cocycle equal the $\theta$-Jordan projection: $\rfr(\g,\g^+)=\lambda_\t(\g)$. For any $\sroot\in\t$ the real valued cocycle $\peso_\sroot\rfr$ has finite entropy.
\end{enonce}

We import the following concepts of cocycles to the setting of Anosov representations:
\begin{itemize}
\item The limit cone of $\rfr$ will be denoted by $\calL_{\t,\rho}$ and referred to as \emph{the $\t$-limit cone of $\rho$}; it is the smallest closed cone that contains the projected Jordan projections $\{\lambda_\t(\g):\g\in\G\}$.

\item The \emph{interior of the dual cone} $\inte(\calL_{\t,\rho})^*\subset\a_\t^*$ consists of linear forms whose \emph{entropy}
\[
\hJ\varphi= \lim_{t\to\infty}\frac1t\log\#\bigl\{[\g]\in[\G]:\varphi(\lambda_\t(\g))\leq t\bigr\}
\]
is finite.

\item The $\t$-\emph{critical hypersurface}, resp. $\t$-\emph{convergence domain}, of $\rfr$ will be denoted by
\begin{align*}\label{qh2}
\calQ_{\t,\rho} & =\left\{\varphi\in\inte\cdc{\t,\rho}:\hJ\varphi=1\right\},\\
\cal D_{\t,\rho} & = \left\{\varphi\in\inte\cdc{\t,\rho}:\hJ\varphi\in(0,1)\right\}.
\end{align*}
\item If $\calL_{\t,\rho}$ has non-empty interior, then we have a \emph{duality} diffeomorphism between $\calQ_{\t,\rho}$ and $\inte\P(\calL_{\t,\rho})$ given by
\[
\varphi\mapsto\sf u_\varphi=\sf T_\varphi\calQ_\rho.
\]
\end{itemize}
Observe that $\sf u_\varphi\in\inte\P(\calL_{\t,\rho})$ and $\varphi|\calL_{\t,\rho}-\{0\}>0$ so $\varphi(\sf u_\varphi)\neq0$. More information on these objets can be found on~\cite[Section~5.9]{dichotomy}.

\begin{rema}\label{r.entropy=crit}
It it proven in~\cite[Theorem~2.31$\MK$(2)]{GMT}
%Glorieux--Monclair--Tholozan
(see also~\cite[Corollary~5.5.3]{dichotomy}) that if $\rho$ is $\t$-Anosov then for every $\varphi\in\inte(\calL_{\t,\rho})^*$ the entropy $\hJ\varphi$ equals the critical exponent
\[
\hC\varphi:= \lim_{t\to\infty}\frac1t\log\#\left\{\g\in\G:\varphi(\cartan(\g))\leq t\right\}.
\]
In particular the $\t$-convergence domain is also given by
\[
\cal D_{\t,\rho} =\left\{\varphi\in (\a_\t)^*:\sum_{\g\in\G}e^{-\varphi(\cartan(\g))}<\infty\right\},
\]
see~\cite[Section~5.7.2]{dichotomy}.
\end{rema}

The following is essentially a consequence of~\cite[Proposition~3.3.3]{Quint-Div},
%Quint's,
we outline a~proof.


\begin{coro}\label{quintgrowththeta}
Consider a norm $N$ on $\a_\t$, then the critical exponent $h^N_\rho$ of the series $s\mapsto\sum_{\g\in\G}e^{-sN(\cartan_\t(\g))}$ is given by
\[
h^N_\rho=\inf\left\{N^*(\varphi):\varphi\in\calQ_{\t,\rho}\right\},
\]
where $N^*$ is the associated dual norm on $\a_\t^*$.
\end{coro}

\begin{proof}
We consider the counting measure $\nu=\sum_{\g\in\G}\delta_{\cartan_\t(\g)}$ on $\a_\t$. We then have, in the notation of~\cite[Section~3]{Quint-Div}, that $\tau^N_\nu=h^N_\rho$ and, by Remark~\ref{r.entropy=crit}, $\sigma^N_\nu=\inf_{\varphi\in\calQ_{\t,\rho}}N^*(\varphi)$. Thus, in order to deduce the corollary from~\cite[Proposition~3.3.3]{Quint-Div} it is enough to verify that the counting measure $\nu$ is of concave growth as in~\cite[Section~3.2]{Quint-Div}. In turn this is a consequence of Lemma~\ref{l.Quint} below, an adaptation of~\cite[Proposition~2.3.1]{Quint-Div} (see also~\cite[Lemma~3.8]{KimOhWang}
%Kim--Oh--Wang
where similar arguments are explained for the $\a_\t$ counting measure).
\end{proof}

\begin{lemm}\label{l.Quint}
Let $\|\,\|$ be a norm on $\a_\t$. Let $\grupo<\sfG$ be Zariski-dense and $\t$-Anosov. Then there exists a product map $m:\grupo\times\grupo\to\grupo$ with the following properties:
\begin{enumerate}
\item\label{lemm4.16.1} there exists a real number $\kappa\geq 0$ such that, for all $\g_1,\g_2\in\grupo$,
\[
\bigl\|\cartan_\t(m(\g_1,\g_2))-\cartan_\t(\g_1)-\cartan_\t(\g_2)\bigr\|\leq\kappa;
\]
\item\label{lemm4.16.2} for every real $R\geq 0$ there exists a finite subset $H$ of $\grupo$ such that, for $\g_1,\g_2\g_1',\g_2'$ in $\grupo$ with $\|\cartan_\t(\g_i)-\cartan_\t(\g_i')\|<R$ for $i=1,2$, then
\[
m(\g_1,\g_2)=m(\g_1',\g_2') \implies \g_i'\in\g_i H, \quad\text{for }i=1,2.
\]
\end{enumerate}
\end{lemm}

\begin{proof}
It is enough to consider the \emph{generic product} map $\pi:\grupo\times\grupo\to\grupo$ constructed in~\cite[Proposition~2.3.1]{Quint-Div}, which satisfies the analogous properties with respect to the Cartan projection $\cartan:\sfG\to\a$ and a norm $\|\,\|$ on $\a$. The first property is satisfied since we can assume that the projection $\pi_\theta:\a\to \a_\t$ is norm non-increasing. The second follows from the Anosov property: by the construction in~\cite[Proposition~2.3.1]{Quint-Div} one can choose $H$ to be the set of elements $\g$ such that $\|\cartan_\t(\g)\|<R'$ for some $R'$ depending on $R$. Such set is finite because, by definition of $\Pi$, there exists $R''$ depending on $R'$ and the norm $\|\,\|$ such that if $\|\cartan_\t(\g)\|<R'$ then $\sroot(\cartan(\gamma))<R''$, which in turn implies by Definition~\ref{AnosovDefi} that $|\gamma|<R''/\mu+C$, and thus $\g$ belongs to a finite subset.
\end{proof}


We observe that for $\varphi\in\inte(\calL_{\t,\rho})^*$ Assumptions~\ref{Pat-SulEx} are guaranteed for $\rfr_\varphi:=\varphi\circ\rfr$. Indeed the cocycle
\[
\ov\rfr{}(\g,x)=\ii\bus_{\ii\t}\big(\g,\xi^{\ii\t}(x)\big)
\]
is dual to $\rfr$, from Eq.~\eqref{forGr} the function $[\cdot,\cdot]_\varphi:\bord^2\G\to\R$
\[
[x,y]_\varphi=\varphi\Bigl(\Gr_\t\left(\xi^\t(y),\xi^{\ii\t}(x)\right)\Bigr)
\]
is a Gromov product for the pair $(\ov\rfr_{\varphi},\rfr_{\varphi})$, and we have the following result guaranteeing existence of Patterson--Sullivan measures $\ps^\varphi$ and $\ov\ps^\varphi$, as well as their values on Cartan basins defined in Eq.~\eqref{e.CartBasin}.

\begin{coro}[{\cite[Corollary~5.5.3, Lemma~5.7.1]{dichotomy}}]
\label{existe}
For every $\varphi\in\inte(\calL_{\t,\rho})^*$ there exists a $\rfr_\varphi$-Patterson--Sullivan measure $\ps^\varphi$ of exponent $\hJ\varphi$, moreover for every $\alpha$ there exists a constant $C$ such that for every $\g\in\G$ one has
\[
\ps^\varphi\left((\xi^\theta)^{-1}(\g B_{\theta,\alpha}(\g))\right) \leq Ce^{-\hJ\varphi\varphi(\cartan(\g))}.
\]
\end{coro}



\subsection{Subspace-conicality}\label{sandwich}

In this section we are interested in a notion of conicality along higher dimensional subspaces of the ambient Levi space.

\begin{defi}\label{W-conical}
Let $\rho:\G\to\sfG$ be $\t$-Anosov and consider a subspace $\sf W\subset\a_\t$. A point $x\in\bord\G$ is \emph{$\sf W$-conical} if there exists a conical sequence $\{\g_n\}_0^\infty\subset\G$ converging to $x$, a constant $K$ and $\{w_n\}_0^\infty\subset\sf W$ such that for all $n$ one has
\[
\big\|\cartan_\t\big(\g_n)- w_n\big\|\leq K.
\]
The set of such points will be denoted by $\bord_{\sf W,\rho}\G=\bord_{\sf W}\G$.
\end{defi}


Assume from now on that $\sf W$ intersects the relative interior of $\calL_{\t,\rho}$, and consider $\varphi\in\calQ_{\t,\rho}$ with $\sf u_\varphi\subset\sf W$. The intersection $\sf W_\varphi=\sf W\cap\ker\varphi$ has co-dimension~$1$ in $\sf W$ and has trivial intersection with the limit cone $\calL_{\t,\rho}$. Consider the quotient space
\[
V=\a_\t/\sf W_\varphi
\]
equipped with the quotient projection $\Pi:\a_\t\to V$. We say that $\rho$ is $(\sf W,\varphi)$-\emph{non-arithmetic} if the group spanned by $\{\Pi(\lambda_\t(\g)):\g\in\G\}$ is dense in $V$. In this section we prove the following.

\begin{theo}\label{hyper}
\label{ps-sandwich}
Let $\rho:\G\to\sfG$ be $\t$-Anosov, $\sf W$ be a subspace of $\a_\t$ intersecting non-trivially the relative interior of $\calL_{\t,\rho}$, and $\varphi\in(\a_\t)^*$ with $\sf u_\varphi\subset \sf W$. Assume $\rho$ is $(\sf W,\varphi)$-non-arithmetic, then:
\begin{itemize}
\item if $\sf W$ has codimension~$1$ then $\mu^\varphi(\bord_\sf W\G)=1;$
\item if $\codim\sf W\geq3$ then $\mu^\varphi(\bord_{\sf W}\G)=0$.
\end{itemize}
\end{theo}

\begin{rema}
If $\rho$ is Zariski-dense then Theorem~\ref{densidad} (\cite{benoist2}) guarantees $(\sf W,\varphi)$-non-arithmeticity for every $\varphi\in(\a_\t)^*$ with $\sf u_\varphi\in\P(\sf W)$, thus Theorem~\ref{hyper} readily implies Theorem~\ref{Ahyper}.
\end{rema}

The remainder of the section is devoted to the proof of Theorem~\ref{hyper}. Let
\[
V^*=\Ann(\sf W_\varphi)=\left\{\psi\in(\a_\t)^*\,:\,\psi|\sf W_\varphi\equiv0\right\},
\]
with a slight abuse of notation we will identify the dual of $V$ with $V^*\subset(\a_\t)^*\subset \a^*$ (recall from Section~\ref{s.Levi} that we are identifying $(\a_\t)^*$ with the subspace of $\a^*$ consisting of $\pi_\t$-invariant linear forms).

The composition of the refraction cocycle of $\rho$ with $\Pi$ is a $V$-valued H\"older cocycle $\cc:\G\times\bord\G\to V$,
\[
\cc(\g,x)=\Pi(\rfr(\g,x)).
\]
Its periods are $\cc(\g,\g_+)=\Pi(\lambda_\t(\g))$, and thus its limit cone is $\calL_\cc=\Pi(\calL_{\t,\rho})$. By $(\sf W,\varphi)$-non-arithmeticity, $\cal L_{\cc}\subset V$ has non-empty interior.


The heart of the proof of Theorem~\ref{hyper} consists on relating $(\sf W,\varphi)$-conical points with elements of $\widetilde{\cal K}(\df^{\varphi})$, where $\df^{\varphi}$ is the directional flow on $\G\/\bord^2\G\times V$ associated to the cocycle $\cc$ as in Section~\ref{directional}. The first step is thus to observe that we can apply Corollary~\ref{d=2} to $\cc$, a task we enter at this point.



Since $\varphi\in\calQ_{\t,\rho}$, it has in particular finite entropy. Moreover, by definition of $V^*$ one has $\varphi\in V^*$. Consequently, the cocycle $\cc$ verifies assumptions in Corollary~\ref{strictly}. One can moreover transfer existence properties from $\rfr$ to $\cc$, indeed one has the following.



\begin{prop}
The cocycle $\ovcc=\Pi\circ\ov\rfr$ is a dual cocycle for $\cc$. For each $\psi\in\calQ_\cc$ there exist Paterson--Sullivan measures for $\cc$ and $\ov \cc$ and the projection $\psi(\Pi([\cdot,\cdot]))$ is a Gromov product for the pair $\psi\circ\cc,\psi\circ\ovcc$.
\end{prop}

\begin{proof}
Since $\psi\in\calQ_\cc=\calQ_{\t,\rho}\cap V^*$ we can apply Corollary~\ref{existe} to $\psi$ to obtain the desired Patterson--Sullivan measure, the remaining statements follow trivially as the equalities are linear.
\end{proof}

Since we are assuming $(\sf W,\varphi)$-non-arithmeticity, the cocycle $\cc$ is non-arithmetic and thus Corollary~\ref{d=2} gives the following dynamical information, observe that $\dim V=\codim \sf W+1$.

\begin{coro}\label{dicoCC}
If $\codim \sf W\leq1$ then the directional flow $\df^\varphi$ is $\BM^\varphi$-ergodic, in particular $\cal K(\df^\varphi)$ has total mass. If $\codim \sf W\geq3$ then $\cal K(\df^\varphi)$ has measure zero.
\end{coro}

Observe that modulo the understood identifications $\calQ_\cc=\calQ_{\t,\rho}\cap V^*$, hence
\[
\sf T_\varphi\calQ_\cc=(\sf T_\varphi\calQ_{\t,\rho})\cap V^*
\]
and thus the map $\sf u^\cc:\calQ_\cc\to\inte\P(\cal L_\cc)$ from Corollary~\ref{strictly} verifies $\sf u^\cc_\varphi=\Pi(\sf u_\varphi)$. So measuring $\sf W$-conicality with respect to $\mu^\varphi$ translates to directional conicality along the direction $\sf u^\cc_\varphi$, which we now recall. We fix an arbitrary norm $\|\,\|$ on $V$ and define, for $\ell\in\P(V)$ and $r>0$, the $r$-\emph{tube about $\ell$} by
\[
\T_r(\ell):=\left\{v\in V\,\middle|\, \exists w\in\ell, \|v-w\|<r\right\}.
\]

\begin{defi}\label{defCon}
A point $y\in\bord\G$ is \emph{$\sf u_\varphi^\cc$-conical} if there exists $r>0$ and a~conical sequence $\{\g_n\}_0^\infty\subset\G$ with $\g_n\to y$ such that for all $n$ one has $\Pi(\cartan_\theta(\rho(\g_n)))\in\T_r(\sf u_\varphi^\cc)$.
\end{defi}

The next statement follows from the definitions.

\begin{lemm}
A point $y\in\bord\G$ is $\sf W$-conical if and only if it is $\sf u^\cc_\varphi$-conical.
\end{lemm}


If we are allowed to worsen the constants, we can replace, in Definition~\ref{defCon}, the conical sequence $(\g_n)$ with an infinite subset of a geodesic ray:

\begin{lemm}\label{conGeo}
A point $y\in\bord\G$ is $\sf u_\varphi^\cc$-conical if and only if there exists $r>0$, a~geodesic ray $(\alpha_i)_0^\infty$ converging to $y$ and an infinite set of indices $\I\subset\N$ such that for all $k\in\I$ one has
\[
\Pi(\cartan_\t(\alpha_k))\in\T_r\left(\sf u_\varphi^\cc\right).
\]
\end{lemm}


\begin{proof}
Assume $y$ is $\sf u_\varphi^\cc$-conical, then since $\{\g_n\}_0^\infty$ is conical, for any geodesic ray $(\alpha_n)_0^\infty$ converging to $y$ there exists $K>0$ and a subsequence $\{\alpha_{n_k}\}$ such that for all $k$ one has $d_\G(\alpha_{n_k},\g_{k})<K$ (Remark~\ref{r.conicalgeodesic}). Proposition~\ref{p.Ben} implies then that for all $k$ one has
\[
\|\cartan(\alpha_{n_k})-\cartan(\g_k)\|
\]
is bounded independently of $k$. This implies the result.
\end{proof}

We now relate $\sf u^\cc_\varphi$-conicality with the recurrence set $\cal K(\df^\varphi)$. By definition of $\cal K(\df^\varphi)$, a point $(x,y,v)\in\bord^2\G\times V$ projects to $\cal K(\df^\varphi)$ if and only if there exist divergent sequences $(\g_n)\subset\G$ and $t_n\to+\infty$ in $\R$ such that
\begin{equation}\label{sequence}
\df^{\varphi}_{t_n}\g_n^{-1}(x,y,v)=\left(\g_n^{-1}x,\g_n^{-1}y,v-\cc\left(\g_n^{-1},y\right)-t_nu_{\varphi}\right)
\end{equation}
is contained in a subset of the form $\{(z,w)\in\bord^2\G:d(z,w)\geq\kappa\}\times B(v,K)$ for some distance $d$ on $\bord\G$. One has the following


\begin{prop}\label{conK}
A point $y\in\bord\G$ is $\sf u_\varphi^\cc$-conical if and only if for every $x\in\bord\G-\{y\}$ and $v\in V$ the point $(x,y,v)$ projects to $\cal K(\df^\varphi)$.
\end{prop}



\begin{proof}
The implication ($\implies$) follows exactly as in the proof of~\cite[Proposition~5.13.4]{dichotomy}. The other implication also follows similarly but with a minor difference we now explain.

Assume that $(x,y,v)$ projects to $\cal K(\df^\varphi)$ and consider sequences $\{\g_n\}$ and $t_n$ as in Eq.~\eqref{sequence}. Since $(\g_n^{-1}x,\g_n^{-1}y)$ remains in a compact subset of $\bord^2\G$, the sequence $\{\g_n\}$ is conical, we will show now that $\g_n\to y$. Indeed, since $t_n\to+\infty$ necessarily $\cc(\g_n^{-1},y)\to-\infty$.

Consider now any root $\sroot\in\t$, with associated fundamental weight $\peso_{\sroot}\in(\calL_{\t,\rho})^*$, and Tits representation $\Fund_{\sroot}:\sfG\to V$. Since $\rho$ is $\theta$-Anosov, the H\"older cocycle $\rfr_{\peso_{\sroot}}$ has positive periods and finite entropy. Since $\cc(\g_n^{-1},y)\to-\infty$ Proposition~\ref{-infty} implies that
\[
\rfr_{\peso_{\sroot}}\left(\g_n^{-1},y\right)\to-\infty.
\]
By definition of the cocycle $\rfr_{\peso_{\sroot}}$ and Eq.~\eqref{busnorma} we have
\begin{equation}\label{normto0}
\frac{\left\|\Fund_{\sroot}\left(\g_n^{-1}\right)v\right\|}{\|v\|}\to0
\end{equation}
for a non-zero $v\in\zeta_\sroot(\xi(y))$, (recall that the map $\zeta_\sroot:\calF_{\sroot}(\sfG)\to \P(V)$ was defined in Eq.~\eqref{maps}). Setting $\dim V=d$, a standard linear algebra computation (for example~\cite[Lemma~A.3]{BPS})
%in Bochi--Potrie-S.
gives
\begin{align*}
\frac{\left\|\Fund_{\sroot}\left(\g_n^{-1}\right)v\right\|}{\|v\|} & \geq\left\|\Fund_{\sroot}\left(\g_n^{-1}\right)\right\|\sin\angle\left(\zeta_\sroot\xi(y),U_{d-1}(\Fund_{\sroot}\g_n)\right)\\
& \geq e^{l_\sroot\peso_\sroot\Gr_\t\left(y,U_{\ii\t}(\g_n)\right)}
\end{align*}
and thus, by Eq.~\eqref{normto0} and Lemma~\ref{UdU1} one has $\g_n\to y$, as desired.

The point $\xi(y)$ lies then in the pushed Cartan basin $\g_nB_{\t,\alpha}(\g_n)$ for an $\alpha$ independent of $n$ (Lemma~\ref{l.conicbasin}), and thus Eq.~\eqref{comparision-shadow} gives a constant $K$ such that for all $n$ one has
\[
K\geq\left\|\cartan_\t(\g_n)-\rfr\left(\g_n,\g_n^{-1}y\right)\right\|=\left\|\cartan_\t(\g_n)+\rfr\left(\g_n^{-1},y\right)\right\|
\]
implying, by Eq.~\eqref{sequence}, that $y$ is $\sf u_\varphi^\cc$-conical, as desired.
\end{proof}

The proof of Theorem~\ref{hyper} follows now along the same lines as in~\cite[Theorem~5.13.3]{dichotomy}. We include the arguments here for completeness.

For $y\in\bord_{\sf W,\rho}\G, x\in\bord\G-\{y\}$ we fix neighbourhoods $A^-$ and $A^+$ of $x$ and $y$ respectively and $T>0$ small enough so that the quotient projection $\sf p:\bord^2\G\times V\to\G\backslash\bord^2\G\times V$ is injective on $\tilde{\rmB}=A^-\times A^+\times B(0,T)$. We can thus use Eq.~\eqref{forBM} to compute the measure of $\rmB=\sf p(\tilde{\rmB})$.

For $\tilde{\cal K}({\df}^\varphi)=\sf p^{-1}(\cal K({\df}^\varphi))$, Proposition~\ref{conK} asserts
\[
A^-\times\left(A^+\cap\bord_{\sf W,\rho}\G\right)\times B(0,T)=\tilde{\cal K}\left({\df}^\varphi\right)\cap\tilde{\rmB}.
\]
If $\codim \sf W=1$ by Corollary~\ref{dicoCC} $\BM^\varphi(\tilde{\rmB})=\BM^\varphi(\tilde{\cal K}({\df}^\varphi)\cap\tilde{\rmB})$, which implies that $\mu^\varphi(A^+\setminus\bord_{\sf W,\rho}\G)=0$ and thus $\mu^\varphi(\bord_{\sf W,\rho}\G)=1$. On the other hand, if $\codim\sf W\geq3$, then we have $\BM^\varphi(\tilde{\cal K}({\df}^\varphi))=0$ so $\mu^\varphi(A^+\cap\bord_{\sf W,\rho}\G)=0$ and the theorem is proved.



\section{Locally conformal representations: Hausdorff dimension of \texorpdfstring{$\hol$}{hol}-conical points}\label{s.Hffbcon}
In this section we let $\K=\R, \C$ or $\H$, the non-commutative field of Hamilton's quaternions. A Cartan subspace $\a$ of $\SL(d,\K)$ is the subspace of $\R^d$ consisting of vectors whose coordinates sum 0. For $g\in\SL(d,\K)$ we denote by
\[
\cartan(g)=(\cartan_1(g),\dots,\cartan_d(g))\in\a^+
\]
the coordinates of the Cartan projection. We recall Definition~\ref{d.hyp}.

\begin{defi}\label{locConf}
Let $p\in\lb2,d-1\rb$. A $\{\slroot_1,\slroot_{d-p}\}$-Anosov representation $\rho:\G\to\SL(d,\K)$ is \emph{$(1,1,p)$-hyperconvex} if, for every pairwise distinct triple $(x,y,z)\in\bord\G^{(3)}$, one has
\[
\left(\xi^1(x)+\xi^1(y)\right)\cap\xi^{d-p}(z)=\{0\}.
\]
If in addition one has $\cartan_2(\rho(\g))=\cartan_p(\rho(\g))$ $\forall\g$, we say that $\rho$ is \emph{locally conformal}. As before, we identify from now on $\g$ and $\rho(\g)$.
\end{defi}

The terminology is justified by Proposition~\ref{conetypesBalls} below stating that for such representations pushed coarse cone types are coarsely balls, a small refinement of an analogous result from~\cite{PSW1}. %P.-S.-Wienhard

In this section we will study conicality from Section~\ref{sandwich} on a specific situation that we now explain. Later, in Section~\ref{teoLCdiff}, we will relate this section to the notion of $\hol$-concavity and in Section~\ref{s.Hffproof} to differentiability properties of the map $\ovxi\circ\xi^{-1}$.


Consider $\bar{\K}\in\{\R,\C,\H\}$ and two locally conformal representations $\rho:\G\to\SL(d,\K)$ and $\ovrho:\G\to\SL(\ovd,\bar{\K})$, with projective equivariant maps
\begin{align*}
\xi:\bord\G &\to\P\left(\K^d\right)\\
\ovxi:\bord\G &\to\P\left(\bar{\K}{}^{\ovd}\right).
\end{align*}
The product representation $(\rho,\ovrho):\G\to\SL(d,\K)\times\SL(\ovd,\bar{\K})$ is $\t$-Anosov for $\t=\{\slroot_1,\slroot_p,\ov{\slroot}_1,\ov{\slroot}_p\}$ with $\{\slroot_1,\ov{\slroot_1}\}$-limit map the ``graph map''
\[
\varXi=\left(\xi,\ov{\xi}\right):\bord\G\to\P\left(\K^{d}\right)\times\P\left(\bar{\K}^{\ovd}\right).
\]
We consider a Cartan subspace of the product group $\SL(d,\K)\times\SL(\ovd,\bar{\K})$ and let $\a_\t$ be the associated Levi space. Its dual $(\a_\t)^*$ is spanned by the fundamental weights of roots in $\t$. We let
\begin{align*}
\slroot &:=\frac{ p\peso_{\slroot_1}-\peso_{\slroot_p}}{p-1},\\
\ovsloot &:=\frac {p\peso_{\ovsloot_1}-\peso_{\ovsloot_p}}{p-1}.
\end{align*}
Both $\slroot,\ovsloot\in(\a_\t)^*$ and under the assumption $\cartan_2(\g)=\cartan_p(\g)$ for all $\g$ of Definition~\ref{locConf}, it holds on $\calL_{\rho} $ that $\slroot_1=\slroot$ and $\ovsloot=\ov{\slroot}_1$ (if $p=2$ the equality holds on $\a$).


\begin{defi}\label{flat}
Fix $\hol\in(0,1]$. A point $x\in\bord\G$ is \emph{$\hol$-conical} if it is conical as in Definition~\ref{W-conical} for the product representation $(\rho,\ovrho)$ with respect to the hyperplane
\[
\{v\in\a_\t:\hol\slroot(v)=\ovsloot(v)\}=\ker(\hol\slroot-\ovsloot).
\]
Equivalently, there exist $R$, a geodesic ray $(\alpha_j)_0^\infty\subset\G$ with $\alpha_j\to x$, and a subsequence $\{j_k\}$ such that for all $k$ one has
\[
\bigl|\hol\slroot\left(\cartan(\alpha_{j_k})\right)-\ov{\slroot}\left(\cartan(\ov\alpha_{j_k})\right)\bigr|\leq R.
\]
\end{defi}



Consider also the critical exponent
\[
\hC{\infty,\hol}=\lim_{t\to\infty}\frac 1t\log\#\bigl\{\g\in\G:\max\left\{\hol\slroot(\cartan(\g)),\ovsloot(\cartan(\ov\g))\right\}\leq t\bigr\},
\]
and recall from Eq.~\eqref{defiII} the dynamical intersection defined by
\begin{equation}\label{IIreps}
\II_{\slroot}(\ovsloot)=\lim_{t\to\infty}\frac1{\#\sf R_t(\slroot)}\sum_{\g\in \sf R_t(\slroot)}\frac{\ovsloot(\lambda(\ov\g))}{\slroot(\lambda(\g))},
\end{equation}
where $\sfR_t(\slroot)=\{[\g]\in[\G]:\slroot(\lambda(\g))\leq t\}$.

In this section we compute the Hausdorff dimension of the image under the graph map $\varXi$ of the set of $\hol$-conical points with respect to a Riemannian metric:

\begin{theo}\label{Hffconical}
Let $\rho,\ovrho$ be locally conformal representations over $\K$ and $\bar{\K}$ respectively. Assume the group generated by $\{(\slroot(\lambda(\g)),\ovsloot(\lambda(\ov\g))):\g\in\G\}$ is dense in~$\R^2$. Then, for every $\hol\in(0,1]$ with
\[
\II_{\slroot}(\ovsloot)>\hol>1/\II_{\ovsloot}(\slroot),
\]
one has
\begin{align*}
\hol\hC{\infty,\hol}\leq\Hff\varXi(\{\hol\text{-conical\ points}\})
&\leq\min\left\{\hC{\infty,\hol},\hol\hC{\infty,\hol}+(1-\hol) \right\}\\
&<\min\left\{\hJ{\ovsloot},\hJ{\slroot}/\hol\right\}\\
&\leq\Hff(\varXi(\bord\G))\\
& =\max\left\{\hJ{\slroot},\hJ{\ovsloot}\right\}.
\end{align*}
\end{theo}

The proof of the above result is completed in Section~\ref{proofHffconical}.

Recall that if $\hC{\slroot_1}=\hC{\ovsloot_1}$ and the representations are not gap-isospectral, then Proposition~\ref{ineq} gives $\II_{\ovsloot_1}(\slroot_1)>1$. Theorem~\ref{Hffconical} studies then $\hol$-conical points for any $\hol$ with $\II_{\ovsloot_1}(\slroot_1)>1/\hol\geq1$. As the following result shows, the equality between entropies is rather natural for $\K=\R$.

\begin{theo}[{\cite{PSW1}}]\label{hyperh=1}
Let $\rho:\G\to\SL(d,\K)$ be locally conformal, then
\[
\hJ{\slroot}=\Hff(\xi(\bord\G)).
\]
Moreover, when $\K=\R$ and $\bord\G$ is homeomorphic to a $p-1$-dimensional sphere, $\scr h_{\slroot}=p-1$.
\end{theo}



When $\G$ is a surface group we can also weaken the assumption on the density of periods:

\begin{coro}\label{flatsurface}
Assume $\bord\G$ is homeomorphic to a circle and let $\rho$ and $\ovrho$ be non-gap-isospectral real $(1,1,2)$-hyperconvex representations of $\G$. Then
\[
\Hff\varXi\left(\{1\text{-conical\ points}\}\right)=\hC{\infty,1}<1.
\]
\end{coro}


\begin{proof}
Proposition~\ref{nonA} below states that under our assumptions the group generated by $\{(\slroot(\lambda(\g)),\ovsloot(\lambda(\ov\g))):\g\in\G\}$ is dense in $\R^2$. Theorem~\ref{hyperh=1} guarantees that $\hC{\slroot}=\hC{\ovsloot}$ and Proposition~\ref{ineq} then gives $\II_{\ovsloot}(\slroot)>1$. The equality thus follows from Theorem~\ref{Hffconical}.
\end{proof}


Kim--Minsky--Oh~\cite{KMO-HffDir} have established related Hausdorff dimension computations when $\rho$ and $\ovrho$ are convex co-compact representations in $\SO(n,1)$ without any assumption on $\II$.



\subsection{Cone types are coarsely balls}

In~\cite{PSW1} Pozzetti, Sambarino and Wienhard gave a concrete description of the images under the boundary map of the cone types at infinity. We discuss here a slight extension of that result adapted to our needs. We denote by $d_\P$ the distance on $\P(\K^d)$ induced by the choice of an inner (Hermitian) product on $\K^d$ and by $B(\ell,r)\subset\P(\K^d)$ the associated ball of radius $r$ about $\ell$.

\begin{prop}\label{conetypesBalls}
Let $\rho:\G\to\SL(d,\K)$ be locally conformal. Then there exist positive constants $ c, k_1,k_2$ and $L\in\N$ such that for every $x\in\bord\G$, every geodesic ray $(\alpha_j)_0^\infty$ with endpoint $x$ and every $j>L$ one has
\[
B\left(\xi(x),k_1e^{-\slroot_1(\cartan(\alpha_j))}\right)\cap\xi(\bord\G)\subset\xi\left(\alpha_j\cone_\infty^{c}(\alpha_j)\right) \subset B\left(\xi(x),k_2e^{-\slroot_1(\cartan(\alpha_j))}\right).
\]
\end{prop}

\begin{proof}
The desired inclusions are proven in~\cite{PSW1} for thickened cone types at infinity. We briefly explain here how to deduce from it the result we need.

Following~\cite{PSW1} we denote by $X_\infty(\g)$, for $\g\in\G$, the \emph{thickened cone type at infinity}, namely the tubular neighborhood in $\P(\K^d)$ of $\xi(\cone_\infty(\g))$ of radius $\delta_\rho/2$, where $\delta_\rho$ is the fundamental constant from Definition~\ref{FConstant}. In~\cite[Corollary~5.10]{PSW1} it is established that there exist $c_1>0$ and $L_0>0$ only depending on the domination constants of $\rho$ such that for all $j\geq L_0$ one has
\[
B\left(\xi(x),c_1e^{-\slroot_1(\cartan(\alpha_j))}\right)\cap\xi(\bord\G)\subset \alpha_iX_\infty(\alpha_i).
\]

By definition the thickened cone type $X_\infty(\g)$ is contained in the Cartan basin $B_{\{\slroot_1\},\alpha}(\g)$ for $\alpha=-2\log\delta_\rho$. So~\cite[Proposition~3.3]{PSW2} provides the existence of $c$ and $L_0$ such that for $\g\in\G$ with $|\g|>L_0$, one has
\[
X_\infty(\g)\cap\xi(\bord\G)\subset\xi\left(\cone_\infty^{c}(\g)\right).
\]
Combining both equations one has, for all $j\geq L_0$ that
\begin{equation}\label{bola}
B\left(\xi(x),c_1e^{-\slroot_1(\cartan(\alpha_j))}\right)\cap\xi(\bord\G)\subset\xi\left(\alpha_j\cone_\infty^{c}(\alpha_j)\right) \subset B\left(\xi(x), Ke^{-\slroot_1(\cartan(\alpha_j))}\right),
\end{equation}
the second inclusion following from Proposition~\ref{p.coneinBall}. This concludes the proof.
\end{proof}



\subsection{Hausdorff dimension and related concepts} Recall that, given a metric space $(X,d)$ and a real number $s>0$, the \emph{$s$-capacity} of $X$ is
\[
\cal H^s(X,d)=\lim_{\epsilon\to 0}\inf\left\{\sum_{U\in\cal U} \diam U^s\,\middle|\, \cal U \text{ open covering of $X$, } \sup_{U\in\cal U}\diam U<\epsilon\right\}
\]
and that
\begin{equation}\label{HffDef}
\Hff(X)=\inf \left\{s\,\middle|\, \cal H^s(X)=0\right\}=\sup \left\{s\,\middle|\,\cal H^s(X)=\infty\right\}.
\end{equation}
The following can be verified directly from the definition:

\begin{lemm}\label{l.Hffunion}
If $X=\bigcup_{n\in\N} X_n$ then
\[
\Hff(X)=\sup \Hff(X_n).
\]
\end{lemm}

We will use the following consequence of~\cite[Theorem~1.5.14]{Edgar} from Edgar's book:


\begin{coro}\label{edgar}
Let $E\subset\R^d$ be a measurable subset equipped with a probability measure $\nu$. If the upper density
\[
\ov{D}^\alpha(x)=\limsup_{r\to0}\frac{\nu\left(B(x,r)\cap E\right)}{r^\alpha}
\]
is $\nu$-essentially bounded above, then $\Hff(E)\geq\alpha$.
\end{coro}

\subsection{The lower bound \texorpdfstring{$\Hff(\varXi\{\hol\text{-conical\ points}\})\geq\scr{b}\hC{\infty,\hol}$}{Hff(varXi{ hol-conical points})>= scrb hC infty, hol}}\label{s.lower}


We import some tools from the proof of Theorem~\ref{hyper}. Consider the vector space
\[
V^*:=\spa\{\slroot,\ov{\slroot}\}
\]
together with its radical $\ann(V^*)=\ker\slroot\cap\ker\ovsloot$ and the quotient vector space $V=\a_\t/\ann(V^*)$. Any element of $V^*$ vanishes on $\ann(V^*)$ and thus $V^*$ is naturally identified with the dual space of $V$. Using the preferred basis $\{\slroot,\ovsloot\}$ of $V^*$ we identify $V$ and $\R^2$ via the isomorphism $v\mapsto (\slroot(v),\ovsloot(v))$ and we let
\[
\Pi:\a_\t\to\R^2
\]
be the quotient projection (composed with the above isomorphism). The image of the hyperplane $\ker\hol\slroot-\ovsloot$ under the composition of $\Pi$ and the identification of $V$ with $\R^2$ is the line passing through $(1,\hol)$,
\[
\Pi\left(\ker(\hol\slroot-\ovsloot)\right)=\left\{v\in V:\hol\slroot(v)=\ovsloot(v)\right\}.
\]
We consider the quadrant
\[
V^+=\{\slroot\geq0\}\cap\{\ov{\slroot}\geq0\}.
\]
Let $\cc=\cc_{(\rho,\ovrho)}:\G\times\bord\G\to V$ be the composition of the refraction cocycle $\rfr_{(\rho,\ovrho)}$ of the pair with $\Pi$. Its periods are
\[
\cc(\g,\g_+)=\left(\slroot(\lambda(\g)),\ovsloot(\lambda(\ov\g))\right),
\]
so by assumption $\cc$ is non-arithmetic. As in Section~\ref{sandwich} one has $\calQ_\cc=V^*\cap\calQ_{\t,(\rho,\ovrho)}$; by non-arithmeticity, the cone $\calL_\cc$ has non-empty interior and thus Corollary~\ref{strictly} gives that $\calQ_\cc$ is a strictly convex curve. We consider the max norm $\|v\|_{\infty,\hol}=\max\{\hol|\slroot(v)|,|\ov \slroot(v)|\}$ on $V$, and its dual (operator) norm on $V^*$ denoted by $\|\,\|^{1,\hol}$. Let $\vi\in\calQ_\cc$ be the unique form such that
\[
\left\|\vi\right\|^{1,\hol}=\inf\left\{\|\varphi\|^{1,\hol}\,:\,\varphi\in\calQ_\cc\right\}.
\]



In the following lemma the role of the assumptions on dynamical intersection in Theorem~\ref{Hffconical} becomes clear:

\begin{lemm}\label{paraphi}
The functional $\vi/\|\vi\|^{1,\hol}$ is a convex combination $s\hol\slroot+(1-s)\ov\tau$ with $s\in(0,1)$ if and only if
\begin{equation}\label{e.Iassumption}
\II_\slroot(\ovsloot)> \hol>1/\II_{\ovsloot}(\slroot).
\end{equation}
In this case one has $\sf T_{\vi}\calQ_\cc=\spa\{ \hol\slroot-\ov{\slroot}\}$.
\end{lemm}


\begin{proof}
Recall from Corollary~\ref{strictly} that $\sf T_{\hJ\slroot\slroot}\calQ_\cc=\ker\II_{\hJ\slroot\slroot}$ and $\calQ_\cc$ is strictly convex.
Furthermore, by definition the functional $\vi$ is the point of $\calQ_\cc$, that minimizes the norm $\|\; \|^{1,\hol}$. The level set $\{\|\varphi \|^{1,\hol}=1\}$ is a rhombus with vertices $(\hol\slroot,\ovsloot)$ (in blue in Figure~\ref{phinfty}), the tangent to $\calQ_\cc$ at $\hJ\slroot\slroot$, in red in Figure~\ref{phinfty}, is the level set $\II_{\hJ\slroot\slroot}(\cdot)=1$, whence its intersection with the $\ovsloot$-axis is $\ovsloot/\II_{\hJ\slroot\slroot}(\ovsloot)$, and the tangent to $\calQ_\cc$ at $\hJ{\ovsloot}\ovsloot$ is the level set $\II_{\hJ{\ovsloot}\ovsloot}(\cdot)=1$, and it intersects the $\slroot$-axis is $\slroot/\II_{\hJ{\ovsloot}\ovsloot}(\slroot)$.

Eq.~\eqref{e.Iassumption} is thus equivalent to the fact that the slope of the side of the rhombus, equal to $-1/\hol$, is between the slope of the tangent at $\hJ\slroot\slroot$, which is equal to $-\hJ\slroot/\II_{\hJ\slroot\slroot}(\ovsloot)=-1/\II_{\slroot}(\ovsloot)$, and the slope of the tangent at $\hJ{\ovsloot}\ovsloot$, which is equal to $-\II_{\hJ{\ovsloot}\ovsloot}(\slroot)/\hJ{\ovsloot}=-\II_{\ovsloot}(\slroot)$.

Strict convexity of $\calQ_\cc$ ensures that this is equivalent to having a unique point in $\calQ_\cc\cap\{t\ovsloot:t>0\}\times\{s\slroot:s>0\}$ tangent to the side of the rhombus, which is the desired functional $\vi$.
\end{proof}

\begin{figure}[!h]
\centering
\begin{tikzpicture}

\begin{scope}[scale=0.9, rotate=5]

\begin{scope}[rotate=-30,scale=0.5]

\begin{scope}[shift={(0,3.5)},scale=0.85]
\draw[name path=Q, thick] (-13,3).. controls (0,-2.5).. (13,3);
\node at (13,3) [right] {$\calQ_\cc$};
\end{scope}
\end{scope}


\draw [name path=slroot2] (0,0) -- (120:6cm);
\draw (0,0) -- (180:2cm);
\draw (0,0) -- (300:2cm);
\draw [name path=slroot1] (0,0) -- (360:6cm);
\node at (5,.5) [above] {$V^*$};

\fill [name intersections={of=Q and slroot2, by={a}}] (a) circle (2pt);
\node [above right] at (a) {$\hJ{\ov\tau}\ov\tau$};

\fill [name intersections={of=Q and slroot1, by={b}}] (b) circle (2pt);
\node [below] at (b) {$\hJ\slroot\slroot$};

\coordinate (s) at (0.5,0);

\node [below] at (s) {$ \hol\slroot$};

\coordinate (tg) at (-1.2,0.9);

\draw [name path=tangente, red] (b) -- (tg);
\fill [name intersections={of=tangente and slroot2, by={I}}] (I) circle (2pt);

\coordinate (hpsi) at (-1.6,2.5);

\draw [name path=ig, white] (b) -- (hpsi);
\fill [name intersections={of=ig and slroot2, by={hs}}] (hs) circle (2pt);
\node [left] at (hs) {$\frac{\hJ\slroot}{\hol}\ov\tau$};

\coordinate (c) at ($(hs)-(b)$);

\draw [name path=bar, white] (s)--($(s)+0.5*(c)$);
\fill [name intersections={of=bar and slroot2, by={bs}}] (bs) circle (2pt);
\node [left] at (bs) {$\ov\tau$};

\draw [blue] (s) --(bs)--($(0,0)-(s)$)--($(0,0)-(bs)$)--(s);

\draw [name path=normh, blue] (b) --(hs);
\fill (s) circle (2pt);
\fill [name intersections={of=normh and Q, by={barpsi}}] (barpsi) circle (2pt);
\node [above right] at (barpsi) {$\ov{\psi}$};
\node [above left] (def) at (-3.5,0.5) {$\frac{\hJ\slroot}{\II_\slroot(\ov\tau)}\ov\tau$};
\node [left] (I') at (I) {};
\draw[->, orange] (def).. controls (0,1) and (-3.5,-.5).. (I');

\coordinate (phi) at (3.95,-0.2);

\draw [name path=phinfty, blue] (phi)--($(phi)+0.85*(c)$);

\fill [name intersections={of=phinfty and Q, by={varphi}}] (varphi) circle (2pt);

\node [above] at (varphi) {$\vi$};
\draw [gray] (0,0) -- (varphi);
\fill (varphi) circle (2pt);

\node [above] (tt) at (tg) {};

\node [above] at (-4,3.4) {$\hJ\slroot\slroot+\sf T_{\hJ\slroot\slroot}\calQ_\cc$};
\draw[->, red] (-4,3.4).. controls (-2,2) and (-3.5,0).. (tt);

\fill (hs) circle (2pt);
\fill (b) circle (2pt);
\fill (bs) circle (2pt);
\end{scope}
\end{tikzpicture}
\caption{The situation of Lemma~\ref{paraphi}.} \label{phinfty}
\end{figure}


We thus obtain the following key properties of $\vi$:

\begin{lemm}\label{phi>min}
Under the assumptions of Theorem~\ref{Hffconical} one has
\begin{enumerate}\romanenumi
\item\label{lemm5.10.1} $\sf u_\vi^\cc=\Pi(\ker(\hol\slroot-\ovsloot))$;
\item\label{lemm5.10.2} for any $v\in V^+$ one has
\[
\vi(v)\geq \hC{\infty,\hol}\hol\min\{\slroot(v),\ov{\slroot}(v)\}.
\]
\end{enumerate}
Moreover one has $\hC{\infty,\hol}<\min\{\hJ{\ovsloot},\hJ{\slroot}/\hol\}$.
\end{lemm}



\begin{proof}
Lemma~\ref{paraphi} implies that
\begin{itemize}
\item[\eqref{lemm5.10.1}] $\sf T_{\vi}\calQ_\cc=\spa\{\hol\slroot-\ov{\slroot}\}$ and thus
\[
\sf u_\vi^\cc=\Ann\bigl(\R\cdot(\hol\slroot-\ov{\slroot})\bigr)=\Pi(\ker(\hol\slroot-\ovsloot)).
\]
\item[\eqref{lemm5.10.2}] $\vi/\|\vi\|^{1,\hol}=s\hol\slroot+(1-s)\ov{\slroot}$ for some $s\in(0,1)$ and hence\footnote{Indeed, if $x,y\geq0$, $s\in(0,1)$ and $\hol\in(0,1]$ one has: $s\hol x+(1-s)y\geq\hol\min(x,y)$. Assume for example that $y\geq x$ (the other case follows similarly), then
\[
s\hol x+(1-s)y-\hol x\geq (1-s)(1-\hol)x\geq0.
\]
}, since $\hol\in(0,1]$,
\[
\vi(v)\geq\|\vi\|^{1,\hol}\hol\min\{\slroot(v),\ov{\slroot}(v)\}
\]
for all $v\in V^+$.
\end{itemize}

In order to prove item~\eqref{lemm5.10.2}, we need to show that $ \hC{\infty,\hol}\leq\|\vi\|^{1,\hol}$. Since 
\[
\vi\bigl(\cartan_\t((\rho,\ovrho)\g)\bigr)\leq \bigl\|\Pi(\cartan_\t((\rho,\ovrho)\g))\bigr\|_{\infty,\hol}\left\|\vi\right\|^{1,\hol},
\]
we deduce, for all $s>\|\vi\|^{1,\hol}$,
\[
\sum_{\g\in\G}e^{-s\left\|\Pi(\cartan_\t((\rho,\ovrho)\g))\right\|_{\infty,\hol}}\leq \sum_{\g\in\G}e^{-\left(s/\left\|\vi\right\|^{1,\hol}\right)\vi\left(\cartan_\t((\rho,\ovrho)\g)\right)}<\infty
\]
where last inequality holds as $\hC\vi=1$ (by Eq.~\eqref{qh1} and Remark~\ref{r.entropy=crit}).


The last assertion follows directly from the definitions:
\begin{align*}
\hC{\infty,\hol}
&=\lim_{t\to\infty}\frac 1t\log\#\left\{\g\in\G:\max\bigl\{\hol\slroot(\cartan(\g)),\ovsloot(\cartan(\ov\g))\bigr\}\leq t\right\}\\
&<\lim_{t\to\infty}\frac 1t\log\#\left\{\g\in\G:\bigl\{\hol\slroot(\cartan(\g))\big\}\leq t\right\}\\
&=\hC{\slroot}/\hol=\hJ{\slroot}/\hol,
\end{align*}
where the strict inequality comes from Corollary~\ref{quintgrowththeta} together with strict convexity of $\calQ_{\t,(\rho,\ovrho)}$ and the last equality follows from Remark~\ref{r.entropy=crit}. The inequality $\hC{\infty,\hol}\leq \hJ{\ovsloot}$ is analogous.
\end{proof}



Let $\ps^\vi$ be the Patterson--Sullivan measure associated to $\vi$ by Corollary~\ref{existe}. Combining Eqs.~\eqref{e.CartBasin}, \eqref{e.conesinCartan} and Corollary~\ref{existe} we deduce that, for every $\g\in\G$,
\begin{equation}\label{shadow}
\ps^\vi\left(\g\cone_\infty^{c}(\g)\right)\leq Ce^{-\vi\left(\cartan_\t((\rho,\ovrho)\g)\right)}\leq Ce^{-\hC{\infty,\hol}\hol\min\left\{\slroot(\cartan(\g)),\ov{\slroot}(\cartan(\ov\g))\right\}},
\end{equation}
where the last inequality comes from Lemma~\ref{phi>min}.


By Proposition~\ref{conetypesBalls} there exist constants $ c,k_1$ and $\ov k_1$ such that if $(\alpha_i)_0^\infty$ is a geodesic ray from $\id$ to $x$ then for all $i$ the subsets
\[
\xi\left(\alpha_i\cone_\infty^{c}(\alpha_i)\right)\quad\text{and}\quad\ov{\xi}\left(\alpha_i\cone_\infty^{c}(\alpha_i)
\right)
\]
contain balls on the corresponding projective spaces of radii
\[
k_1e^{-\slroot(\cartan(\alpha_i))} \quad\text{and}\quad \ov k_1e^{-\ov{\slroot}(\cartan(\ov\alpha_i))}
\]
respectively where $k_1, \ov k_1$ depend on the representations but not on $i$. Since $\varXi(\bord\G)$ is a graph, the preceding radius computation implies that the image of the cone type $\varXi(\alpha_i\cone_\infty^{c}(\alpha_i))$ contains the intersection of $\xi(\bord\G)\times\ov{\xi}(\bord\G)$ with a ball, for the product metric on $\P(\K^d)\times\P(\K^{\ovd})$, of radius
\begin{equation}\label{radiusgraph}
ke^{-\min\left\{\slroot(\cartan(\alpha_i)),\ov{\slroot}(\cartan(\ov\alpha_i))\right\}},
\end{equation}
for some uniform constant $k$. This set of balls forms a fine set of neighbourhoods around any point $x\in\partial\G$. Combining this with Eq.~\eqref{shadow} and the fact that $\mu^\vi$ is supported on $\bord\G$, one has, possibly enlarging the constant $C$, that for all $r$ the measure of the ball of radius $r$ about $\varXi(x)$ is
\[
\ps^\vi\left(B(x,r)\right)\leq Cr^{-\hC{\infty,\hol}\hol}.
\]

Since $\dim V^*=2$ and $\cc_{(\rho,\ovrho)}$ is assumed non-arithmetic, Theorem~\ref{ps-sandwich} states that the subset of $\hol$-conical points has full $\ps^\vi$ measure. Applying Corollary~\ref{edgar} one concludes that
\[
\Hff\left(\varXi\{\hol-\text{conical points}\}\right)\geq\hol\hC{\infty,\hol}.
\]


\subsection{The upper bound} We now prove the second inequality.

\begin{prop}\label{upper}
Let $\rho,\ovrho$ be locally conformal representations over $\K$ and~$\bar{\K}$. For every $\hol\leq 1$,
\[
\Hff\bigl(\varXi\{\hol-\text{conical\ points}\}\bigr)\leq\min\left\{\hC{\infty,\hol},\hol\hC{\infty,\hol}+(1-\hol) \right\}.
\]
\end{prop}

\begin{proof}
We say that a point $x$ is \emph{$(R,\hol)$-conical} if there exists a geodesic ray $(\alpha_i)_{i\in\N}$ converging to $x$ such that for an infinite subset $\I\subset \N$ of indices and for every $k\in \I$
\begin{equation}\label{e.5.11}
\bigl|\hol\slroot\left(\cartan(\alpha_k)\right)-\ov{\slroot}\left(\cartan(\ov\alpha_k)\right)\bigr|\leq R.
\end{equation}
We denote by $\sf C_\hol^R$ the set of $(R,\hol)$-conical points. By Definition~\ref{flat} one has
\[
{\displaystyle\bigcup_{R>0}\sf C_\hol^R=\{x\in\bord\G: x\text{ is }\hol-\text{conical}\}},
\]
and thus by Lemma~\ref{l.Hffunion} it suffices to show that for every $R$ one has
\[
\Hff\left(\sf C_\hol^{R}\right)\leq \hC{\infty,\hol}.
\]

For any constant $K>0$ and any $\g\in\G$ we denote by $B_{\g}^{\max, K}$ the open ball of $\P(\K^{d})\times \P(\K^{\ovd})$ given by:
\[
B_{\g}^{\max, K}:= B\left(\left(U_1(\g),U_1(\ov\g)\right), K e^{-\max{\{\hol\slroot(\cartan(\g)), \ov{\slroot}(\cartan(\ov\g))\}}}\right),
\]
and denote by
\[
\cal U_T^K:=\big\{B_{\g}^{\max, K}|\, |\g|\geq T\big\}.
\]
Let $K$, resp. $\ov K$, be the constants given by Proposition~\ref{p.coneinBall} for the representation $\rho$ (resp. $\ovrho$).

We first observe that for $C=2e^R\max\{K,\ov K\}$ and every $T>0$, the set $\cal U_T^C$ covers $\varXi(\sf C_\hol ^R)$. Indeed, if $x\in\sf C_\hol^R$ consider the geodesic ray $(\alpha_i)_{i\in\N}$ converging to $x$, and the set $\I$ of indices for which Eq.~\eqref{e.5.11} holds. Then for every $k\in\I$ one has, since $\hol\leq1$, that
\begin{align}
\slroot\left(\cartan(\rho\alpha_k)\right) &\geq\hol\slroot\left(\cartan(\alpha_k)\right)>\max\left\{\hol\slroot\left(\cartan(\alpha_k)\right), \ovsloot\left(\cartan(\ov\alpha_k)\right)\right\}-R, \label{e.5.12.a}\\
\ovsloot\left(\cartan(\ov\alpha_k)\right)  &>\max\left\{\hol\slroot\left(\cartan(\alpha_k)\right), \ovsloot\left(\cartan(\ov\alpha_k)\right)\right\}-R. \label{e.5.12}
\end{align}
Let now $T$ be fixed and choose $k\in \I$, $k>T$. Since $x\in\alpha_k\cone_\infty^{c}(\alpha_k)$, Proposition~\ref{p.coneinBall} together with Eq.~\eqref{e.5.12} give
\begin{align*}
d\left(\xi(x),U_1(\alpha_k)\right) &\leq Ce^{-\max\left\{\hol\slroot\left(\cartan(\alpha_k)\right),\, \ovsloot\left(\cartan(\ov\alpha_k)\right)\right\}}\\
d\left(\ovxi(x),U_1(\ov\alpha_k)\right) &\leq Ce^{-\max\left\{\hol\slroot\left(\cartan(\alpha_k)\right),\, \ovsloot\left(\cartan(\ov\alpha_k)\right)\right\}},
\end{align*} as desired.

Furthermore, by definition of $\hC{\infty,\hol}$, for every $s>\hC{\infty,\hol}$,
\[
\sum_{U\in \cal U_T^C}\diam U^s\leq2^sC^s\sum_{|\g|\geq T}e^{-s\max\left \{\hol\slroot(\cartan(\g)),\,\ovsloot(\cartan(\ov\g))\right\}}<+\infty,
\]
whence, Eq.~\eqref{HffDef} yields $\Hff(\sf C_\hol^R)\leq\hC{\infty,\hol}$. In order to obtain the second upper bound we observe that, if $\alpha\in\G$ satisfies Eq.~\eqref{e.5.11}, the set $\varXi(\alpha\cone_\infty^{c}(\alpha))$ can be covered with $e^{(1-\hol)\slroot(\cartan(\alpha))}$ balls of radius $2Ce^{-\slroot(\cartan(\alpha))}$. We denote by $\cal U_T$ the collection of open balls, that only take into account elements $\alpha\in\G$ with $|\alpha|>T$ that verify~\eqref{e.5.11}, which in particular covers the set $\sf C_\hol^R$. Using Eq.~\eqref{e.5.12.a} we obtain
\begin{align*}
\sum_{U\in \cal U_T}\diam U^s &\leq 2^sC^s\sum_{|\g|\geq T}e^{(1-\hol)\slroot(\cartan(\g))}e^{-s\slroot(\cartan(\g))}\\
&\leq 2^sC^s\sum_{|\g|\geq T}e^{-(s-(1-\hol))\slroot(\cartan(\g))}\\
&\leq 2^sC^se^{\frac{R(s-1+\hol)}\hol}\sum_{|\g|\geq T}e^{-\frac{(s-(1-\hol))}\hol\max\left\{\hol\slroot(\cartan(\g)),\,\ovsloot(\cartan(\ov\g))\right\}}.
\end{align*}

Since the latter quantity is finite whenever $\frac{(s-(1-\hol))}\hol>\hC{\infty,\hol}$, we deduce
\[
\Hff\left(\sf C_\hol^R\right)<\hol\hC{\infty,\hol}+(1-\hol).
\]
\end{proof}



We conclude this subsection computing the Hausdorff dimension of the image of the whole boundary through the graph map. See~\cite{HffGraph} for examples of homeomorphisms between Cantor sets for which the Hausdorff dimension of the graph exceeds the maximal Hausdorff dimension of the factors.

\begin{prop}\label{p.uppergeneral}
Let $\rho:\G\to\SL(d,\K)$, $\ovrho:\G\to\SL(\ovd,\bar{\K})$ be locally conformal. Then
\[
\Hff(\varXi(\partial\G))=\max\left\{\hC{\tau},\hC{\ov \tau}\right\}
\]
\end{prop}
\begin{proof}
This follows as in the proof of Proposition~\ref{upper} considering the covers of $\varXi(\partial\Gamma)$ given by $\cal U_T^C:=\{B_{\g}^{\min, C}\mid |\g|\geq T\}$ with
\[
B_{\g}^{\min, K}:= B\left(\left(U_1(\g),U_1(\ov\g)\right),\; K e^{-\min{\left\{\slroot(\cartan(\g)),\, \ov{\slroot}(\cartan(\ov\g))\right\}}}\right),
\]
and $C=2\max\{K, \ov K\}$ where $K$ (resp. $\ov K$) is the constant given by Proposition~\ref{p.coneinBall} for the representations $\rho$ (resp. $\ov \rho$). To conclude it is enough to observe that
\[
\hC{\min \{\slroot, \ovsloot\}}=\max\left\{\hC{\slroot}, \hC{\ov \slroot}\right\},
\]
a fact proven for example in~\cite[Lemma~5.1]{PSW2}.
\end{proof}

It is easy to generalize Proposition~\ref{p.uppergeneral} to an arbitrary number of factors. as an application we get.
\begin{coro}
Let $\rho:\G\to\SL(d,\K)$ and $\theta\subset\Delta$ be such that for all $\slroot_i\in\theta$, $\Phi_{\slroot_i}\circ\rho$ is $(1,1,2)$-hyperconvex. Then
\[
\Hff\left(\xi^\theta_\rho(\partial\G)\right)=\max_{\slroot_i\in\theta}\hC{\slroot_i}.
\]
\end{coro}

\subsection{Proof of Theorem~\ref{Hffconical}}\label{proofHffconical}
The first inequality is established in Section~\ref{s.lower}, the second inequality is proven in Proposition~\ref{upper}, the third inequality follows from Lemma~\ref{phi>min} and the fourth from Theorem~\ref{hyperh=1}.The last equality was established in Proposition~\ref{p.uppergeneral}.\qed


\section{\texorpdfstring{$\hol$}{hol}-concavity and \texorpdfstring{$\hol$}{hol}-conicality: Final steps for the proof of Theorem~\ref{LCdiff}}\label{teoLCdiff}



The goal of this section is to prove the following more general version of Theorem~\ref{LCdiff}. As before, fix $\{\K,\bar{\K}\}\subset\{\R,\C,\H\}$ together with locally conformal representations $\rho:\G\to\SL(d,\K)$ and $\ovrho:\G\to\SL(\ovd,\bar{\K})$ of an arbitrary word-hyperbolic group~$\G$. For $\hol\in(0,1]$ recall that $\Xi:\xi(\bord\G)\to\ovxi(\bord\G)$ is \emph{$\hol$-concave} at $x\in\bord\G$ if there exists $y_k\to x$ such that the incremental quotients
\begin{equation}\label{incrlimit}
\frac{d_{\P}\left(\ovxi(x),\ovxi(y_k)\right)}{d_{\P}\left(\xi(x),\xi(y_k)\right)^\hol}
\end{equation}
are bounded away from $0$ and $\infty$ (independently of $k$). We also let $\Ext_{\rho,\ovrho}^\hol$ be the set of $x\in\bord\G$ that are $\hol$-concavity points of $\Xi$. Finally, recall that $\rho$ and $\ovrho$ are \emph{not gap-isospectral} if there exists $\g\in\G$ such that $\slroot(\cartan(\g))\neq\ovsloot(\cartan(\ov\g))$.



\begin{theo}\label{tutti.LCdiff}
Let $\rho, \ovrho$ be locally conformal representations acting irreducibly on $\K^d$ and ${\bar{\K}}{}^{\ovd}$ respectively as real vector spaces, that are not gap-isospectral. Consider any $\hol\in(0,1]$ with $\II_{\slroot}(\ovsloot)>\hol>(\II_{\ovsloot}(\slroot))^{-1}$, then:
\begin{itemize}
\item if $\{\K,\bar{\K}\}\subset\{\R,\C\}$ one has
\begin{equation}\label{holdineq}
\begin{aligned}[t]
\hol\hC{\infty,\hol}\leq\Hff\left(\varXi({\Ext}_{\rho,\ovrho}^\hol)\right)
&\leq\min\left\{\hC{\infty,\hol},\hol\hC{\infty,\hol}+(1-\hol) \right\}\\
&<\min\left\{\hJ{\ovsloot},\hJ{\slroot}/\hol\right\}\\
&\leq\Hff(\varXi(\bord\G))\\
&=\max\{\hJ{\slroot},\hJ{\ovsloot}\};
\end{aligned}
\end{equation}
\item if $\K=\H$ (resp. $\bar{\K}=\H$), Eq.~\eqref{holdineq} holds if we further assume that the real Zariski closure of $\rho(\G)$ (resp. of $\ovrho(\G)$) does not have compact factors.
\end{itemize}
\end{theo}



\subsection{Hyperplane conicality and the concavity condition}



We commence with a lemma relating $\hol$-conicality to the desired concavity properties of the equivariant map $\Xi:\xi(\bord\G)\to\ovxi(\bord\G)$.


\begin{lemm}\label{generalcase}
Let $\rho$ and $\ovrho$ be locally conformal representations over $\K$ and $\bar{\K}$ respectively, and $\hol\in(0,1]$. Then one has $\{\hol\text{-conical points of }
(\rho,\ovrho)\}=\Ext_{\rho,\ovrho}^\hol$.
\end{lemm}



\begin{proof}
Let $(\alpha_i)_{i\in\N}$ denote a geodesic ray converging to $x$. Proposition~\ref{conetypesBalls} gives constants $C_1,C_2, \bar{C}_1, \bar{C}_2$ and $L\in\N$ such that, for every $n\in\N$ and every $y_n\in \alpha_n\cone_{\infty}^{c}(\alpha_n)\setminus\alpha_{n+L}\cone_\infty^{c}(\alpha_{n+L})$, it holds
\begin{equation}\label{e.nestcone}
\begin{aligned}
C_1e^{-\slroot(\cartan(\alpha_n))} &<d_\P\bigl(\xi(y_n),\xi(x)\bigr) <C_2e^{-\slroot(\cartan(\alpha_n))},  \\
\bar{C}_1e^{-\ov{\slroot}(\cartan(\ov\alpha_n))} &<d_\P\bigl(\ovxi(y_n),\ovxi(x)\bigr) <\bar{C}_2e^{-\ovsloot(\cartan(\ov\alpha_n))}.
\end{aligned}
\end{equation}



Assume first that $x$ is $\hol$-conical. By Definition~\ref{flat} we obtain a geodesic ray $(\alpha_i)_0^\infty$, an infinite set of indices $\I\subset\N$ and a number $R$, such that for all $k\in\I$ one has
\begin{equation}\label{e.conical}
\bigl|\hol\slroot(\cartan(\alpha_k))-\ovsloot(\cartan(\ov\alpha_k))\bigr|<R.
\end{equation}

For each such $k$ we choose a point $y_k\in \alpha_k\cone_\infty^{c}(\alpha_k)\setminus \alpha_{k+L}\cone_\infty^{c}(\alpha_{k+L})$. By construction $y_k$ converges to $x$. Combining both equations, for every $k\in \I$ it holds
\[
e^{-R}\frac{\bar{C}_1}{ {C_2}^\hol}\leq \frac{d_\P\left(\ovxi(y_k),\ovxi(x)\right)}{d_\P\left(\xi(y_k),\xi(x)\right)^\hol} \leq e^{R}\frac{\bar{C}_2}{ {C_1}^\hol},
\]
so the incremental quotient~\eqref{incrlimit} is uniformly far from $0$ and $\infty$. Whence $\{\hol-\text{conical points}\}\subset\Ext_{\rho,\ovrho}^\hol$.



Conversely, assume that $x$ is not $\hol$-conical. The Cartan projections of two consecutive elements $\alpha_i$ and $\alpha_{i+1}$ make uniformly bounded gaps (Proposition~\ref{p.Ben}), and thus there exists $C$ such that for all $n\in\N$ one has
\[
\bigl|\slroot(\cartan(\alpha_{n+1}))-\slroot(\cartan(\alpha_{n}))\bigr|<C.
\]
As a consequence, we can assume, up to switching the roles of $\rho$ and $\ovrho$, that for any $R$ there exists $n_R$ such that for every $n> n_R$ one has
\[
-\hol\slroot(\cartan(\alpha_{n}))+\ov{\slroot}(\cartan(\ov\alpha_{n}))>R.
\]
In turn this implies, thanks to Eq.~\eqref{e.nestcone}, that for every $y\in \alpha_{n_R}\cone_\infty^{c}(\alpha_{n_R})$,
\[
\frac{d_\P\left(\ovxi(y),\ovxi(x)\right)}{d_\P\left(\xi(y),\xi(x)\right)^\hol}\leq e^{-R}\frac {\bar{C}_2}{ {C_1}^\hol}.
\]

Since $R$ is arbitrary, and the sets $\alpha_{n_R}\cone_\infty^{c}(\alpha_{n_R})$ form a system of neighborhoods of the point $x$, we deduce that the limit in Eq.~\eqref{incrlimit} exists and equals $0$. This concludes the proof of Lemma~\ref{generalcase}.
\end{proof}



\subsection{Non-arithmeticity of periods}



In this section we establish a non-arithme\-ti\-ci\-ty condition, necessary to apply later Theorem~\ref{Hffconical}. This is established in a rather general setting. Recall that a subgroup $\grupo<\SL(d,\K)$ is \emph{$\K$-proximal} if it contains a $\K$-proximal element, i.e. there exists $g\in\grupo$ such that $\slroot_1(\lambda(g))>0$.


\begin{prop}\label{nonA}
Let $\grupo$ be a finitely generated group. Let $\rho:\grupo\to\SL(d,\K)$ and $\ovrho:\grupo\to\SL(\ovd,\bar{\K})$ be two $\K$-proximal representations that act irreducibly on $\K^d$ and $\bar{\K}{}^{\ovd}$ respectively, as real vector spaces. Assume there exists $\g\in\grupo$ proximal for $\rho$ and $\ovrho$ such that $\slroot_1(\lambda(\rho \g))\neq\ovsloot_1(\lambda(\ov \rho \g))$. If $\{\K,\bar{\K}\}\subset\{\R,\C\}$, then the group generated by the pairs
\[
\Bigl\{\bigl(\slroot_1(\lambda(\rho \g)),\ovsloot_1(\lambda(\ov \rho \g))\bigr)\,:\,\g\in\grupo\Bigr\}
\]
is dense in $\R^2$. If $\K=\H$ we further assume that the Zariski closure over $\R$ of $\rho(\grupo)$ has no compact factors, and the same for $\ovrho(\grupo)$ if moreover $\bar{\K}=\H$, then the same conclusion holds.
\end{prop}

To prove the proposition we need Lemmas~\ref{semisimple} and~\ref{semi} below.


\begin{lemm}\label{semisimple}
Let $\K$ be either $\R$ or $\C$. Let $\grupo<\SL(d,\K)$ be a subgroup acting irreducibly on $\K^d$ as a real vector space and assume $\grupo$ contains a $\K$-proximal element. Then the real Zariski closure of $\grupo$ is semi-simple, has finite center and without compact factors.
\end{lemm}



\begin{proof}
If $\K=\R$ Lemma~\ref{semisimple} is the content of~\cite[Lemma~8.6]{entropia} and the proof over $\C$ is a slight modification of the latter. Indeed, let $\sfG$ be the Zariski closure of $\rho(\grupo)$ over the reals, by the irreducibility assumption it is a reductive (real-algebraic) group. By Schur's lemma the elements commuting with $\grupo$ consist only on homotheties, but since we're in special linear group one has that the center of $\sfG$ is finite.


The group $\sfG$ is then semi-simple and we let $K$ be the identity component of the product of all the compact simple factors of $\sfG$. We also let $H$ be the identity component of the product of all the non-compact simple factors of $\sfG$. The groups $H$ and $K$ commute and one has $HK$ has finite index in $\sfG$.

Consider a proximal $g\in\sfG$, up to a fixed power we may write $g=kh$ with $k\in K$ and $h\in H$. Since $K$ is compact, its eigenvalues have modulus one so we conclude that $h$ is proximal and that $g_+=h_+$. The attracting line of $h$ is thus invariant under~$K$. Since $K$ is connected, an element of $ K$ acts on $h_+$ as multiplication by some element of $\circle$.

By irreducibility we may find a basis of $\C^d$ consisting on fixed attracting lines of proximal elements of $H$. This basis simultaneously diagonalizes $K$, so we get an injective map from $K$ to a compact group isomorphic to a $d$-dimensional torus. Consequently $K$ is abelian, and since it commutes with $H$ we conclude that $K$ is contained in the identity component of the center of $\sfG$, which we proved earlier to be trivial.
\end{proof}



\begin{lemm}\label{semi}
Let $G$ be a semi-simple real-algebraic Lie group with finite center and no compact factors. Fix $\vartheta,\ov\vartheta\subset\simple_G$ two non-empty subsets with $\vartheta\cap\ov\vartheta=\emptyset$. Let $\grupo$ be a group and $\scr r:\grupo\to G$ a representation with Zariski-dense image. Then, for every closed cone with non-empty interior $\scr C\subset\inte\calL_{\scr r(\grupo)}$, the group spanned by the pairs
\[
\left\{\left(\min_{\sigma\in\vartheta}\sigma(\lambda(\scr r g)),\min_{\ov\sigma\in\ov\vartheta}\ov\sigma(\lambda(\scr r g))\right)\,:\,g\in\grupo\quad\text{and}\quad\lambda(\scr rg)\in\scr C\right\}
\]
is dense in $\R^2$.
\end{lemm}

\begin{proof}
Define the piecewise linear maps $\slroot,\ovsloot:\a^+\to\R$ by:
\begin{align*}
\slroot(v) & =\min\bigl\{\sigma(v):\sigma\in\vartheta\bigr\}\\
\ovsloot(v) & =\min\bigl\{\ov\sigma(v):\ov\sigma\in\ov\vartheta\bigr\}.
\end{align*}

The vanishing set of the difference $\slroot-\ovsloot$ is contained the union of $\ker(\sroot-\bb)$ for arbitrary $\sroot\in\vartheta$ and $\bb\in\ov\vartheta$. Since $\vartheta$ and $\ov\vartheta$ are disjoint, this is a union of hyperplanes of $\a$, from which we deduce that the set of zeroes of $\slroot-\ovsloot$ has empty interior.



Since $\scr C\subset \inte\calL_{\scr r(\grupo)}$ has non-empty interior, the difference $\slroot-\ovsloot$ does not identically vanish on $\scr C$. Since $\slroot$ and $\ovsloot$ are piecewise linear, we can choose a possibly smaller closed cone with non-empty interior
\[
\scr C'\subset\scr C,
\]
and $\sroot\in\vartheta$, $\bar{\bb}\in\ov\vartheta$ such that for all $v\in\scr C'$ one has
\[
\slroot\times\ovsloot(v):=\left(\slroot(v),\ovsloot(v)\right)=\left(\sroot(v),\bar{\bb}(v)\right).
\]
Since $\sroot$ and $\bb$ are distinct simple roots the map $ (\sroot,\bb):\a\to\R^2$ is surjective.


By~\cite[Proposition~5.1]{limite} there exists a sub-semigroup $\grupo'<\grupo$ such that $\scr r(\grupo')$ is a Zariski-dense Schottky semi-group with $\calL_{\scr r(\grupo')}=\scr C'$. In particular, for all $\g\in\grupo'$ one has
\[
\slroot\times\ovsloot(\lambda(\scr r\g))=\left(\sroot(\lambda(\scr r\g)),\bar{\bb}(\lambda(\scr r\g))\right).
\]
By Benoist's Theorem~\ref{densidad}, stating that the group generated by the Jordan projections $\lambda(\scr r\g)$, for $\g\in\grupo'$, is dense in $\a$, we conclude that the group spanned by
\[
\left\{\left(\left(\sroot(\lambda(\scr r\g)),\;\bar{\bb}(\lambda(\scr r\g))\right)\right):\g\in\grupo'\right\}
\]
is dense in $\R^2$, giving in turn the desired conclusion.
\end{proof}



\begin{proof}[Proof of Proposition~\ref{nonA}]
Denote by $\sfG$ and $\bar{\sfG}$ the Zariski closures of $\rho(\grupo)$ and $\ovrho(\grupo)$ respectively. Both $\sfG$ and $\bar{\sfG}$ are semi-simple, have finite center, and don't have compact factors: if $\{\K,\bar{\K}\}\subset\{\R,\C\}$ then this is the content of Lemma~\ref{semisimple}, if either $\K$ and/or $\bar{\K}$ equals $\H$ then this is an assumption. We let $\iota:\grupo\to\sfG$ and $\ov\iota:\grupo\to\bar{\sfG}$ be the respective inclusions.

If we let $\phi:\sfG\to\SL(d,\K)$ and $\ov\phi:\bar{\sfG}\to\SL(\ovd,\bar{\K})$ be the associated real representations, so that $\rho=\phi\circ\iota$ and $\ovrho=\ov\phi\circ\ov\iota$, we have from Section~\ref{representaciones} two subsets of simple roots $\t:=\t_\phi$ and $\ov\t:=\t_{\ov\phi}$ such that for all $a\in\a_\sfG^+$ and $b\in\a_{\bar{\sfG}}^+$ one has
\begin{equation}\label{piecewise}
\begin{aligned}
\slroot(a)&:=\slroot_1(\phi(a))  =\min\bigl\{\sroot(a):\sroot\in\t\bigr\}\\ \ov\tau(b)&:=\ovsloot_1(\ov\phi (b))  =\min\bigl\{\ov\sroot(b):\ov\sroot\in\ov\t\bigr\}.
\end{aligned}
\end{equation}
In particular, for every $\g\in\grupo$ one has $\slroot_1(\lambda(\g))=\slroot(\lambda_\sfG(\iota\g))$, and similarly for $\ovrho$.


Since $\phi$ and $\ov\phi$ are faithful, $\t$ and $\ov\t$ contain at least one root of each factor of, respectively, $\sfG$ and $\bar{\sfG}$. If $\vartheta\subset\t$ then we let
\[
\slroot^{\vartheta}(v)=\min_{\sigma\in\vartheta}\sigma(v),\ v\in\a_{\sfG}.
\]


If $\sf H$ is a non-trivial product of simple factors of $\sfG$ then we let $\iota_\sf H:\grupo\to\sf H$ be the composition of $\iota$ with the projection of $\sfG$ onto $\sf H$. By Zariski-density of $\iota(\grupo)$, each representation $\iota_\sf H$ has Zariski-dense image (though unlikely to be discrete). We also~let
\[
\t_{\sf H} =\t\cap\simple_{\sf H}.
\]
Each $\t_\sf H$ is non-empty. We analogously define $\ov\iota_{\bar{\sf H}}$, $\t^{\bar{\sf H}}$ and $\ovsloot^{\bar{\sf H}}$.



We now let $\sf L$ be the largest product of simple factors, simultaneously of $\sfG$ and $\bar{\sfG}$, so that $\iota_\sf L$ is conjugated (up to finite index) to $\ov\iota_\sf L$. Let $\sf H$ and $\bar{\sf H}$ be the remaining factors of $\sfG$ and $\bar{\sfG}$ respectively, i.e.
\[
\sfG=\sf L\times \sf H\quad\text{and}\quad\bar{\sfG}=\sf L\times\bar{\sf H},
\]
and moreover, by definition of $\sf L$, the representation $\scr r:\grupo\to \sf L\times\bar{\sf H}\times\sf H$
\begin{equation}\label{repsZ1}
\scr r:g\mapsto \bigl(\iota_{\sf L}(g),\;\ov\iota_{\bar{\sf H}}(g),\;\iota_\sf H(g)\bigr)
\end{equation}
has Zariski-dense image, see for example~\cite[Corollary~11.6]{pressure}. %Bridgeman--Canary--Labourie-S.
We remark that we are not assuming that any of $\sf L$, $\bar{\sf H}$ or $\sf H$ is non-trivial (they can't, of course, be all trivial).

If $(u,v,w)\in\a_{\sf L}\times\a_{\bar{\sf H}}\times\a_{\sf H}$ we naturally think of $(u,v)$ as an element of $\a_{\bar{\sfG}}$ and of $(u,w)$ as an element of $\a_{\sfG}$. We now write
\begin{align*}
\Theta &=\t_\sf L\cap\ov\t_\sf L,\\ 
\Theta_{\sf L} &=\t_\sf L\setminus\Theta,\\
\ov\Theta_{\sf L} & =\ov\t_{\sf L}\setminus\Theta.
\end{align*}



One has, for all $(u,v,w)\in\a_{\sf L}\times\a_{\bar{\sf H}}\times\a_{\sf H}$ that
\begin{equation}\label{tauL}
\begin{aligned}
\slroot(u,w) &=\min\left\{\slroot^{\Theta_\sf L}(u),\slroot^{\Theta}(u),\slroot^{\t_\sf H}(w)\right\}\\
\ovsloot(u,v) &=\min\left\{\slroot^{\ov\Theta_\sf L}(u),\slroot^{\Theta}(u),\ovsloot^{\t_{\bar{\sf H}}}(v)\right\}.
\end{aligned}
\end{equation}



By assumption, there exists $g\in\grupo$ such that $\rho(g)$ and $\ovrho(g)$ are proximal and
$\slroot(\lambda_{\sfG}(\iota g))\neq\ovsloot(\lambda_{\bar{\sfG}}(\ov\iota g))$. Assume, without loss of generality, that
\begin{equation}\label{eqA}
\slroot(\lambda_{\sfG}(\iota g))<\ovsloot\left(\lambda_{\bar{\sfG}}(\ov\iota g)\right).
\end{equation}
By means of Eqs.~\eqref{tauL} we see that in this situation one has
\[
\slroot^{\Theta_{\sf L}\cup\t_\sf H}(\lambda_\sfG(\iota g))=\slroot(\lambda_\sfG(\iota g))<\ovsloot(\lambda_{\bar{\sfG}}(\ov\iota g)),
\]
in particular the union $\Theta_\sf L\cup\t_\sf H$ must be non-empty. Moreover, this strict inequality yields the existence of a small closed cone with non-empty interior $\scr C_0\subset\calL_\rho\subset\a_{\sfG}^+$ about $\R_+\lambda_\sfG(\rho g) $ such that
\begin{equation}\label{tauH=tau1}
\slroot^{{\Theta_{\sf L}\cup\t_\sf H}}(a)=\slroot_1(a)\,\forall a\in\scr C_0.
\end{equation}



Consider now the representation $\scr r:\grupo\to\sf L\times\bar{\sf H}\times\sf H$ from~\eqref{repsZ1} and a closed cone with non-empty interior $\scr C \subset\calL_{\scr r(\grupo)}\subset\a_{\sf L}^+\times\a_{\bar{\sf H}}^+\times\a_{\sf H}^+$ whose natural projection onto $\a_\sfG^+=\a_{\sf L}^+\times\a_{\sf H}^+$ is $\scr C_0$.

Lemma~\ref{semi} applied to the group $G=\sf L\times\bar{\sf H}\times\sf H$, the representation $\scr r$, the disjoint non-empty subsets $\vartheta=\Theta_\sf L\cup\t_\sf H$ and $\ov\vartheta=\ov\t_\sf L\cup\t_{\bar{\sf H}}$ and the cone $\scr C$, provides the desired conclusion.
\end{proof}


We conclude with the following corollary that we don't need but is of independent interest.

\begin{coro}\label{corraiz}
Let $\rho:\G\to\SL(d,\K)$ and $\ovrho:\G\to\SL(\ovd,\bar{\K})$ be $\R$-irreducible and $\{\slroot_1,\slroot_2\}$-Anosov and $\{\ovsloot_1,\ovsloot_2\}$-Anosov respectively. If $\K=\H$ assume moreover the Zariski closure of $\rho(\G)$ does not contain compact factors, and analogously for $\ovrho$. If $\rho$ and $\ovrho$ are not gap-isospectral then
\[
\II_{\ovsloot_1}(\slroot_1)>{\hJ{ \ovsloot_1}}/{\hJ{ \slroot_1}}.
\]
\end{coro}



\begin{proof}
Since both representations are projective-Anosov every element has proximal image in both representations. Pro\-po\-si\-tion~\ref{nonA} implies then that, since they are not gap-isospectral, the group spanned by the pairs $\{(\slroot_1(\lambda(\g)),\ovsloot_1(\lambda(\ov\g))):\g\in\G\}$ is dense.
Since $\slroot_1=2\peso_1-\peso_2$, one has $\slroot_1\in\a_{\{\slroot_1,\slroot_2\}}$ and thus we can compose the refraction cocycle of $\rho$ with $\tau_1$ to obtain a real valued cocycle with periods $\tau_1\circ\lambda$. The same holds for $\ovrho$ and so Proposition~\ref{ineq} applied to this pair of cocycles yields desired strict inequality.
\end{proof}



\subsection{Proof of Theorem~\ref{tutti.LCdiff}}\label{ProofA}

Theorem~\ref{tutti.LCdiff} follows from Proposition~\ref{nonA} giving the desired non-arithmeticity of periods, Lemma~\ref{generalcase} identifying the set $ \Ext_{\rho,\ovrho}^\hol$ with the set of $\hol$-conical points of $(\rho,\ovrho)$ and Theorem~\ref{Hffconical} computing the Hausdorff dimension of the latter when the periods are non-arithmetic. The last equality is a direct consequence of Proposition~\ref{p.uppergeneral}.\qed



\section{Theorem~\ref{t.Zcl}: Zariski closures of real-hyperconvex surface-group representations}\label{s.6}


In this section we prove Theorem~\ref{t.Zcl} giving a preliminary classification of Zariski closures of irreducible real $(1,1,2)$-hyperconvex representations of surface groups. For most of the section we work with a pair of $(1,1,2)$-hyperconvex representations and eventually reduce the proof of Theorem~\ref{t.Zcl} to a situation like this; we will crucially use Theorem~\ref{t.C1}.



\subsection{When \texorpdfstring{$\Xi$}{Xi} has oblique derivative}\label{s.6.1}
We prove here a result of independent interest, albeit possibly known to experts. This subsection only requires Section~\ref{cont} and Section~\ref{s.Anosov} and will be needed not only for Theorem~\ref{t.Zcl} but also for Theorems~\ref{tutti} and~\ref{thm.tuttiC}.


Either $\G$ has boundary homeomorphic to a circle, or it is a Kleinian group. In the first case we let
\[
\rho,\ovrho:\G\to\Diff^{1+\nu}(\circle)
\]
be H\"older conjugated to action of $\G$ on its boundary; if instead $\G<\PSL(2,\C)$ is a~Kleinian group we let $\rho,\ovrho:\G\to\PSL(2,\C)$ be two convex co-compact representations that lie in the same connected component of the subset of the character variety $\fkX(\G,\PSL(2,\C))$ consisting of convex co-compact representations.

We let $X$ be either the circle or $\bord\H^3$. To simplify notation we will denote the action of $\g\in\G$ on $X$ via $\rho$ by $\g$, the action via $\ovrho$ by $\ov{\!\g}$, and the limit sets of $\rho$ and $\ovrho$ by $\bord\G,\bar{\bord\G}\subset X$ respectively.

In both situations there exists a H\"older-continuous map
\[
\Xi:X\to X
\]
conjugating $\rho$ and $\ovrho$. Indeed while in the surface case this holds by definition, in the Kleinian case this is a theorem by Marden~\cite{Marden}, see also Anderson's survey~\cite[p.~32]{SurveyAnderson}: the equivariant limit map $\Xi:\bord\G\to\bord\G$ conjugating the actions $\rho$ and $\ovrho$ on their respective limit sets extends to a $\G$-equivariant, H\"older continuous homeomorphism of the whole Riemann sphere $\bord\H^3$. We study differentiability points of $\Xi$ with oblique derivative.

We let $d$ be either a visual distance on $X$ (in the complex case) or a distance inducing the chosen $\class^1$ structure on the circle $\circle$.

\begin{defi}\label{assuA}
An action $\rho$ admits a \emph{Lipschitz-compatible cover} if there exists a finite open cover $\cal B$ of $X$ and a map $\G\to\cal B$, $\gamma\mapsto\cal B_\infty(\g)$ such that:
\begin{enumerate}\romanenumi
\item\label{defi7.1.1} for any $a,b\in \G$ so that $|ab|=|a|+|b|$ one has
\begin{enumerate}
\item\label{defi7.1.1.a} $b\cal B_{\infty}(ab)\subset \cal B_\infty(a)$,
\item\label{defi7.1.1.b} $\cal B_{\infty}(ab)\subset\cal B_{\infty}(b)$;
\end{enumerate}
\item\label{defi7.1.2} there exist $\lambda>0$, $C$ and $L\in\N$ such that if
$|\g|\geq L$ and $x,y\in\cal B_\infty(\g)$ then
\[
d(\g x,\g y)\leq Ce^{-|\g|\lambda}d(x,y);
\]
\item\label{defi7.1.3} there exist constants $r_1,r_2$ and a function $\slroot:\G\to\RR$ with $\slroot(\g)\geq\lambda|\g|$ such that for every $\g\in\G$ and every $x\in\cone_\infty(\g)$,
\[
B\left(x,r_1e^{-\slroot(\g)}\right)\subset\g\cal B_\infty(\g)\subset B\left(x, r_2e^{-\slroot(\g)}\right).
\]
\end{enumerate}
\end{defi}

The goal of the subsection is to prove the following result, similar arguments can be found in Guizhen~\cite{cui} in the context of conjugacies of expanding circle maps.


\begin{prop}\label{c.indepPer}
Let $\rho,\ovrho$ be as above and assume both admit a Lipschitz compatible cover. If there exists $p\in\bord\G$ such that $\Xi$ has a finite non-vanishing derivative (complex derivative in the Kleinian case) at $p$ then $\Xi|\bord\G$ is bi-Lipschitz.
\end{prop}



We work under the assumptions of Proposition~\ref{c.indepPer} and begin its proof with the following lemma. For $\g\in\G$ we denote its derivative at $x\in X$ by $\g'(x)\in\K$ defined, according our two situations, by:
\begin{itemize}
\item $X=\bbS^1$: the derivative $\tilde\g'(\tilde x)$ of a lift of $\g$ to the universal cover $\R$ of $\bbS^1$, and a lift $\tilde x\in\R$ of $x$, the number $\tilde \g'(\tilde x)$ is independent of these choices;
\item $X=\bord\H^3$: we fix an arbitrary point $\infty\notin\bord\G$, identify $X-\{\infty\}$ with $\C$ via the stereographic projection and let $\g'(x)$ be the standard complex derivative.
\end{itemize}


\begin{lemm}\label{controlC}
Let $\rho:\G\to\Diff^{1+\nu}(X)$ admit a Lipschitz compatible cover. There exists a constant $\kappa>0$ and $N\in\N$ such that for all $\g\in\G$ with $|\g|\geq N$ and $x,y\in\cal B_\infty(\g)$ one has
\[
\bigl\lvert\log\left|\g'(x)\right|-\log\left|\g'(y)\right|\bigr|\leq \kappa d(x,y)^\nu.
\]
\end{lemm}



\begin{proof}
We consider $L$ from Definition~\ref{assuA}, so that for every $\eta\in\G$ with $|\eta|\geq L$ and $x,y\in\cal B_\infty(\eta)$ one has
\begin{equation}\label{lip}
d(\eta x,\eta y)\leq Ce^{-|\eta|\lambda}d(x,y).
\end{equation}

Since the action is $\class^{1+\nu}$ we can find a positive $K$ such that for every $\beta$ with $|\beta|\leq L$ and $u,w\in X$ one has
\begin{equation}\label{derivativeGen}
\bigl\lvert\log\left|\beta'(u)\right|-\log\left|\beta'(w)\right|\bigr|\leq Kd(u,w)^\nu.
\end{equation}

We let then $K'=\max\{K,KC^\nu\}$. We begin by showing, by induction on $k$, that if $|\g|=kL$ then for all $x,y\in\cal B_\infty(\g)$, one has
\begin{equation}\label{induc}
\bigl\lvert\log\left|\g'(x)\right|-\log\left|\g'(y)\right|\bigr|\leq K'\left(\sum_{i=0}^{k-1}e^{-\nu\lambda L i}\right)d(x,y)^\nu.
\end{equation}

Eq.~\eqref{derivativeGen} gives the base case, so assume that the result holds up to $k-1$. We write $\g=\beta\eta$ with $|\beta|=L$, $|\eta|=(k-1)L$. By Definition~\ref{assuA} (ib) we have
\begin{equation}\label{conocont}
\cal B_\infty(\g)\subseteq\cal B_\infty(\eta).
\end{equation}
Applying the chain rule gives that for every $u\in X$ one has
\[
\log\left|\g'(u)\right|=\log\left|(\beta)'(\eta u)\right|+\log\left|(\eta)'(u)\right|
\]
and thus, when $x,y\in\cal B_\infty(\g)$,
\begin{multline*}
\bigl\lvert\log\left|\g'(x)\right|-\log\left|\g'(y)\right|\bigr\rvert\\
\begin{aligned}
& \leq \bigl\lvert\log\left|\beta'(\eta x)\right|-\log\left|\beta'(\eta y)\right|\bigr\rvert+\bigl|\log\left|\eta'(x)\right| -\log\left|\eta'(y)\right|\bigr|\\
& \leq K d(\eta x,\eta y)^\nu + K'\left(\sum_{i=0}^{k-2}e^{-\nu\lambda L i}\right)d(x,y)^\nu\quad (\text{by~\eqref{derivativeGen} and induction})\\
& \leq KC^\nu e^{-|\eta|\nu\lambda}d(x,y)^\nu+ K'\left(\sum_{i=0}^{k-2}e^{-\nu\lambda L i}\right)d(x,y)^\nu\quad (\text{by~\eqref{conocont} and~\eqref{lip}}).
\end{aligned}
\end{multline*}

This shows Eq.~\eqref{induc} which implies that for $\kappa_0=K'/(1-e^{-\nu\lambda L})$, every $\g\in\G$ whose word-length is an integer multiple of $L$, and $x,y\in\cal B_\infty(\g)$ one has
\[
\bigl\lvert\log\left|\g'(x)\right|-\log\left|\g'(y)\right|\bigr|\leq \kappa_0 d(x,y)^\nu.
\]


To conclude the lemma we consider an arbitrary $\g$ with $|\g|=mL+t$ and $t< L$. We write $\g=\beta \eta$ with $|\beta|=mL$. By Definition~\ref{assuA}$\MK$\eqref{defi7.1.1.a} it holds
\begin{equation}\label{e.6.6}
\eta \cal B_\infty(\g)\subset\cal B_\infty(\beta).
\end{equation}
Applying the chain rule gives then
\begin{align*}
\bigl\lvert\log\left|\g'(x)\right|-\log\left|\g'(y)\right|\bigr|& \leq \bigl\lvert\log\left|\beta'(\eta x)\right|-\log\left|\beta'(\eta y)\right|\bigr| +\bigl\lvert\log\left|\eta'(x)\right| -\log\left|\eta'(y)\right|\bigr|\\
& \leq \kappa_0d(\eta x,\eta y)^\nu + Kd(x,y)^\nu \quad (\text{by~\eqref{derivativeGen} and~\eqref{e.6.6}})\\
& \leq \left(\kappa_0C^\nu e^{-mL\lambda}+K\right) d(x,y)^\nu \;\:(\text{by~\eqref{lip}})
\end{align*}
so taking $\kappa=K+\kappa_0C^\nu e^{-L\lambda}$ we conclude the proof of Lemma~\ref{controlC}.
\end{proof}



\begin{proof}[Proof of Proposition~\ref{c.indepPer}]
Let $p\in\bord\G$ be such that $\Xi$ has a finite non-vanishing derivative. Fix a geodesic ray $(\alpha_n)_{0}^\infty$ through the identity with $\alpha_{n}\to p$. By definition for all $n$ one has $p\in\alpha_n\cone_\infty(\alpha_n)$. Without loss of generality we may also assume that
\[
p=0=\Xi(0)
\]
and we may write the derivative as the incremental limit
\[
\Xi'(0)=\lim_{y\to 0}\frac{\Xi(y)}{y}\in\K-\{0\}.
\]
For each $n$ we let $s_n=r_1e^{-\slroot(\alpha_n)}$, so that by Definition~\ref{assuA}$\MK$\eqref{defi7.1.3},
\[
B(0,s_n)\subset\alpha_n\cal B_\infty (\alpha_n).
\]

\pagebreak{}
We consider the scaling map
\[
g_n:B(0,1)\to \alpha_n\cal B_\infty(\alpha_n)
\]
defined by $g_n(z)=s_n z$.

Let $a_n$ be an arbitrary point at distance $s_n$ from $0$ and let $\tilde s_n=\Xi(a_n)$. Observe that since $\Xi$ is differentiable at zero, for $n$ big enough the image $\Xi(B(0,s_n))$ is coarsely a ball around zero of size comparable to that of $\ov\alpha_n\cone_\infty(\ov\alpha_n)$, and in particular we can assume, since the cover $\{\cal B_\infty(\ov{\!\g})\}$ is Lipschitz compatible (Definition~\ref{assuA}$\MK$\eqref{defi7.1.3}, that $\Xi(B(0,s_n))$ is contained in $\ov\alpha_n\cal B_\infty(\ov\alpha_n)$. Furthermore we deduce that there exist positive constants $d,D$ such that for every $n$
\[
d<\frac{\ov r_2e^{-\ovsloot(\ov\alpha_n)}}{|\tilde s_n|}<D.
\]

Here we denote by $\ov r_i, \ov \lambda, \bar{C}, \ovsloot$ the constants and function associated to the Lipschitz compatible cover $\{\cal B_\infty(\ov{\!\g})\}$ for the action $\ov \rho$. We consider the scaling map
\[
\tilde g_n:B(0,D)\to B(0,|\tilde s_n|D)
\]
by $z\mapsto z\tilde s_n$.

Since $s_n\to0$ and $\Xi'(0)\notin\{0,\infty\}$ exists, the composition
\[
\tilde g_n^{-1}\Xi g_n(z)=\frac{\Xi(zs_n)}{\tilde s_n}\cdot\frac{s_nz}{s_nz}=\frac{\Xi(zs_n)}{s_nz}\cdot\frac{s_n}{\tilde s_n}\cdot z=\frac{\Xi(zs_n)}{s_nz}\cdot\frac{s_n}{\Xi(s_n)}\cdot z
\]
converges uniformly on compact subsets to the identity map.



On the other hand, one has
\[
\tilde g_n^{-1}\Xi g_n=\tilde g_n^{-1}\overline\alpha_{n}\Xi\alpha_{n}^{-1} g_n.
\]
We now study the maps $f_n:=\alpha_{n}^{-1}\circ g_n$ and $\tilde f_n:=\tilde g_n^{-1}\circ\overline\alpha_{n}$. Since the coverings $\cal B$ and $\ov {\cal B}$ are finite, we can assume, up to extracting a subsequence that there exists sets $\cal B_\infty\in \cal B$, $\ov{\cal B}_\infty\in \ov {\cal B}$ so that, for every $n$, $\cal B_\infty(\alpha_n)=\cal B_\infty$ (resp. $\cal B_\infty(\ov\alpha_n)=\ov{\cal B}_\infty$).


Observe that for every $x\in B(0,1)$ one has
\[
\log\left|f_n'(x)\right|=\log\left|\left(\alpha_n^{-1}\right)'(g_nx)\right|+\log|s_n|= -\log\left|\alpha_n'\left(\alpha_n^{-1}g_nx\right)\right|+\log|s_n|.
\]
Now by definition of $g_n$, we have that $g_nx \in\alpha_n\cal B_\infty(\alpha_n)$ and thus $\alpha_n^{-1}(g_nx)\in\cal B_\infty(\alpha_n)$. For $n$ large enough we can apply Lemma~\ref{controlC} to $\alpha_n$ to obtain $\kappa$ so that for every pair $x,y\in B(0,1)$ it holds
\[
\bigl\lvert\log\left|f_n'(x)\right|-\log\left|f_n'(y)\right|\bigr|\leq \kappa d(x,y)^\nu.
\]
We conclude that the family of maps $(f_n)$ is uniformly bi-Lipschitz on $B(0,1)$ and thus, since $(f_n0)$ is bounded, Arzela--Ascoli's theorem applies to give a subsequence (still denoted by $f_n$) that converges to a bi-Lipschitz map $f$ defined on $B(0,1)$.


A similar reasoning applies to the maps $\tilde f$ defined on $\ov{\cal B}_\infty$, and we obtain that, about~$0$, $\Xi$ can be written as a composition of bi-Lipschitz maps and is thus bi-Lipschitz. Let $U$ be the open neighborhood where $\Xi$ is bi-Lipschitz, then for every $\g\in\G$, $\Xi$ is bi-Lipschitz on $\g U$ with possibly different constants. However, by compactness we can cover $\bord\G$ by finitely many open sets such that $\Xi$ is bi-Lipschitz restricted to each one of them. A priori the bi-Lipschitz constants of $\Xi$ depend on each open set but we may choose them to work simultaneously for these finitely many translates of~$U$. It remains then to check the bi-Lipschitz property on points which are uniformly far apart, i.e. whose distance is $\geq\eps$ for some $\eps>0$. However, the function
\[
(x,y)\mapsto \frac{d(f(x),f(y))}{d(x,y)}
\]
is continuous and non-vanishing on $\{(x,y):d(x,y)\geq\eps\}$. Since the latter space is compact we obtain the bi-Lipschitz property for points which are far apart, as desired.
\end{proof}


The following Lemma guarantees we can later apply the results of this section to the situation of our interest.


\begin{lemm}\label{usar}
\leavevmode
\begin{itemize}
\item Assume $\bord\G$ is homeomorphic to a circle and let $\rho:\G\to\SL(d,\R)$ be $(1,1,2)$-hyperconvex. Then the induced action of $\rho(\G)$ on the $\class^{1+\nu}$ circle $\xi(\bord\G)$ admits a Lipschitz compatible covering.
\item If $\G$ is a convex co-compact Kleinian group then the action of $\G$ on $\bord_\infty\H^3$ admits a Lipschitz compatible cover.
\end{itemize}
\end{lemm}

\begin{proof}
Recall from Section~\ref{cont} that we have fixed a word metric on $\G$ and we denote by $\cone_\infty(\g)\subset\partial\G\subset X$ the set of endpoints of geodesic rays contained in the cone type $\cone(\g)$.

Let $\delta_\rho$ be the fundamental constant of $\rho$ from Definition~\ref{FConstant}, and let $\cal B_{\infty}(\g)=X_\infty(\gamma)$ be the $\delta_\rho/2$-neighbourhood of $\cone_\infty(\g)$ inside $X$. This is the thickened cone type at infinity considered in~\cite[Section~5]{PSW1} (see also the proof of Proposition~\ref{conetypesBalls}). It is a proper subset of $X$ by Corollary~\ref{c.fundamentalConstant}. The cover $\cal B$ is finite since there are only finitely many cone types~\cite[p.~455]{BH}.

Definition~\ref{assuA}$\MK$\eqref{defi7.1.1} holds since the same property holds for $\cone_\infty(\g)$, property~\eqref{defi7.1.2} is a~consequence of Proposition~\ref{Lipschitz-compatibleB}. Finally, Property~\eqref{defi7.1.3} was proven in~\cite[Corollary~5.10]{PSW1} choosing $\slroot(\g):=\slroot_1(\cartan\rho(\g))$ (see also the proof of Proposition~\ref{conetypesBalls}). Observe that in the real case by considering $X=\circle$ we are implicitly considering only the intersection with the limit set, while in the Kleinian group case it is not necessary to intersect with the limit set since the $\G$-action on the whole $X$ is conformal.
\end{proof}


We now establish the following corollary that will be used in the sequel.

\begin{coro}\label{diffpointreal}
Assume $\bord\G$ is homeomorphic to a circle. Let $\rho:\piS\to\PGL(d,\R)$ and $\ovrho:\piS\to\PGL(\ovd,\R)$ be $(1,1,2)$-hyperconvex, consider the map between $\class^{1+\nu}$ circles
\[
\Xi=\ovxi\circ\xi^{-1}:\xi(\bord\piS)\to\ovxi(\bord\piS).
\]
If $\Xi$ has a differentiability point with finite non-vanishing derivative then $\rho$ and $\ovrho$ are gap-isospectral.
\end{coro}

\begin{proof}
By Lemma~\ref{usar} we can apply Proposition~\ref{c.indepPer} to obtain that $\Xi$ is bi-Lipschitz. The following standard lemma from linear algebra (see for example~\cite{convexes1} and~\cite[Lemma~3.4]{exponential}) gives the period computation completing the proof.
\end{proof}



\begin{lemm}\label{lambda2}
Let $g\in\PGL(d,\R)$ be proximal with attracting point $g_+\in\P(\R^d)$ and repelling hyperplane $g_-\in\P((\R^d)^*)$. Let $V_{\lambda_2(g)}$ be the sum of the characteristic spaces of $g$ whose associated eigenvalue is of modulus $\exp\lambda_2(g)$, Then for every $v\notin\P(g_-)$, with non-zero component in $V_{\lambda_2(g)}$, one has
\[
\lim_{n\to\infty}\frac{\log d_\P(g^n(v),g_+)}n= -\slroot_1(\lambda(g)).
\]
\end{lemm}



\subsection{Limit curves in non-maximal flags}

We proceed with another intermediate step for the proof of Theorem~\ref{t.Zcl} describing differentiability points of boundary maps in partial flag manifolds $\calF_{\{\sroot,\bb\}}$ for $\{\sroot,\bb\}$-Anosov representations. This step follows from the combination of Theorem~\ref{t.C1} and Corollary~\ref{diffpointreal}.



Let $\sfG$ be real-algebraic and semi-simple. Let $\{\sroot,\bb\}\subset\simple$ be two distinct simple roots. The partial flag space $\calF_{\{\sroot,\bb\}}$ carries two transverse foliations that are the level sets of the natural projections $\calF_{\{\sroot,\bb\}}\to\calF_{\{\sroot\}}$ and $\calF_{\{\sroot,\bb\}}\to\calF_{\{\bb\}}$. We will refer to these as the \emph{canonical foliations} of $\calF_{\{\sroot,\bb\}}$.


\begin{coro}\label{nonC1}
Let $\sfG$ be real-algebraic and semi-simple and let $\{\sroot,\bb\}\subset\simple$ distinct. Let $\rho:\piS\to\sfG$ be Zariski-dense and $\{\sroot,\bb\}$-Anosov. If both curves $\xi^\sroot(\bord\piS)$ and $\xi^{\bb}(\bord\piS)$ are $\class^{1+\alpha}$ then every differentiability point of $\xi^{\{\sroot,\bb\}}(\bord\piS)$ is tangent to one of the canonical foliations of $\calF_{\{\sroot,\bb\}}$.
\end{coro}

\begin{proof}
By Benoist's Theorem~\ref{densidad} the limit cone of $\rho$ has non-empty interior, in particular there exists $\g\in\piS$ such that
\begin{equation}\label{isospe}
\sroot(\lambda(\g))\neq\bb(\lambda(\g)).
\end{equation}

Consider the Tits representations $\Fund_\sroot$ and $\Fund_\bb$ associated to $\sroot$ and $\bb$. Since $\rho(\piS)$ is Zariski-dense, both representation $\Fund_\sroot\rho$ and $\Fund_\bb\rho$ are irreducible and since $\rho$ is $\{\sroot, \bb\}$-Anosov both representation $\Fund_\sroot\rho$ and $\Fund_\bb\rho$ are projective Anosov. Recall that by definition of $\Fund_\sroot$, for every $g\in\sfG$ one has
\[
\slroot_1(\lambda(\Fund_\sroot(g)))=\sroot(\lambda(g)),
\]
so by Eq.~\eqref{isospe} the representations $\Fund_\sroot\rho$ and $\Fund_\bb\rho$ are not gap-isospectral.

Since the maps $\zeta_\sroot$ and $\zeta_\bb$ are analytic, both projective curves $\zeta_\sroot\xi^{\sroot}(\bord\piS)$ and $\zeta_\bb\xi^{\bb}(\bord\piS)$ are $\class^{1+\alpha}$ and thus by Zhang--Zimmer's Theorem~\ref{t.C1} the representations $\Fund_\sroot\rho$ and $\Fund_\bb\rho$ are $(1,1,2)$-hyperconvex.

The natural embedding $\calF_{\{\sroot,\bb\}}\to\P(V_\sroot)\times\P(V_\bb)$ sends $\xi^{\{\sroot,\bb\}}$ to the graph of the map $\Xi$ from Corollary~\ref{diffpointreal} and thus the corollary implies the result.
\end{proof}



\subsection{Proof of Theorem~\ref{t.Zcl}}\label{s.pB}
\label{s.Zd}
The goal of the section is to prove Theorem~\ref{t.Zcl}, stating that the Zariski closure $\sfG$ of the image of an irreducible $(1,1,2)$-hyperconvex representation $\rho:\piS\to\PGL(d,\R)$ is simple and the highest weight of the induced representation $\Fund:\sfG\to\PGL(d,\R)$ is a multiple of a fundamental weight associated to a root whose root-space is one-dimensional.


It is known that an irreducible subgroup $\sfG<\PGL(d,\R)$ containing a proximal element is semi-simple without compact factors (see~\cite[Lemma~8.6]{entropia} for an explicit argument following a suggestion by Quint).

We consider the induced representation $\rho_0:\G\to\sfG$ and denote by $\Fund:\sfG\to\PGL(d,\R)$ the linear representation so that $\rho=\Fund\rho_0$. Let $\chi=\chi_\Fund\in\a^*$ be the highest weight of $\Fund$. As in Definition~\ref{trep} we consider
\[
\t=\t_\Fund=\left\{\sroot\in\Delta\,:\, \chi-\sroot \text{ is a weight of }\Fund \right\}
=\left\{\sroot\in\simple\,:\,\<\chi,\sroot\>\neq0\right\}.
\]
It is enough to show that $\t$ is reduced to a single root $\{\sroot_0\};$ indeed, if this is the case, upon writing $\chi$ in the basis of fundamental weights $\{\peso_\sroot:\sroot\in\simple\}$ (recall their defining Eq.~\eqref{pesoFund}) one has
\[
\chi=\sum_{\sroot\in\simple}\<\chi,\sroot\>\peso_\sroot=\<\chi,\sroot_0\>\peso_{\sroot_0},
\]
Moreover this gives:
\begin{itemize}
\item $\sfG$ is simple by Lemma~\ref{simple};
\item the weights on the first level consist solely on $\chi-\sroot_0$ and its associated weight space is $\phi(\ge_{-\sroot_0})V_{\chi_\Fund}$. By definition, if $a\in\inte\a^+$ then $\Fund(\exp(a))$ is proximal with associated eigenline $V^\chi$ and the generalized eigenspace associated with the second eigenvalue, in modulus, is $\phi(\ge_{-\sroot_0})V_{\chi_\Fund}$. Since $\rho(\G)$ is $\{\slroot_1,\slroot_2\}$-Anosov, we know that for all $\g\in\G$, $\rho(\g)$ and $\wedge^2\rho(\g)$ are proximal, thus the generalized eigenspace associated to the second eigenvalue of $\rho(\g)$ is $1$-dimensional. Combining both facts we conclude that $\phi(\ge_{-\sroot_0})V_{\chi_\Fund}$ is one-dimensional, but by Lemma~\ref{not} no element of $\ge_{-\sroot_0}$ acts trivially on $V_{\chi_\Fund}$ so $\ge_{-\sroot_0}$ is $1$-dimensional, as desired.
\end{itemize}



We proceed now to show that in the present situation $\t$ consists of only one element.

By definition of $\t$ one has, for every $g\in\sfG$, that
\[
\slroot_1\bigl(\lambda(\Fund(g))\bigr)=\min_{\sroot\in\theta}\bigl\{\sroot(\lambda_\sfG(g))\bigr\}.
\]
Consequently, the limit cone $\calL_{\rho_0}\subset\a^+_\sfG$ does not intersect the walls of elements in $\t$ and, since $\rho_0:\G\to\sfG$ is a quasi-isometry, Remark~\ref{wallsAnosov} implies that the representation $\rho_0$ is $\t$-Anosov.

Recall from Eq.~\eqref{maps} that we have a $\Fund$-equivariant analytic embedding $\zeta_\t:\sfG/\sf P_{\theta}\to \P(\R^{d})$. One has moreover that $\xi^1_\rho=\zeta_\t\circ\xi^\t_{\rho_0}$. Since $\xi^2_\rho$ is H\"older-continuous we conclude that the boundary map $\xi^\theta$ has $\class^{1+\alpha}$-image. Composing with the projections $\calF_\t\to\calF_{\t'}$ one sees that, for any $\t'\subset\t$ the curve $\xi_{\rho_0}^{\t'}(\bord\G)$ is a also a $\class^{1+\alpha}$ circle.


Assume now there exists two distinct roots $\sroot,\bb$ in $\t$ and consider the subset $\t'=\{\sroot,\bb\}$. By the previous paragraph the curve $\xi^{\t'}(\bord\G)$ is $\class^{1+\alpha}$. Corollary~\ref{nonC1} gives then that $\xi^{\t'}(\bord\G)$ is necessarily contained in one of the leaves of the canonical foliations of $\calF_{\t'}$, thus giving that one of the maps $\xi^{\sroot}$ or $\xi^\bb$ is constant, achieving a~contradiction. \qed



\section{Non-differentiability and \texorpdfstring{$1$}{1}-conicality: the proof of Theorem~\ref{tutti}}\label{s.Hffproof}


\subsection{Non-differentiability and \texorpdfstring{$1$}{1}-conicality}



By means of Section~\ref{s.6.1} we can improve Lemma~\ref{generalcase} when we deal with a pair of real hyperconvex representations of surface groups, this is the missing ingredient for Theorem~\ref{tutti}:



\begin{coro}\label{nondiffbconicalR}
Assume $\bord\G$ is homeomorphic to a circle. Let $\rho,\ovrho$ two $(1,1,2)$-hyperconvex representations over $\R$ of $\G$ that are not gap-isospectral. Then, the set of non-differentiability points of $\Xi$ coincides with the set of $1$-conical points.
\end{coro}



\begin{proof}
We choose a $\class^1$ identification of the $\class^1$ torus $\xi(\bord\G)\times\ovxi(\bord\G)\subset\P(\R^d)\times \P(\R^{\ovd})$ with the quotient of the square $[-1,1]\times[-1,1]$ preserving the product structure, and such that the point $(x,\Xi(x))$ corresponds to $(0,0)$. In these coordinates the graph of $\Xi$ is a monotone curve $[-1,1]\to[-1,1]$ passing through the origin. Since the chosen identification is $\class^1$, it is in particular $K$-bi-Lipschitz for some $K$, so we can write (coarsely in a small neighbourhood of $x$) $d(\xi(y),\xi(x))=|y|$ and $d(\ovxi(y),\ovxi(x))=|\Xi(y)|$.

From Lemma~\ref{generalcase} we know that $x$ is $1$-conical if and only if either $\lim_{y\to x}\frac{|\Xi(y)|}{|y|}$ exists and is far from $0$ and $\infty$, either it does not exist. The proposition is settled if we show that the first situation cannot happen, so let's assume it does. However, since $\Xi$ is monotone we can remove the $|\,|$ and we get that $x$ is a differentiability point of $\Xi$ with oblique derivative. Corollary~\ref{diffpointreal} implies then that for all $\g\in\G$ one has $\slroot_1(\lambda(\g))=\ovsloot_1(\lambda(\ov\g))$, contradicting our assumption.
\end{proof}



\subsection{Proof of Theorem~\ref{tutti} and an analogous for Kleinian groups}

We begin with the proof of Theorem~\ref{tutti} by recalling the following result from \cite{BeyP2}.

\begin{coro}[{\cite{BeyP2}}]\label{l.weaklyirred}
Assume $\bord\G$ is homeomorphic to a circle and let $\rho:\G\to \PGL(d,\R)$ be $(1,1,2)$-hyperconvex. Then there exists an irreducible $(1,1,2)$-hyperconvex representation $\rho_0:\G\to \PGL(m,\R)$ such that, for every $\g\in\G$ one~has
\[
\slroot_1(\lambda(\g))=\slroot_1(\lambda(\rho_0\g)).
\]
\end{coro}


We now prove Theorem~\ref{tutti}. Since there exists $\g\in\G$ with $\slroot_1(\lambda(\g))\neq\ov{\slroot}_1(\lambda(\ov\g))$, Corollary~\ref{l.weaklyirred} allows us to apply Proposition~\ref{nonA} to obtain the density assumption in Theorem~\ref{Hffconical}, so one has
\[
\Hff\Xi(\{1\text{-conical points}\})=\hC{\max\{\slroot,\ov{\slroot}\}}.
\]
Corollary~\ref{nondiffbconicalR} states that the set of $1$-conical points coincides with the set of non-differentiability points of $\Xi$. The inequality $\hC{\infty,1} <1$ then follows from Theorem~\ref{Hffconical}. This completes the proof of Theorem~\ref{tutti}.



\subsection{Proof of Corollary~\ref{cor.B}} We conclude the paper proving Corollary~\ref{cor.B}. Recall from Section~\ref{representaciones} that for every simple root $\sroot$ of $\sfG$ we chose a Tits representation $\Phi_\sroot:\sfG\to\PSL(V_\sroot)$.

\begin{coro}\label{c.Hk}
Assume $\bord\G$ is homeomorphic to a circle and let $\sfG$ be a simple Lie group. Let $\rho:\G\to\sfG$ have Zariski-dense image. If for $\sroot,\bb\in\simple$ the representations $\Phi_\sroot\circ\rho$ and $\Phi_\bb\circ \rho$ are $(1,1,2)$-hyperconvex,
then
\begin{enumerate}\romanenumi
\item\label{coro8.3.1} the image of the limit curve $\xi^{\{\sroot,\bb\}}:\bord\G\to\calF_{\{\sroot,\bb\}}$ is Lipschitz and the Hausdorff dimension of the points where it is non-differentiable is $\hC{\max\{\sroot,\bb\}}$.
\item\label{coro8.3.2} If the opposition involution $\ii$ on $\ge$ is non-trivial and $\bb=\ii\sroot$ then
\[
\hC{\max\{\sroot,\bb\}}=\hC{(\sroot+\bb)/2}.
\]
\end{enumerate}
\end{coro}

\begin{proofc}
\begin{proof}[{\eqref{coro8.3.1}}]
Since the map $\Phi_\sroot:\calF_{\sroot}\to\P(V_{\sroot})$ is analytic, and $\Phi_\sroot\circ\xi^\sroot(\bord\G)$ is a $\class^1$-submanifold (Theorem~\ref{t.C1}), $\xi^\sroot(\bord\G)$ is a $\class^1$ submanifold as well. The curve $\xi^{\{\sroot,\bb\}}:=\calF_{\{\sroot,\bb\}}\cap (\xi^{\sroot}(\bord\G)\times\xi^{\bb}(\bord\G))$ is the graph of the homeomorphism $\Xi$ and is thus a Lipschitz curve. The second claim is then a direct consequence of Theorem~\ref{tutti}.\let\qed\relax
\end{proof}

\begin{proof}[{\eqref{coro8.3.2}}]
Assume the opposition involution $\ii$ of $\ge$ is non-trivial and that $\bb=\ii\sroot$. Using notation from Section~\ref{s.lower} with $\sroot=\slroot$ and $\bb=\ii\sroot=\ov{\slroot}$ we let $V^*=\spa\{\sroot,\bb\}$, $V=\a_\t/\Ann(V^*)$, $\Pi:\a_\t\to V$ the quotient projection, $\|\,\|_\infty=\max\{|\sroot|,|\bb|\}$, $\|\,\|^1$ its dual norm on $V^*$ and ${\varphi^\infty_1}\in\calQ_\cc$ the only form minimizing $\|\,\|^1$.

Since $\ii\sroot=\bb$, the space $V^*$ is preserved by $\ii$ and the fact that $\lambda(g^{-1})=\ii\lambda(g)$ (for all $g\in \sfG$) implies that $\calQ_\cc$ is $\ii$-invariant. Moreover, the norm $\|\,\|^1$ is $\ii$-invariant and by definition of ${\varphi^\infty_1}$ one has $\ii{\varphi^\infty_1}={\varphi^\infty_1}$. However, $(\sroot+\bb)/2$ is also $\ii$-invariant and $\hC{(\sroot+\bb)/2}(\sroot+\bb)/2\in\calQ_\cc$ whence
\[
{\varphi^\infty_1}=\hC{(\sroot+\bb)/2}(\sroot+\bb)/2.
\]
%\end{proof}


We now consider a refinement of Corollary~\ref{quintgrowththeta}. Indeed, examining its proof one readily sees that the statement holds replacing $\a_\t$ with $V$ and $\cartan_\t$ with $\Pi\cartan_\t$, thus one has $\hC{\max\{\sroot,\bb\}}=\|{\varphi^\infty_1}\|^1$.
\end{proof}
\let\qed\relax
\end{proofc}

\subsection{The \texorpdfstring{$\PSL(2,\C)$}{PSL(2, C)}-case}\label{klein}

If $\rho,\ovrho:\G\to\PSL(2,\C)$ are convex co-compact representations that are connected by convex co-compact representations, it was proven by Marden~\cite{Marden} that the natural map $\Xi:\Lambda_\rho\to\Lambda_{\ovrho}$ conjugating the respective actions extends to a H\"older homeomorphism $\Xi:\C\P^1\to\C\P^1$ that is $(\rho,\ovrho)$-equivariant. We consider in this case the complex derivative of such an extension $\Xi$ and say that $\Xi$ is $\C$-\emph{differentiable} at a given $x\in\Lambda_\rho$ if, conformally identifying $\bord\H^3-\{\text{point}\}$ to $\C$, the limit
\[
\Xi'(x):=\lim_{y\to x} \frac{\Xi(x)-\Xi(y)}{x-y}
\]
exists or is infinite. We let now $\Ndiff_{\rho,\ovrho}$ be the set of points $x\in\Lambda_\rho$ where the extended conjugating map $\Xi$ is not $\C$-differentiable.


The proof of the following works verbatim as in Corollary~\ref{nondiffbconicalR}.

\begin{prop}\label{nondiffbconicalC}
Let $\rho,\ovrho:\G\to\PSL(2,\C)$ be non-gap-isospectral and in the same connected component of
\[
\bigl\{\varrho:\G\to\PSL(2,\C):\,\varrho\text{\ is\ convex\ co-compact}\bigr\}.
\]
Then, the set of non-$\C$-differentiability points of $\Xi$ coincides with the set of $1$-conical points.
\end{prop}



Density of the group generated by the pairs $\{(\lambda(\g),\lambda(\bar\g)):\g\in\G\}$ follows readily from Benoist~\cite{benoist2} (see Theorem~\ref{densidad}), from this point on the exact same proof of Theorem~\ref{tutti} gives the following.

\begin{theo}\label{thm.tuttiC}
Let $\rho,\ovrho:\G\to\PSL(2,\C)$ be non-gap-isospectral convex co-compact representations that are connected by convex co-compact representations. Assume without loss of generality that $\hJ{\ovsloot}\geq \hJ{\slroot}$. If $\II_{\slroot}(\ovsloot)>1$, then
\[
\Hff\left(\Ndiff_{\rho,\ovrho}\right)=\hC{\infty,1}.
\]
\end{theo}



\bibliography{20251124d_pozzetti}
\end{document}
